\chapter{模空间}\label{chap:moduli}
\SourcePage{136}
本章的目标是描述复结构模空间与 Kähler 模空间，并解释如何在 Calabi--Yau 流形 $V$ 的 Kähler 模空间与其镜像 $V^\circ$ 的复结构模空间之间构造镜像映射。

本章的讨论将揭示，复结构模与 Kähler 模具有出人意料的丰富内容。我们将特别关注环面簇中的 Calabi--Yau 超曲面，因为在这一情形，可以显式描述这两类模空间的形态，尽管其中有些结论仍属猜想。实际上，只研究这些模空间的多项式部分和环面部分，但它们已经足够有趣。

复结构模与 Kähler 模之间的密切联系，通过镜像映射建立起来；镜像映射将在本章末定义。困难在于，不了解 Kähler 模就无法充分理解复结构模，反之亦然。此外，理解两者都需要了解镜像映射，而要充分理解镜像映射，又必须等到\hyperref[chap:quantum-cohomology]{第~\ref*{chap:quantum-cohomology} 章}和\hyperref[chap:mirror-theorem]{第~\ref*{chap:mirror-theorem} 章}研究量子上同调及量子微分方程以后。因此，我们的讨论虽然表面上依次从复结构模走向 Kähler 模，再走向镜像映射，实际上却会在这些主题之间往返穿插。

\section{复结构模}\label{sec:complex-moduli}\index{复结构模空间|hyperpage}
本节讨论 Calabi--Yau 流形的复结构模。\SourceCite{Viehweg} 表明，\emph{极化} Calabi--Yau 流形 $(V,\mathcal L)$ 的模可由光滑拟射影簇参数化。我们将关注该模空间的一个连通分支 $\mathcal M$。

目标是为 $\mathcal M$ 找到一个方便使用的紧化 $\overline{\mathcal M}$。我们采取实用的立场，因此会灵活理解“方便”的含义。最低限度的要求是：在与镜像对称有关的最大幺幂边界点附近，紧化边界除子具有正规交叉。由奇点消解，具有正规交叉边界的紧化总是存在，却远不唯一。不过，为什么这样的紧化应具有最大幺幂单值性，目前并不清楚。

先讨论最大幺幂单值性，随后主要关注\hyperref[sec:batyrev-mirror]{第 4.1 节}所述 Batyrev 构造得到的 Calabi--Yau 超曲面的模空间，再详细说明一些例子。

\subsection{最大幺幂边界点}\label{subsec:mum-boundary-points}
设在去心多圆盘 $(\Delta^*)^r\subset\Delta^r$ 上有一个 Calabi--Yau 流形族，且在原点的单值性为最大幺幂。这就是所谓的\emph{最大幺幂边界点}。令 $z_1,\ldots,z_r$ 为 $\Delta^r$ 的坐标，$g_0,\ldots,g_r$ 如最大幺幂单值性的定义~\ref{def:maximal-unipotent-point} 所述。
\SourcePage{137}
再用下式定义全纯函数 $q_1,\ldots,q_r$ \SourceCite{Morrison6}：
\begin{equation}\label{eq:mum-canonical-coordinates}
\frac1{2\pi i}\log q_k=\frac1{\langle g_0,\Omega\rangle}
\sum_{j=1}^r\langle g_j,\Omega\rangle m^{jk},
\end{equation}
其中 $(m^{jk})$ 为定义~\ref{def:maximal-unipotent-point} 中矩阵 $(m_{ij})$ 的逆矩阵。

计算 $q_k$ 在原点附近的行为。由 $N_i g_j=m_{ij}g_0$，注意到沿 $z_i=0$ 有
$\langle g_j,\Omega\rangle\sim m_{ij}\langle g_0,\Omega\rangle\log(z_i)/(2\pi i)$。
把所有 $z_i$ 对应的渐近式代入式~\eqref{eq:mum-canonical-coordinates}，得到在原点附近 $q_k\sim cz_k$，其中 $c\ne0$ 为某个常数。因此，$(q_1,\ldots,q_r)$ 在最大幺幂边界点附近定义局部坐标。此外，这些坐标几乎唯一。若把 $g_0,\ldots,g_r$ 换成 $g'_j=\sum_{k=0}^r c_{jk}g_k$，就得到新的局部坐标 $(q'_1,\ldots,q'_r)$；可以证明
\begin{equation}\label{eq:mum-coordinate-basis-change}
q'_k=e^{2\pi i(c^{k0}/c_{00})}q_k,\qquad
c^{k0}=\sum_{j=1}^r c_{j0}m'^{jk},
\end{equation}
其中 $(m'^{jk})$ 是由 $N_j(g'_k)=m'_{jk}g'_0$ 确定的矩阵的逆矩阵。\hyperref[sec:mirror-map]{第~\ref*{sec:mirror-map} 节}将讨论如何唯一确定 $q_k$。

遗憾的是，上述“唯一性”假定始终使用 $(\Delta^*)^r$ 的同一个正规交叉紧化。因此，需要确定改用另一紧化时，$q_k$ 如何变化。只需局部处理此事，所以考虑全纯映射 $\phi:\Delta^r\to\Delta^r$，在源和目标的去心部分采用环面坐标；假定它在稠密环面上为同构，并把坐标超平面映到坐标超平面组成的边界内。必要时缩小圆盘后，很容易看出：可以在源上选择坐标 $w=(w_1,\ldots,w_r)$，在目标上选择坐标 $z=(z_1,\ldots,z_r)$，使映射为
\begin{equation}\label{eq:mum-monomial-coordinate-change}
z_j=\phi_j(w)=\prod_i w_i^{a_{ij}},
\end{equation}
其中 $a_{ij}$ 为非负整数，矩阵 $(a_{ij})$ 在 $\mathbb Z$ 上可逆。

我们断言，源中的 $w=0$ 是最大幺幂边界点，当且仅当目标中的 $z=0$ 是最大幺幂边界点。为此，先比较绕除子 $z_i=0$ 与绕除子 $w_i=0$ 的单值性。分别用 $T_i^z,N_i^z$ 表示绕除子 $z_i=0$ 的单值性及其对数，类似地用 $T_i^w,N_i^w$ 表示绕除子 $w_i=0$ 的单值性及其对数。沿一条绕 $w_i=0$ 一周的路径，例如
\[
\theta\longmapsto(p_1,\ldots,p_{i-1},p_i e^{2\pi i\theta},p_{i+1},\ldots,p_r),
\]
代入式~\eqref{eq:mum-monomial-coordinate-change} 后，可以看到：与 $\phi$ 复合所得路径绕 $z_j=0$ 共 $a_{ij}$ 周。因此，对每个 $i$，得到
\[
T_i^w=\prod_j(T_j^z)^{a_{ij}}.
\]
现在假定 $T_j^z$ 为幺幂算子。由于它们彼此交换，$T_i^w$ 也为幺幂；反向结论则由 $(a_{ij})$ 的可逆性推出。于是可以取对数，得到
\begin{equation}\label{eq:mum-log-coordinate-change}
N_i^w=\sum_j a_{ij}N_j^z.
\end{equation}

\SourcePage{138}
为便于后面使用，注意该式可以写成
\begin{equation}\label{eq:mum-log-matrix-change}
\begin{pmatrix}N_1^w\\\vdots\\N_r^w\end{pmatrix}
=A\begin{pmatrix}N_1^z\\\vdots\\N_r^z\end{pmatrix},\qquad A=(a_{ij}).
\end{equation}
接着比较单值性权滤过。在源上，它由 $N=\sum_i a_iN_i^w$ 确定，其中 $a_i>0$。不过，由式~\eqref{eq:mum-log-coordinate-change}，有
\[
N=\sum_i a_iN_i^w=\sum_j\left(\sum_i a_i a_{ij}\right)N_j^z.
\]
已假定对所有 $i,j$ 都有 $a_{ij}\geq0$；由于 $(a_{ij})$ 可逆，$\sum_i a_i a_{ij}>0$。这说明在目标上也可使用 $N$，所以源和目标具有相同的单值性权滤过。由此很容易看出，在 $w=0$ 处满足定义~\ref{def:maximal-unipotent-point}，当且仅当在 $z=0$ 处满足它，断言得证。

在这种情况下，用 $q_i^z,q_i^w$ 表示不同局部紧化中由式~\eqref{eq:mum-canonical-coordinates} 确定的坐标。计算得
\begin{equation}\label{eq:mum-q-monomial-change}
q_j^z=\prod_i(q_i^w)^{a_{ij}}.
\end{equation}
注意它与式~\eqref{eq:mum-monomial-coordinate-change} 的相似性。

所以，一个最大幺幂边界点通过式~\eqref{eq:mum-monomial-coordinate-change} 可以产生\emph{许多}其他最大幺幂边界点，进而产生许多组局部坐标 $(q_1,\ldots,q_r)$；它们都通过式~\eqref{eq:mum-q-monomial-change} 联系起来。在\hyperref[subsec:large-radius-limits]{第 6.2.1 小节}将看到，Kähler 一侧的大半径极限点的局部坐标有\emph{完全相同}的不确定性。

回到整体情形，设 $\overline{\mathcal M}$ 是复结构模空间 $\mathcal M$ 的一个光滑正规交叉紧化。若 $\overline{\mathcal M}$ 有一个最大幺幂边界点，上述讨论就表明：通过爆破，可以得到具有多个最大幺幂边界点的其他紧化。$\overline{\mathcal M}$ 本身也可能是一个爆破所得空间。事实上，$\overline{\mathcal M}$ 可能有多个最大幺幂边界点，在爆缩到一个奇异紧化时，这些点会被识别为同一点。

整理这些可能性的一种办法，是在最大幺幂边界点之间引入等价关系。依照 \SourceCite{Morrison6}，作如下定义。
\begin{definition}\label{def:equivalent-mum-points}
设 $\overline{\mathcal M}$ 是模空间 $\mathcal M$ 的一个光滑正规交叉紧化。若存在包含 $p,q$ 的连通集 $\Xi\subset\overline{\mathcal M}-\mathcal M$，并且在 $\Xi$ 的一个邻域上存在由 $\overline{\mathcal M}$ 的余切丛中、沿边界有对数极点的形式组成的局部系统 $\mathcal L$，使得在 $p,q$ 各自附近，该局部系统由 $1$ 形式 $d\log q_k$ 张成，其中 $q_k$ 如式~\eqref{eq:mum-canonical-coordinates} 所述，则称两个最大幺幂边界点 $p,q\in\overline{\mathcal M}$ \emph{等价}。
\end{definition}
\index{最大幺幂边界点!等价关系|hyperpage}
很容易看出，爆破产生的多个最大幺幂边界点，在这一意义下彼此等价。

\IndexLocation{indexloc:0830:138}镜像对称的一个基本思想是：Calabi--Yau 流形 $V$ 的复结构模空间中，最大幺幂边界点的每个等价类，应当对应于其镜像 $V^\circ$ 的 Kähler 模空间中一个适当定义的
\SourcePage{139}
大半径极限点等价类（将在\hyperref[subsec:large-radius-limits]{第 6.2.1 小节}讨论）。

实际处理中，从一个选定的 $\overline{\mathcal M}$ 出发，通常会尽可能进行爆缩，直到得到一个无法继续爆缩的光滑正规交叉紧化。这样可以使等价类尽可能小。\hyperref[subsec:two-parameter-moduli]{第 6.1.4 小节}将给出一个例子。

还要指出，定义~\ref{def:equivalent-mum-points} 的等价关系，也可在 $\mathcal M$ 的更一般紧化上定义。\SourceCite{Morrison3} 已经做了这项工作。

最后，注意我们还没有讨论最大幺幂边界点的\emph{存在性}。这仍然是猜想，尽管迄今计算的每个例子都有这样的点。很快还会看到，环面超曲面的模空间具有自然出现的紧化，其中有一些特殊边界点；猜想这些点是最大幺幂的（\SourceCite{HLY2} 给出了一个很好的部分结果）。一般而言，镜像对称预言 $\mathcal M$ 存在十分特殊的紧化，它们不仅具有最大幺幂边界点，而且反映镜像的 Kähler 模结构。本章后面还会进一步讨论这一点。

\subsection{环面超曲面的复结构模}\label{subsec:toric-complex-moduli}\index{多项式模空间|hyperpage}
研究 Calabi--Yau 环面超曲面的复结构模。回顾\hyperref[sec:batyrev-mirror]{第 4.1 节}，对自反多胞形 $\Delta$，反典范超曲面 $\overline V\subset\mathbb P_\Delta$ 是 Calabi--Yau，尽管可能有奇点。若 $\Sigma$ 是 $\Delta$ 的法扇的一个极大射影细分，已经证明 $\overline V$ 的严格变换仍为反典范超曲面，从而得到一族极小 Calabi--Yau 超曲面 $V\subset X=X_\Sigma$。所以 $V$ 是一个至多具有终端奇点的 Gorenstein 轨形；当 $\dim V=3$ 时，它是光滑的。

不处理 $V$ 的整个复结构模空间 $\mathcal M$，而集中研究其中的“多项式”部分 $\mathcal M_{\mathrm{poly}}\subset\mathcal M$，即那些可实现为 $X$ 中超曲面的复结构。后面将看到，若 $\Delta$ 的任何余维二的面都没有内部格点，那么 $\mathcal M_{\mathrm{poly}}$ 就是整个复结构模空间。

为构造多项式模空间，从形如 $\sum_{m\in\Delta\cap M}\lambda_m t^m$ 的 Laurent 多项式构成的向量空间 $L(\Delta\cap M)$ 出发。与\hyperref[subsec:toric-polytopes]{第 3.2.1 小节}一样，$t^m$ 是对应于 $m$ 的特征标。（\hyperref[chap:toric]{第~\ref*{chap:toric} 章}中，将 $L(\Delta\cap M)$ 记作 $L(\Delta)$；这里希望更明确地体现其对格的依赖。）由式~\eqref{eq:toric-sections}，有自然同构
\[
H^0(X,\mathcal O_X(-K_X))\simeq L(\Delta\cap M).
\]
可以把 $L(\Delta\cap M)$ 中的一个多项式看作定义 $X$ 中一个反典范超曲面的仿射部分。此外，$-K_X=\sum_\rho D_\rho$ 表明
\[
\Delta\cap M=\{m\in M:\langle m,v_\rho\rangle\geq-1\text{，对所有 }\rho\}.
\]
由于典范微分的拉回仍是典范微分，\IndexLocation{indexloc:1355:139}自同构群 $\operatorname{Aut}(X)$ 自然作用于射影空间 $\mathbb P(L(\Delta\cap M))$，所以多项式模空间是商
\[
\mathcal M_{\mathrm{poly}}=\mathbb P(L(\Delta\cap M))/\operatorname{Aut}(X).
\]
要使这一表述精确，需将 $\mathbb P(L(\Delta\cap M))$ 换成对应于拟光滑超曲面的子集，并考虑商的存在性。后一问题很重要，因为 $\operatorname{Aut}(X)$ 未必为约化群。不过，如 \SourceCite{BC} 第 13 节所指出的，可以找到一个非空不变开集
\SourcePage{140}
$U\subset\mathbb P(L(\Delta\cap M))$，使几何商 $\mathcal M_{\mathrm{poly}}$ 存在。因此，至少在一般位置上可以构造多项式模空间。

若能更清楚地了解这一模空间，当然更好。因此，接着讨论 \SourceCite{AGM1} 引入的\emph{简化多项式模空间} $\mathcal M_{\mathrm{simp}}$。令 $(\Delta\cap M)_0$ 为 $\Delta\cap M$ 中不属于 $\Delta$ 任何刻面内部的点所组成的集合，于是 $L((\Delta\cap M)_0)$ 是形如
\[
\sum_{m\in(\Delta\cap M)_0}\lambda_m t^m
\]
的 Laurent 多项式构成的向量空间。由于环面 $T=T_N$ 通过 $v\cdot t^m=v^m t^m$（$v\in T$）作用于 $t^m$，便在 $L((\Delta\cap M)_0)$ 上诱导出作用
\begin{equation}\label{eq:toric-moduli-coefficient-action}
v\cdot\lambda_m=v^m\lambda_m,\qquad v\in T.
\end{equation}
因此，$T$ 作用于 $\mathbb P(L((\Delta\cap M)_0))$；可以验证，在包含映射 $T\subset\operatorname{Aut}(X)$ 下，这与 $\operatorname{Aut}(X)$ 在 $\mathbb P(L(\Delta\cap M))$ 上的作用相容。

现在将简化多项式模空间 $\mathcal M_{\mathrm{simp}}$ 定义为
\[
\mathcal M_{\mathrm{simp}}=\mathbb P(L((\Delta\cap M)_0))/T.
\]
与前面一样，实际上须限制到 $\mathbb P(L((\Delta\cap M)_0))$ 中对应于拟光滑超曲面的子集，并考虑商的存在性。幸运的是，后一问题不太棘手，因为 $T$ 是约化群，而式~\eqref{eq:toric-moduli-coefficient-action} 给出了作用的线性化。因此，不仅能够构造 $\mathcal M_{\mathrm{simp}}$，还可用 GIT 商将它紧化 \SourceCite{Mumford1}：
\[
\overline{\mathcal M}_{\mathrm{simp}}=\mathbb P(L((\Delta\cap M)_0))/\!/T.
\]
不过，对我们的目的而言，最有趣的紧化来自 \SourceCite{KSZ1} 的 \emph{Chow 商}。后面将看到，这种紧化具有由次级扇给出的自然环面结构。

基本断言是：显然定义的映射 $\phi:\mathcal M_{\mathrm{simp}}\to\mathcal M_{\mathrm{poly}}$ 至少在一般位置上是局部同构。这就是 \SourceCite{AGM1} 中的\emph{占优猜想}，现在加以证明。
\begin{proposition}\label{prop:simplified-moduli-etale}\index{简化多项式模空间|hyperpage}
映射 $\phi:\mathcal M_{\mathrm{simp}}\to\mathcal M_{\mathrm{poly}}$ 在一般位置上为平展映射。
\end{proposition}
\begin{proof}
用 $\operatorname{Aut}_0(X)$ 表示 $\operatorname{Aut}(X)$ 的单位元所在连通分支。如 \SourceCite{Cox} 所指出的，$\operatorname{Aut}_0(X)$ 在 $\operatorname{Aut}(X)$ 中具有有限指数，因此只需证明自然映射
\[
\mathbb P(L((\Delta\cap M)_0))/T\longrightarrow
\mathbb P(L(\Delta\cap M))/\operatorname{Aut}_0(X)
\]
为占优映射。

令 $I=(\Delta\cap M)-(\Delta\cap M)_0$，即 $\Delta$ 的刻面内部格点所组成的集合。由命题~\ref{prop:automorphism-facets} 的证明，每个 $m\in I$ 确定一个自同构的单参数族 $y_{m,u}:X\to X$，其中参数 $u\in\mathbb C$。为描述 $y_{m,u}$，使用\hyperref[subsec:homogeneous-ring]{第 3.2.3 小节}的齐次坐标环 $S=\mathbb C[x_\rho]$。由于 $m$ 位于某个刻面的内部，存在 $\rho$，使 $\langle m,v_\rho\rangle=-1$，而对所有 $\rho'\ne\rho$ 都有 $\langle m,v_{\rho'}\rangle>-1$。置
$\mathbf x^D=\prod_{\rho'\ne\rho}x_{\rho'}^{\langle m,v_{\rho'}\rangle}$，相应的自同构单参数族为
\[
y_{m,u}(x_\rho,\mathbf x)=(x_\rho+u\mathbf x^D,\mathbf x).
\]
\SourcePage{141}
这里与式~\eqref{eq:toric-root-automorphism} 一样，将 $\mathbb C^{\Sigma(1)}$ 的点写成 $(x_\rho,\mathbf x)$，其中 $\mathbf x$ 是以 $\Sigma(1)-\{\rho\}$ 为指标的向量。

考虑具有 $|I|$ 个参数的自同构族
\begin{equation}\label{eq:toric-root-product}
y_{\mathbf u}=\prod_{m_i\in I}y_{m_i,u_i}.
\end{equation}
为定义式~\eqref{eq:toric-root-product} 的右端，需要给 $I$ 选定一个次序，因为 $y_{m_i,u_i}$ 未必彼此交换。将包含映射 $T\subset\operatorname{Aut}_0(X)$ 与映射 $\mathbf u\mapsto y_{\mathbf u}$ 结合，得到映射
\[
T\times\mathbb C^I\longrightarrow\operatorname{Aut}_0(X),
\]
很容易看出它为占优映射。因此，只需考虑 $T\times\mathbb C^I$ 在 $\mathbb P(L(\Delta\cap M))$ 上的作用；只要证明在一般位置上，$\mathbb C^I$ 的轨道维数为 $|I|$，且与 $\mathbb P(L((\Delta\cap M)_0))$ 横截相交，命题就成立。

先计算 $y_{m,u}$ 在 Laurent 单项式 $t^\alpha\in L(\Delta\cap M)$ 上的作用。为此，需将 $t^\alpha$ 映到单项式 $\prod_\rho x_\rho^{\langle\alpha,v_\rho\rangle+1}$（见\hyperref[subsec:homogeneous-ring]{第 3.2.3 小节}），施以 $y_{m,u}$，再映回 $L(\Delta\cap M)$。若 $\langle m,v_\rho\rangle=-1$ 确定 $m$ 所在的刻面，则计算得
\[
y_{m,u}(t^\alpha)=
\begin{cases}
t^\alpha,&\langle\alpha,v_\rho\rangle=-1,\\
t^\alpha+u(\langle\alpha,v_\rho\rangle+1)t^{\alpha+m}
+\text{高阶项},&\langle\alpha,v_\rho\rangle>-1.
\end{cases}
\]
可以验证，$\langle\alpha,v_\rho\rangle>-1$ 蕴含 $\alpha+m\in\Delta\cap M$。\hyperref[ex:degree-four-moduli]{例题~\ref*{ex:degree-four-moduli}} 将给出一次具体计算，包括上面省略的高阶项。

由于 $\Delta$ 自反，原点 $0$ 属于 $(\Delta\cap M)_0$，所以 $1=t^0\in L((\Delta\cap M)_0)$。于是很容易得到
\[
y_{\mathbf u}(1)=1+\sum_{m_i\in I}u_i t^{m_i}+\text{高阶项}.
\]
这说明，$\mathbb C^I$ 作用于 $1$ 所得轨道的维数正确；而且它在 $1$ 处的切空间显然与 $L((\Delta\cap M)_0)$ 互补。由于 $y_{\mathbf u}(1)$ 的常数项恒为 $1$，它恒不为零；射影化后上述一阶结论仍然成立，命题随即得证。
\end{proof}

这一命题表明，在一般位置上，$\mathcal M_{\mathrm{simp}}$ 是 $\mathcal M_{\mathrm{poly}}$ 的有限覆盖。了解 $X$ 的自同构群后，就能显式计算作用于 $\mathcal M_{\mathrm{simp}}$ 的有限群。\SourceCite{CFKM} 已在一个例子中完成这一计算；后面的\hyperref[ex:degree-four-moduli]{例题~\ref*{ex:degree-four-moduli}} 将在稍简单的情形中说明这种技巧。

接着计算 $\mathcal M_{\mathrm{poly}}$ 与整个复结构模空间 $\mathcal M$ 的维数，并确定两者何时相等。
\begin{proposition}\label{prop:toric-moduli-dimensions}
设 $V\subset X=X_\Sigma$ 为极小 Calabi--Yau 环面超曲面，其中 $\Sigma$ 是 $n$ 维自反多胞形 $\Delta$ 的法扇的一个极大射影细分。则：
\begin{enumerate}[label=(\roman*)]
\item $V$ 的多项式模空间 $\mathcal M_{\mathrm{poly}}$ 的维数为
\[
\dim\mathcal M_{\mathrm{poly}}=l(\Delta)-n-1-\sum_\Gamma l^*(\Gamma),
\]
其中 $l(\Delta)=|\Delta\cap M|$，求和遍及 $\Delta$ 的所有刻面 $\Gamma$，$l^*(\Gamma)$ 为 $\Gamma$ 相对内部的格点数。
\SourcePage{142}
\item $V$ 的复结构模空间 $\mathcal M$ 的维数为
\[
\dim\mathcal M=l(\Delta)-n-1-\sum_\Gamma l^*(\Gamma)
+\sum_\Theta l^*(\Theta)l^*(\widehat\Theta^\circ),
\]
其中第二个求和遍及 $\Delta$ 的所有余维二的面，$\widehat\Theta^\circ$ 为 $\Delta^\circ$ 中与 $\Theta$ 对偶的一维面。
\item 若 $\Delta\cap M$ 中没有余维二的面的内部格点，则 $\mathcal M=\mathcal M_{\mathrm{poly}}$。
\end{enumerate}
\end{proposition}
\begin{proof}
$\mathcal M_{\mathrm{simp}}$ 的维数很容易计算：由于 $\Delta$ 的维数为 $n$，$T$ 在 $\mathbb P(L((\Delta\cap M)_0))$ 上的作用，在一般位置上具有有限稳定子。因此
\[
\dim\mathcal M_{\mathrm{simp}}=|(\Delta\cap M)_0|-n-1,
\]
这也就是 $\mathcal M_{\mathrm{poly}}$ 的维数。由 $l(\Delta)=|(\Delta\cap M)_0|+\sum_\Gamma l^*(\Gamma)$，得到命题的第一部分。

对于第二部分，注意 $V$ 为极小 Calabi--Yau 轨形，所以由命题~\ref{prop:minimal-cy-orbifold-deformations}，它的无穷小形变由 $H^1(V,\Theta_V)\simeq H^1(V,\widehat\Omega_V^{n-2})$ 分类（记住 $V$ 的维数为 $n-1$）。此外，由命题~\ref{prop:terminal-gorenstein-orbifold-codimension}，$V$ 的奇点集余维至少为 $4$，所以 Ran 的结果 \SourceCite{Ran2} 表明这些形变无障碍。于是 $\mathcal M$ 的维数为 $h^{n-2,1}(V)$。式~\eqref{eq:batyrev-hn21} 计算了这个 Hodge 数，从而证明了第二部分。最后一部分直接由前两部分推出，命题得证。
\end{proof}

这一命题与\hyperref[chap:mirror-constructions]{第 4 章}的定理~\ref{thm:batyrev-hodge} 有很好的联系。在一般点 $V$ 处，$\mathcal M$ 的切空间为 $H^1(V,\Theta_V)\simeq H^{n-2,1}(V)$；\hyperref[sec:batyrev-mirror]{第 4.1 节}中记为 $H_{\mathrm{poly}}^{n-2,1}(V)\subset H^{n-2,1}(V)$ 的空间，就是 $\mathcal M_{\mathrm{poly}}$ 在 $V$ 处的切空间。因此，命题~\ref{prop:toric-moduli-dimensions} 的第一部分给出了式~\eqref{eq:polynomial-hodge-number} 的严格证明。

现在对 $\mathcal M_{\mathrm{poly}}$ 有了更多了解，下一项任务就是将它紧化，并考察能否找到一些\IndexLocation{indexloc:0766:142}最大幺幂边界点。实际上，将使用 $\mathcal M_{\mathrm{simp}}$，因为它有一个特别好的紧化。如前所述，采用 \SourceCite{KSZ1} 的 Chow 商，并取其正规化：
\begin{equation}\label{eq:moduli-chow-quotient}\index{Chow 商!正规化|hyperpage}
\overline{\mathcal M}_{\mathrm{simp}}=\mathbb P(L((\Delta\cap M)_0))/\!/_{\mathrm{Ch}}T.
\end{equation}
它是一个环面簇，先描述其中的环面。注意，$L((\Delta\cap M)_0)$ 包含环面 $(\mathbb C^*)^{(\Delta\cap M)_0}$。位似给出自然映射 $\mathbb C^*\to(\mathbb C^*)^{(\Delta\cap M)_0}$（射影化时需要它）；而式~\eqref{eq:toric-moduli-coefficient-action} 表明，$T$ 的作用来自由特征标 $t^m$（$m\in(\Delta\cap M)_0$）给出的映射 $T\to(\mathbb C^*)^{(\Delta\cap M)_0}$。所以有环面
\begin{equation}\label{eq:moduli-quotient-torus}
T_0=(\mathbb C^*)^{(\Delta\cap M)_0}/\operatorname{im}(T\times\mathbb C^*),
\end{equation}
它包含在上述商中。利用特殊点 $0\in(\Delta\cap M)_0$，可以去掉 $\mathbb C^*$。具体地，令 $\Xi=(\Delta\cap M)_0-\{0\}$，则
\[
T_0=(\mathbb C^*)^\Xi/\operatorname{im}(T).
\]
$T_0$ 的单参数子群构成格
\begin{equation}\label{eq:moduli-quotient-lattice}
N_0=(\mathbb Z^\Xi/N)/\text{挠元},
\end{equation}
其中当 $\Xi=\{m_1,\ldots,m_s\}$ 时，包含映射 $N\to\mathbb Z^\Xi$ 为 $u\mapsto(\langle u,m_1\rangle,\ldots,\langle u,m_s\rangle)$。

\SourcePage{143}
格 $N_0$ 很重要，因为\hyperref[sec:gkz-decomposition]{第 3.4 节}定义的 $\Xi=(\Delta\cap M)_0-\{0\}$ 的 GKZ 分解，自然位于向量空间 $N_0\otimes\mathbb R$ 中。此外，由定理~\ref{thm:gkz-fan}，该扇的极大锥是所有 $M_{\mathbb R}$ 中射影、单纯且满足 $\Sigma^\circ(1)\subset\Xi$ 的扇 $\Sigma^\circ$ 对应的锥 $\operatorname{cpl}(\Sigma^\circ)$。

这与 $V$ 的 Batyrev 镜像密切相关。回顾这一构造：若 $\Delta^\circ\subset N_{\mathbb R}$ 为极多胞形，$\Sigma^\circ$ 为 $\Delta^\circ$ 的法扇的一个极大射影细分，那么 Batyrev 镜像就是反典范超曲面族 $V^\circ\subset X_{\Sigma^\circ}$。按定义~\ref{def:projective-subdivision}，$\Sigma^\circ$ 射影、单纯，细分 $\Delta^\circ$ 的法扇，并满足 $\Sigma^\circ(1)=\Delta\cap M-\{0\}$。不过，在\hyperref[subsec:enlarged-kahler-moduli]{第 6.2.3 小节}将看到，可以将最后一个条件换成
$\Sigma^\circ(1)=(\Delta\cap M)_0-\{0\}=\Xi$，而不改变 $V^\circ$。这样，$V^\circ\subset X_{\Sigma^\circ}$ 就给出 $\Xi$ 的 GKZ 分解中的一个锥 $\operatorname{cpl}(\Sigma^\circ)$。

因此，每一种 Batyrev 镜像 $V^\circ\subset X_{\Sigma^\circ}$ 的选择，都在 GKZ 分解中给出一个特殊极大锥。在对应于 GKZ 分解的环面簇中，这些锥给出环面边界上的特殊点。它们可能是奇异点；不过，很快将看到，猜想其消解能够给出所需的最大幺幂边界点。此外，作为\hyperref[subsec:enlarged-kahler-moduli]{第 6.2.3 小节}的预告，还应指出：$\operatorname{cpl}(\Sigma^\circ)$ 的内部就是 $X_{\Sigma^\circ}$ 的 Kähler 锥，并且自然同构于 $V^\circ$ 的 Kähler 锥的环面部分。这是 $V$ 的复结构模与其镜像 $V^\circ$ 的 Kähler 模之间同构关系的开端。

现在可以仔细描述 Chow 商~\eqref{eq:moduli-chow-quotient}。定理~\ref{thm:gkz-fan} 中 $\Xi$ 的 GKZ 分解并不完备，但如\hyperref[sec:gkz-decomposition]{第 3.4 节}所述，可以自然扩充为一个完备扇，称为 $\Xi$ 的\emph{次级扇}。为便于后面使用，注意次级扇由以下点集构造：
\begin{equation}\label{eq:moduli-secondary-configuration}\index{次级扇|hyperpage}
\Xi^+=(\Xi\cup\{0\})\times\{1\}=(\Delta\cap M)_0\times\{1\}\subset M\times\mathbb Z.
\end{equation}
\SourceCite{KSZ2} 表明，正规化的 Chow 商 $\overline{\mathcal M}_{\mathrm{simp}}$ 就是由 $\Xi$ 的次级扇确定的环面簇。\SourceCite{KSZ1} 探究了 Chow 商与各种 GIT 商之间的关系。还应指出，$\overline{\mathcal M}_{\mathrm{simp}}$ 有一个由 \SourceCite{GKZ2} 的\emph{\IndexLocation{indexloc:1242:143}次级多胞形}给出的自然射影嵌入。

环面紧化 $\overline{\mathcal M}_{\mathrm{simp}}$ 包含一些特殊的环面边界点，对应于镜像族的不同选择。不过，$\overline{\mathcal M}_{\mathrm{simp}}$ 不一定光滑，因此讨论最大幺幂边界点以前，需要先消解奇点。于是选取次级扇的一个细分，使它给出光滑环面簇。所得 $\mathcal M_{\mathrm{simp}}$ 的光滑紧化记为 $\widetilde{\mathcal M}_{\mathrm{simp}}$。

这一光滑紧化具有如下性质：每个 Batyrev 镜像 $V^\circ\subset X_{\Sigma^\circ}$ 都在 $\widetilde{\mathcal M}_{\mathrm{simp}}$ 中给出有限多个特殊边界点，对应于细分 $\operatorname{cpl}(\Sigma^\circ)$ 的锥。猜想这些正是所寻求的最大幺幂边界点。
\begin{conjecture}\label{conj:toric-mum-boundary}
设 $\widetilde{\mathcal M}_{\mathrm{simp}}$ 是 Chow 商 $\overline{\mathcal M}_{\mathrm{simp}}$ 的一个光滑环面消解，$V^\circ\subset X_{\Sigma^\circ}$ 是 $V$ 的一个 Batyrev 镜像。则 $\widetilde{\mathcal M}_{\mathrm{simp}}$ 中对应于细分 $\operatorname{cpl}(\Sigma^\circ)$ 的锥的点，全部都是最大幺幂边界点。此外，这些点在定义~\ref{def:equivalent-mum-points} 的等价关系下全部等价。
\end{conjecture}
\SourceCite{BK1} 曾给出这一猜想的一个不完整证明。若 $H^{1,1}(V)$ 在杯积下生成 $\bigoplus_pH^{p,p}(V)$，那么只要证明
\SourcePage{144}
式~\eqref{eq:toric-complete-intersection-i-function} 定义的形式函数 $I_V$ 的系数确实是镜像族 $V^\circ$ 的周期，正如 \SourceCite{Givental4} 所断言的（见\hyperref[subsec:nef-gkz]{第 5.5.3 小节}的讨论），就很容易推出这一猜想。特别是，这适用于 $V$ 为三维簇的情形。可以证明，猜想~\ref{conj:toric-mum-boundary} 对后面的\hyperref[ex:weighted-complex-moduli]{例题~\ref*{ex:weighted-complex-moduli}} 和\hyperref[ex:degree-four-moduli]{例题~\ref*{ex:degree-four-moduli}} 成立。例如，前一个例题的证明可由论文 \SourceCite{CdFKM} 中的计算推出。

猜想~\ref{conj:toric-mum-boundary} 的动机之一是：在镜像对称下，绕某个边界除子的单值性的对数，应与同细分 $\operatorname{cpl}(\Sigma^\circ)$ 的锥的相应生成元的负元作杯积识别起来；最大幺幂边界点的定义~\ref{def:maximal-unipotent-point} 后面的讨论已经提到这一点。\hyperref[chap:quantum-cohomology]{第~\ref*{chap:quantum-cohomology} 章}讨论 A 模型 Hodge 结构变分时，将从更深入的角度理解它。

\emph{本章余下部分都假定猜想~\ref{conj:toric-mum-boundary} 成立。}

猜想~\ref{conj:toric-mum-boundary} 描述的最大幺幂边界点来自次级扇中一些十分特殊的锥。GKZ 分解中还有其他锥：它们对应的 $\Sigma^\circ$ 不使用 $\Xi$ 的全部点，或者不细分 $\Delta^\circ$ 的法扇；将 GKZ 分解扩充为次级扇时，还会得到一些锥。猜想~\ref{conj:toric-mum-boundary} 没有说明 $\overline{\mathcal M}_{\mathrm{simp}}$ 中对应于这些锥的边界点的性质。不过，在\hyperref[subsec:enlarged-kahler-moduli]{第 6.2.3 小节}将看到，所有这些锥都有有趣的物理解释——每个锥都给出一种不同的物理理论。

还应提到，\emph{Gröbner 扇}给出了次级扇的一个典范细分 \SourceCite{MR,Sturmfels1}。它给出 $\mathcal M_{\mathrm{simp}}$ 的一个紧化，是 $\overline{\mathcal M}_{\mathrm{simp}}$ 的爆破，因而奇点可能更少。\SourceCite{HLY1,HLY2} 使用了这一紧化，并证明存在\emph{最大退化}边界点，即在局部存在唯一的全纯周期积分，常数倍除外。这是最大幺幂性的一个必要条件，因此应由猜想~\ref{conj:toric-mum-boundary} 推出。

目前的讨论着重于爆破 $\overline{\mathcal M}_{\mathrm{simp}}$，似乎 $\overline{\mathcal M}_{\mathrm{simp}}$ 是 $\mathcal M_{\mathrm{simp}}$ 的“极小”紧化。不过，在某些情况下，将次级扇中的若干锥合并成一个较大的锥是完全自然的。这会给出一个紧化 $\overline{\mathcal M}'_{\mathrm{simp}}$，使 $\overline{\mathcal M}_{\mathrm{simp}}$ 是 $\overline{\mathcal M}'_{\mathrm{simp}}$ 的爆破。\hyperref[ex:degree-four-moduli]{例题~\ref*{ex:degree-four-moduli}} 将展示这种现象；\hyperref[subsec:enlarged-kahler-moduli]{第 6.2.3 小节}将仔细研究应如何以及何时这样做。

还要注意，除了存在多个镜像，原来的族 $V\subset X_\Sigma$ 也有多个版本，因为极大射影细分 $\Sigma$ 可能有许多种选择。我们对 $\mathcal M_{\mathrm{simp}}$ 的描述没有提到 $\Sigma$，原因很明确：这些族具有相同的多项式模。变化的是 $V$ 的 Kähler 锥，因为每个 $\Sigma$ 都对应于\hyperref[subsec:enlarged-kahler-moduli]{第 6.2.3 小节}所讨论的 $V$ 的扩大 Kähler 模空间的一部分。\hyperref[chap:mirror-constructions]{第~\ref*{chap:mirror-constructions} 章}中，$\Sigma$ 和 $\Sigma^\circ$ 的多种选择似乎是个麻烦，因为它们使我们无法典范地选择镜像。不过，现在逐渐了解到：“多个镜像”的现象，对从整体上理解复结构模与 Kähler 模至关重要。

紧化 $\overline{\mathcal M}_{\mathrm{simp}}$ 的环面边界包含刚才讨论的最大幺幂边界点。不过，$\overline{\mathcal M}_{\mathrm{simp}}-\mathcal M_{\mathrm{simp}}$ 中也有不在环面边界上的点，因为环面 $T_0\subset\overline{\mathcal M}_{\mathrm{simp}}$ 中有许多点对应于奇异超曲面 $V\subset X_\Sigma$。实际上，由于 $V$ 和 $X_\Sigma$ 可能
\SourcePage{145}
有奇点，通常用\emph{非退化超曲面}来表述更方便。对 $V\subset X_\Sigma$，这意味着：对每个环面轨道 $\mathcal O\subset X_\Sigma$，交集 $V\cap\mathcal O$ 或者为空，或者光滑且余维为一。由于 $X_\Sigma$ 是轨形，很容易推出 $X_\Sigma$ 的每个非退化超曲面都是子轨形。

在 $L((\Delta\cap M)_0)$ 中，\emph{判别式轨迹} $\mathcal D$ 是所有使 $f=0$ 不是非退化超曲面的 Laurent 多项式 $f$ 构成的集合。由 \SourceCite[第 10 章]{GKZ2}，$\mathcal D$ 的定义方程由\emph{主 $\mathcal A$ 行列式}给出。置 $\mathcal A=(\Delta\cap M)_0=\Xi\cup\{0\}$，则它是 $f$ 的系数的多项式，定义为
\[
E_{\mathcal A}(f)=\operatorname{Res}_{\mathcal A}
\left(t_1\frac{\partial f}{\partial t_1},\ldots,t_n\frac{\partial f}{\partial t_n},f\right),
\]
其中 $\operatorname{Res}_{\mathcal A}$ 为 \SourceCite[第 8 章]{GKZ2} 定义的稀疏结式，$t_1,\ldots,t_n$ 为环面 $T$ 的坐标。这个多项式实际上可约，并具有因式分解
\[
E_{\mathcal A}(f)=\prod_{\Gamma\subset\Delta}
D_{\mathcal A\cap\Gamma}(f_\Gamma)^{m(\Gamma)}.
\]
右端可如下理解。$\mathcal A$ 判别式 $D_{\mathcal A}$ 定义了 $L((\Delta\cap M)_0)$ 中那些在环面 $T$ 上具有奇异点的多项式所组成集合的 Zariski 闭包（因此，若该集合的余维大于 $1$，则 $D_{\mathcal A}=1$）。对每个面 $\Gamma\subset\Delta$，再将 $f_\Gamma$ 定义为 $f$ 中对应于 $\Gamma\cap\mathcal A$ 的各项之和。最后，上式的指数 $m(\Gamma)$ 是正整数，在良好情形下等于 $1$。

由这一描述可知，判别式轨迹可能有许多连通分支。幸运的是，许多 $D_{\mathcal A\cap\Gamma}$ 常常等于 $1$。例如，后面的\hyperref[ex:weighted-complex-moduli]{例题~\ref*{ex:weighted-complex-moduli}} 中，判别式轨迹恰好有两个分支。还应指出，$\mathcal D$ 恰好也是下面将描述的 GKZ 超几何方程组的判别式轨迹。

从 $L((\Delta\cap M)_0)$ 下降到 Chow 商时，判别式轨迹的像仍记为 $\mathcal D$。判别式的\emph{主分支}是由 $D_{\mathcal A}=0$ 定义的超曲面。当 $V$ 为三维簇时，该分支的一般点参数化具有一个结点的超曲面 $V$（三维结点在物理文献中常称为\emph{锥形奇点}）。判别式的其他非空分支，则参数化带有可能更复杂的奇点的超曲面。

除刚才描述的判别式轨迹外，$\overline{\mathcal M}_{\mathrm{simp}}$ 的环面边界也参数化奇异超曲面，尽管参数值 $\lambda_i$ 可能没有良好定义。从 GIT 商的角度看，问题在于这些参数值可能对应于非稳定超曲面。因此，作为 $\overline{\mathcal M}_{\mathrm{simp}}$ 中的对象，超曲面本身可能没有良好定义。例如，五次镜像族~\eqref{eq:quintic-pencil} 在 $\psi\to\infty$ 时的极限，可以取为 $x_1x_2x_3x_4x_5=0$；然而，进行坐标变换 $x_5\mapsto x_5/\psi$ 后，极限却为
\[
x_1^5+x_2^5+x_3^5+x_4^5+x_1x_2x_3x_4x_5=0.
\]

\subsection{超几何方程与模空间}\label{subsec:hypergeometric-moduli}\index{GKZ方程组@GKZ 方程组!模空间|hyperpage}
\hyperref[sec:hypergeometric]{第 5.5 节}定义了 GKZ 超几何方程，并且以相当朴素的方式讨论了它们在 $V$ 的模空间上的行为。现在，对复结构模有了更多了解，可以更清楚地认识 GKZ 方程组。

\hyperref[sec:hypergeometric]{第 5.5 节}构造超几何方程时使用了 $\Delta\cap M$。这里，由于要在模空间 $\overline{\mathcal M}_{\mathrm{simp}}$ 上工作，将使用 $(\Delta\cap M)_0$
\SourcePage{146}
来代替它。按照式~\eqref{eq:toric-gkz-configuration}，这意味着使用从式~\eqref{eq:moduli-secondary-configuration} 中的点集
\[
(\Delta\cap M)_0\times\{1\}=\Xi^+
\]
构造的 GKZ 方程组。与\hyperref[chap:hodge-yukawa]{第~\ref*{chap:hodge-yukawa} 章}一样，将 $(\Delta\cap M)_0$ 的点记为 $m_i$，其中 $m_0=0$。于是 $\Xi^+$ 中各点的关系模为
\begin{equation}\label{eq:moduli-relation-lattice}
M_0=\left\{\ell=(\ell_i)\in\mathbb Z^{\Xi^+}:\sum_i\ell_i m_i=0,\ \sum_i\ell_i=0\right\}.
\end{equation}
由式~\eqref{eq:moduli-quotient-torus}，知道 $M_0$ 可识别为环面 $T_0$ 的特征标群，所以式~\eqref{eq:moduli-quotient-lattice} 中的 $N_0$ 是 $M_0$ 的对偶。GKZ 超几何方程组由算子 $\Box_\ell$（$\ell\in M_0$）及 $Z_j$（$j=1,\ldots,n+1$）组成。$\Box_\ell,Z_j$ 的精确定义见\hyperref[sec:hypergeometric]{第 5.5 节}。

\hyperref[chap:hodge-yukawa]{第~\ref*{chap:hodge-yukawa} 章}用 $(\mathbb C^*)^{\Delta\cap M}/\operatorname{im}(T\times\mathbb C^*)$ 粗略近似模空间。子空间 $(\mathbb C^*)^{(\Delta\cap M)_0}/\operatorname{im}(T\times\mathbb C^*)$ 恰好就是环面 $T_0\subset\overline{\mathcal M}_{\mathrm{simp}}$。$f\in L((\Delta\cap M)_0)$ 在由 $T_0$ 参数化的族上给出相对 $n$ 形式
\[
\widetilde\omega=\frac{\lambda_0}{f}\frac{dt_1}{t_1}\wedge\cdots\wedge\frac{dt_n}{t_n},
\]
已经证明它满足方程 $\Box_\ell\lambda_0^{-1}\widetilde\omega=0$ 及 $Z_j\lambda_0^{-1}\widetilde\omega=0$。后一个方程只是表达 $\widetilde\omega$ 在 $T$ 和 $\mathbb C^*$ 的作用下不变（这是因为 $\widetilde\omega$ 中额外加入了因子 $\lambda_0$），所以它下降为 $\overline{\mathcal M}_{\mathrm{simp}}$ 所参数化的族上的有理相对 $n$ 形式。因此，$\overline{\mathcal M}_{\mathrm{simp}}$ 上的 GKZ 方程组为
\[
\Box_\ell\lambda_0^{-1}\widetilde\omega=0,\qquad\ell\in M_0.
\]
为使用这些方程，需要以 $\overline{\mathcal M}_{\mathrm{simp}}$ 上的局部坐标表达它们。每个 $\ell=(\ell_i)\in M_0$ 给出
\begin{equation}\label{eq:moduli-monomial-coordinate}
z=\prod_i\lambda_i^{\ell_i},
\end{equation}
它在 $T\times\mathbb C^*$ 下不变，所以选取 $M_0$ 的一组基，就给出环面 $T_0$ 的坐标 $z_1,\ldots,z_r$。\hyperref[sec:hypergeometric]{第 5.5 节}的方法于是能将方程 $\Box_\ell\lambda_0^{-1}\widetilde\omega=0$ 写成 $z_i$ 的方程，不过该节没有指定应使用 $M_0$ 的哪组基。

现在，在猜想~\ref{conj:toric-mum-boundary} 的前提下，已经确切知道如何选择坐标。固定 $\overline{\mathcal M}_{\mathrm{simp}}$ 的一个光滑环面消解 $\widetilde{\mathcal M}_{\mathrm{simp}}$，令 $\sigma$ 为消解的扇中的一个极大锥，它细分来自 $V$ 的某个 Batyrev 镜像 $V^\circ\subset X_{\Sigma^\circ}$ 的锥 $\operatorname{cpl}(\Sigma^\circ)$。由于 $\widetilde{\mathcal M}_{\mathrm{simp}}$ 光滑，$\sigma$ 的生成元构成 $N_0$ 的一组基；$M_0$ 中的对偶基给出坐标 $z_1,\ldots,z_r$，它们是 $\widetilde{\mathcal M}_{\mathrm{simp}}$ 中对应于 $\sigma$ 的点的局部坐标。

后面的\hyperref[subsec:two-parameter-moduli]{第 6.1.4 小节}将给出一个很好的例子。还应指出，上述特殊坐标 $z_1,\ldots,z_r$ 是镜像映射的一部分；\hyperref[sec:mirror-map]{第~\ref*{sec:mirror-map} 节}将进一步详细讨论。

\subsection{例题}\label{subsec:two-parameter-moduli}
现在给出两个 Calabi--Yau \IndexLocation{indexloc:0099:146}环面超曲面复结构模的例子。
\SourceExampleNumber{6.1.4}{1}
\begin{example}\label{ex:weighted-complex-moduli}
首先考虑 $\mathbb P(1,1,2,2,2)$ 的环面爆破中反典范\IndexLocation{indexloc:1020:146}超曲面的镜像族，前面
\SourcePage{147}
\hyperref[sec:hodge-examples]{第 5.4 节}、\hyperref[sec:hypergeometric]{第 5.5 节}及\hyperref[sec:yukawa-couplings]{第 5.6 节}已研究过它。沿用\hyperref[ex:weighted-pf-method]{例题~\ref*{ex:weighted-pf-method}} 的记号，自反多胞形 $\Delta^\circ$ 的顶点为
\[
(-1,-2,-2,-2),\ (1,0,0,0),\ (0,1,0,0),\ (0,0,1,0),\ (0,0,0,1).
\]
要研究 $V^\circ\subset X_{\Sigma^\circ}$ 的复结构模，其中 $\Sigma^\circ$ 是 $\Delta^\circ$ 的法扇的一个极大射影细分。

先证明 $\overline{\mathcal M}=\overline{\mathcal M}_{\mathrm{poly}}=\overline{\mathcal M}_{\mathrm{simp}}$。\hyperref[ex:weighted-pf-method]{例题~\ref*{ex:weighted-pf-method}} 指出，$\Delta^\circ\cap N$ 由 $\Delta^\circ$ 的五个顶点，以及
$(0,-1,-1,-1)=\frac12(-1,-2,-2,-2)+\frac12(1,0,0,0)$ 和 $(0,0,0,0)$ 组成。这给出 $\Delta^\circ$ 的一条具有内部格点的边；与之对偶的 $\Delta$ 的二维面没有内部格点。$\Delta^\circ$ 的余维二的面也没有内部格点，所以由命题~\ref{prop:toric-moduli-dimensions}，$\overline{\mathcal M}=\overline{\mathcal M}_{\mathrm{poly}}$。

还注意到，$\Delta^\circ$ 的刻面没有内部格点。因此，由命题~\ref{prop:automorphism-facets}，$\operatorname{Aut}(X_{\Sigma^\circ})$ 没有根；这意味着 $X_{\Sigma^\circ}$ 的环面 $T^\circ$ 就是自同构群的单位元所在连通分支。不过，由 $\Delta^\circ$ 的法扇的自同构，得到
$\operatorname{Aut}(X_{\Sigma^\circ})/T^\circ\simeq\mathbb Z/2\mathbb Z\times S_3$
（交换 $\Delta^\circ$ 的前两个顶点给出 $\mathbb Z/2\mathbb Z$，置换其余三个顶点给出 $S_3$）。需证明这些自同构在 $\overline{\mathcal M}_{\mathrm{simp}}$ 上平凡作用。注意 $f\in L((\Delta^\circ\cap N)_0)$ 可写为
\[
f=\lambda_1t_1^{-1}t_2^{-2}t_3^{-2}t_4^{-2}+\lambda_2t_1
+\lambda_3t_2+\lambda_4t_3+\lambda_5t_4
+\lambda_6t_2^{-1}t_3^{-1}t_4^{-1}+\lambda_0.
\]
交换 $t_2,t_3$ 的自同构（$S_3$ 中的一个）在 $L((\Delta^\circ\cap N)_0)$ 上诱导一个交换 $\lambda_3,\lambda_4$ 的自同构。不过，可以验证，$(1,\lambda_4/\lambda_3,\lambda_3/\lambda_4,1)\in T^\circ$ 在 $f$ 上具有相同效果，所以这个自同构在 $(T^\circ\times\mathbb C^*)$ 轨道上平凡作用。其余自同构类似处理，因而 $\overline{\mathcal M}_{\mathrm{poly}}=\overline{\mathcal M}_{\mathrm{simp}}$。

下一步构造次级扇，它由 $\Xi^+=(\Delta^\circ\cap N)_0\times\{1\}$ 构造。因此
\[
\begin{split}
\Xi^+=\{&(-1,-2,-2,-2,1),\ (1,0,0,0,1),\ (0,1,0,0,1),\\
&(0,0,1,0,1),\ (0,0,0,1,1),\ (0,-1,-1,-1,1),\ (0,0,0,0,1)\}.
\end{split}
\]
由式~\eqref{eq:moduli-relation-lattice}，$M_0$ 为 $\Xi^+$ 各点的\IndexLocation{indexloc:0701:147}线性关系格。以下矩阵的各行给出一组基：
\begin{equation}\label{eq:weighted-moduli-relations}
\begin{pmatrix}0&0&1&1&1&1&-4\\1&1&0&0&0&-2&0\end{pmatrix}.
\end{equation}
也得到 $N_0$ 的对偶基。为描述次级扇，使用\hyperref[sec:gkz-decomposition]{第 3.4 节}：它表明，在秩为二的情形中，次级扇是 $N_0\otimes\mathbb R$ 中的完备扇，其一维锥由矩阵~\eqref{eq:weighted-moduli-relations} 的各列生成。因此，在 GKZ 分解~\eqref{eq:weighted-gkz-picture} 中加入由 $(-1,0)$ 生成的锥，并补成完备扇，就得到次级扇。

回顾式~\eqref{eq:weighted-gkz-picture}，知道 $V^\circ$ 有唯一的 Batyrev 镜像，即\hyperref[ex:weighted-resolution]{例题~\ref*{ex:weighted-resolution}} 描述的 $\mathbb P(1,1,2,2,2)$ 的环面爆破中的超曲面族。GKZ 分解中对应的锥为第一象限，因此所需局部坐标 $z_1,z_2$ 由矩阵~\eqref{eq:weighted-moduli-relations} 的行给出。利用式~\eqref{eq:moduli-monomial-coordinate}，各行给出
\[
z_1=\frac{\lambda_3\lambda_4\lambda_5\lambda_6}{\lambda_0^4},\qquad
z_2=\frac{\lambda_1\lambda_2}{\lambda_6^2}.
\]
\SourcePage{148}
如\hyperref[sec:hypergeometric]{第 5.5 节}所述，可利用 $T^\circ\times\mathbb C^*$，将一般的 $f\in L((\Delta\cap M)_0)$ 化为
\begin{equation}\label{eq:weighted-moduli-laurent}
f=z_2t_1^{-1}t_2^{-2}t_3^{-2}t_4^{-2}+t_1+z_1t_2+t_3+t_4
+t_2^{-1}t_3^{-1}t_4^{-1}+1.
\end{equation}
因此，式~\eqref{eq:weighted-pf-family} 首次使用的这种写法并非偶然：按猜想~\ref{conj:toric-mum-boundary}，它恰好给出最大幺幂边界点所需的形式。

在环面边界以外，本例中的判别式轨迹 $\mathcal D\subset\overline{\mathcal M}_{\mathrm{simp}}$ 有两个分支：由 $D_{\mathcal A}=0$（$\mathcal A=(\Delta\cap M)_0$）定义的主分支，以及与 $\Delta$ 的边 $\Gamma$ 对应的分支 $D_{\Gamma\cap\mathcal A}=0$，其中
\[
\Gamma\cap M=\{(-1,-2,-2,-2),(1,0,0,0),(0,-1,-1,-1)\}
=\{m_1,m_2,m_6\}.
\]
为确定 $D_{\mathcal A}$，将 $f$ 写为式~\eqref{eq:weighted-moduli-laurent}，并确定使方程 $f=t_i\partial f/\partial t_i=0$ 在 $T^\circ$ 中有解的 $z_1,z_2$ 的值。可以手算完成，结果是
\[
D_{\mathcal A}(f)=512^2z_1^2z_2-(1-256z_1)^2.
\]
再看边 $\Gamma$，有 $f_\Gamma=z_2t^{m_1}+t^{m_2}+t^{m_6}$。由于 $m_1+m_2=2m_6$，$t^{m_1+m_2}=(t^{m_6})^2$，计算得 $D_{\mathcal A\cap\Gamma}(f_\Gamma)=1-4z_2$。Picard--Fuchs 方程~\eqref{eq:weighted-pf-second} 在 $1-4z_2=0$ 时退化，因此它选出了判别式轨迹的一个分支。对于另一个分支 $D_{\mathcal A}=0$，\hyperref[ex:weighted-yukawa]{例题~\ref*{ex:weighted-yukawa}} 中计算的 Yukawa 耦合表明，GKZ 方程组也在该分支上退化。
\end{example}

\SourceExampleNumber{6.1.4}{2}
\begin{example}\label{ex:degree-four-moduli}
考虑自反多胞形 $\Delta$，它是 $\mathbb Z^4=M$ 中以下各点的凸包：
\begin{equation}\label{eq:second-moduli-polytope}
\begin{gathered}
(-1,-2,-3,-7),\ (1,0,0,0),\ (0,1,0,0),\\
(0,0,1,0),\ (0,0,0,1),\ (0,0,-1,-2).
\end{gathered}
\end{equation}
Batyrev 构造给出镜像族 $V\subset X_\Sigma$ 和 $V^\circ\subset X_{\Sigma^\circ}$。这里的目标是描述 $V$ 的复结构模。将在\hyperref[subsec:enlarged-kahler-moduli]{第 6.2.3 小节}考虑 $V^\circ$ 的 Kähler 模。

首先观察到，$\Delta$ 的格点是式~\eqref{eq:second-moduli-polytope} 的六个顶点，加上原点 $(0,0,0,0)$ 和点 $(0,0,0,-1)$；后一点位于由除第五个顶点以外的所有顶点张成的刻面的内部。由于余维二的面没有内部格点，命题~\ref{prop:toric-moduli-dimensions} 表明 $\overline{\mathcal M}=\overline{\mathcal M}_{\mathrm{poly}}$。刻面内部恰好只有一个格点，所以同一命题还表明模空间为二维。

刻面有内部格点，还意味着 $\operatorname{Aut}(X)$ 大于环面 $T$。不过，仍有 $\overline{\mathcal M}_{\mathrm{poly}}=\overline{\mathcal M}_{\mathrm{simp}}$，现在来证明。先描述由 $(0,0,0,-1)$ 产生的额外自同构在空间 $L(\Delta\cap M)$ 上的作用。虽然可以直接计算，但用齐次坐标来做颇有启发。回顾，$X_\Sigma$ 的扇 $\Sigma$ 满足 $\Sigma(1)=\Delta^\circ\cap N-\{0\}$。可以验证 $\Delta^\circ$ 是以下各点的凸包：
\[
\begin{gathered}
(13,-1,-1,-1),\ (1,-1,3,-1),\ (-1,6,-1,-1),\ (-1,-1,3,-1),\\
(-1,-1,-1,-1),\ (-1,-1,-1,1),\ (-1,0,3,-1).
\end{gathered}
\]
$\Delta^\circ$ 还有许多其他格点。为简化问题，回顾超曲面 $V$ 都是 $\mathbb P_\Delta$ 中反典范超曲面的拉回。因此，可以使用与法扇本身对应的齐次坐标
\SourcePage{149}
进行计算。将按上述次序使用齐次坐标 $x_1,\ldots,x_7$。

由\hyperref[sec:toric-automorphisms]{第 3.6 节}，知道 $\operatorname{Aut}(X_\Sigma)$ 唯一的根来自 $(0,0,0,-1)$，因为它是 $\Delta$ 的刻面内部唯一的格点。包含它的刻面与 $\Delta^\circ$ 的顶点 $(-1,-1,-1,1)$ 对偶，因此很容易看出，相应的根为 $(x_6,x_1x_2x_3x_4x_5x_7)$。自同构的相应单参数族为
\[
y_u(x_6)=x_6+u x_1x_2x_3x_4x_5x_7,\qquad y_u(x_i)=x_i\quad(i\ne6).
\]
回顾命题~\ref{prop:automorphism-facets} 之前的讨论：Laurent 单项式 $t^m\in L(\Delta\cap M)$ 在齐次坐标环中给出单项式
\[
t^m=\prod_i x_i^{\langle m,v_i\rangle+1}.
\]
这里 $v_i$ 为上面列出的 $\Delta^\circ$ 的顶点。将 $\Delta\cap M$ 的点记为：式~\eqref{eq:second-moduli-polytope} 中依次列出的点为 $m_1,\ldots,m_6$，$m_7=(0,0,0,-1)$，$m_0=(0,0,0,0)$。于是
\[
\begin{aligned}
t^{m_0}&=x_1x_2x_3x_4x_5x_6x_7,\\
t^{m_1}&=x_4^2x_5^{14},&t^{m_2}&=x_1^{14}x_2^2,\\
t^{m_3}&=x_3^7x_7,&t^{m_4}&=x_2^4x_4^4x_7^4,\\
t^{m_5}&=x_6^2,&t^{m_6}&=x_1^4x_3^4x_5^4,\\
t^{m_7}&=x_1^2x_2^2x_3^2x_4^2x_5^2x_7^2.
\end{aligned}
\]
自同构的单参数族 $y_u$ 在这些单项式上的作用为
\[
\begin{aligned}
y_u(t^{m_0})&=t^{m_0}+u t^{m_7},\\
y_u(t^{m_5})&=t^{m_5}+2u t^{m_0}+u^2t^{m_7},\\
y_u(t^{m_i})&=t^{m_i},\qquad i\ne0,5.
\end{aligned}
\]
现在证明，自然映射 $\overline{\mathcal M}_{\mathrm{simp}}\to\overline{\mathcal M}_{\mathrm{poly}}$ 在一般位置上是一一的。第一步是说明保持 $\Delta$ 的格自同构在 $\overline{\mathcal M}_{\mathrm{simp}}$ 上平凡作用；这可用\hyperref[ex:weighted-complex-moduli]{例题~\ref*{ex:weighted-complex-moduli}} 的方法完成。所以只需考虑 $\operatorname{Aut}(X_\Sigma)$ 的单位元所在连通分支，它由 $T$ 和自同构 $y_u$ 生成。取 $\overline{\mathcal M}_{\mathrm{simp}}$ 中一点，以 $f=\sum_{i=0}^6\lambda_i t^{m_i}$ 表示，并施以 $y_u$。$y_u(f)$ 中 $t^{m_7}$ 的系数为 $u\lambda_0+u^2\lambda_5$。令该系数等于零，就能找到 $\overline{\mathcal M}_{\mathrm{simp}}$ 中与 $f$ 在 $\overline{\mathcal M}_{\mathrm{poly}}$ 中有相同像的所有点。有两个解：$u=0$ 和 $u=-\lambda_0/\lambda_5$。非平凡解 $u=-\lambda_0/\lambda_5$ 给出
\[
y_u(f)=-\lambda_0t^{m_0}+\sum_{i=1}^6\lambda_i t^{m_i}.
\]
不过，它与 $f$ 表示 $\overline{\mathcal M}_{\mathrm{simp}}$ 的同一点，因为可通过 $(-1,-1,-1,-1,-1)\in T\times\mathbb C^*$ 的作用由 $f$ 得到它。因此，$\overline{\mathcal M}_{\mathrm{simp}}=\overline{\mathcal M}_{\mathrm{poly}}$。

所以可用次级扇紧化 $V$ 的复结构模空间；这里次级扇由 $\Xi^+=(\Delta\cap M)_0\times\{1\}$ 构造。为线性关系格 $M_0$ 选取一组基，由以下矩阵的各行给出：
\[
\begin{pmatrix}-2&-2&-4&1&0&7&0\\1&1&2&0&1&-3&-2\end{pmatrix}.
\]
与上一例一样，可从矩阵的各列读出次级扇的一维锥。因此，$\overline{\mathcal M}_{\mathrm{simp}}$ 是与
\SourcePage{150}
下图所示的扇对应的环面簇：
\begin{figure}[H]
\centering
\refstepcounter{equation}\label{eq:second-moduli-secondary-fan}
\makebox[\textwidth]{\makebox[0pt][l]{(\theequation)}\hfill
\raisebox{-.5\height}{\includegraphics[width=0.53975\textwidth]{6-1-secondary-fan.png}}\hfill\mbox{}}
\end{figure}
在\hyperref[subsec:enlarged-kahler-moduli]{第 6.2.3 小节}计算镜像 $V^\circ$ 的 Kähler 模时，将看到 GKZ 分解由锥 $\sigma_1,\sigma_2,\sigma_3$ 组成（记住，GKZ 分解是次级扇的一个子扇）。不过，只有 $\sigma_1,\sigma_2$ 具有 $\operatorname{cpl}(\Sigma^\circ)$ 的形式，其中 $\Sigma^\circ$ 是 $\Delta^\circ$ 的法扇的适当射影细分。因此，只有 $\sigma_1,\sigma_2$ 给出 $\overline{\mathcal M}_{\mathrm{simp}}$ 的最大幺幂边界点。

这两个锥中，只有 $\sigma_2$ 光滑，所以要将理论应用于 $\sigma_1$，必须对它细分。这会在一个消解 $\widetilde{\mathcal M}_{\mathrm{simp}}$ 中给出多个最大幺幂边界点。不过，还有另一种处理办法：将 $\sigma_1,\sigma_2$ \emph{合并}成一个锥 $\sigma=\sigma_1\cup\sigma_2$。将 $\sigma$ 与式~\eqref{eq:second-moduli-secondary-fan} 中其余锥 $\sigma_3,\sigma_4,\sigma_5$ 合用，得到一个爆缩
$\overline{\mathcal M}_{\mathrm{simp}}\to\overline{\mathcal M}'_{\mathrm{simp}}$，而 $\overline{\mathcal M}'_{\mathrm{simp}}$ 中对应于 $\sigma$ 的点 $p$ 光滑。可以证明，$p$ 是最大幺幂边界点；并且在爆破 $\widetilde{\mathcal M}_{\mathrm{simp}}\to\overline{\mathcal M}'_{\mathrm{simp}}$ 下，$\widetilde{\mathcal M}_{\mathrm{simp}}$ 的所有最大幺幂边界点都映到 $p$，在定义~\ref{def:equivalent-mum-points} 的意义下彼此等价。因此，从\IndexLocation{indexloc:0772:150}最大幺幂单值性的角度看，$\overline{\mathcal M}'_{\mathrm{simp}}$ 是最适宜采用的紧化。

\hyperref[subsec:enlarged-kahler-moduli]{第 6.2.3 小节}将从镜像流形 $V^\circ$ 的 Kähler 模空间的角度重新讨论本例。特别是，将了解到为何应将 $\sigma_1,\sigma_2$ 合并为一个锥 $\sigma$ 的系统性原因。
\end{example}

\section{Kähler 模}\label{sec:kahler-moduli}\index{Kähler 模空间|hyperpage}
现在转向 Kähler 模。维数至少为 $2$ 的 Calabi--Yau 流形 $V$ 具有 Kähler 锥 $K(V)$，并且如定义~\ref{def:kahler-moduli} 所述，具有复化 Kähler 空间
\[
K_{\mathbb C}(V)=\{\omega\in H^2(V,\mathbb C):\operatorname{Im}(\omega)\in K(V)\}/\operatorname{im}H^2(V,\mathbb Z).
\]
\emph{复化 Kähler 模空间}则是商
\[
\mathcal{KM}(V)=K_{\mathbb C}(V)/\operatorname{Aut}(V).
\]
不会造成混淆时，将 $\mathcal{KM}(V)$ 简记为 $\mathcal{KM}$。本节的目标是理解 $\mathcal{KM}$ 的结构，特别是找到一个良好的紧化，其中包含与\hyperref[sec:complex-moduli]{第 6.1 节}研究的最大幺幂边界点相对应的点。与该节一样，在构造中保持一定灵活性会更方便。
\SourcePage{151}
\hyperref[sec:mirror-map]{第~\ref*{sec:mirror-map} 节}将讨论镜像映射。猜想它把 $\mathcal{KM}(V)$ 的\emph{大半径极限点}（将在\hyperref[subsec:large-radius-limits]{第 6.2.1 小节}定义）的一个邻域，映到 $V^\circ$ 的复结构模空间中一个最大幺幂边界点的邻域。这将把读者在本节与\hyperref[sec:complex-moduli]{第 6.1 节}之间发现的相似之处联系起来。镜像映射可用来将这些邻域与 SCFT 模空间的子集识别起来。一开始就必须强调，镜像映射只在局部定义，所以目前试图把本节构造的整个紧化，与 SCFT 模空间的一部分识别起来，会有风险。另一方面，将在例子中看到，本节构造的紧化确实为理解 SCFT 模空间的整体结构提供了有用的几何框架，不过还须由镜像映射加以“修正”。如何给出 SCFT 模空间的精确数学描述，仍是一个有待解开的精彩谜题。

\subsection{大半径极限点}\label{subsec:large-radius-limits}
研究 Kähler 模的第一步，是为复化 Kähler 空间 $K_{\mathbb C}(V)$ 找到一些良好的部分紧化。为此，固定一个满维单纯锥 $\sigma$，其内部位于 Kähler 锥 $K(V)$ 中，并令
\[
\mathcal D_\sigma=(H^2(V,\mathbb R)+i\operatorname{Int}(\sigma))/\operatorname{im}H^2(V,\mathbb Z)
\subset K_{\mathbb C}(V).
\]
注意，$\mathcal D_\sigma$ 是 $K_{\mathbb C}(V)$ 的开子集。进一步假定 $\sigma$ 由 $H^2(V,\mathbb Z)/\text{挠元}$ 的一组基 $T_1,\ldots,T_r$ 生成。$T_i$ 位于 Kähler 锥的闭包中。于是
$\operatorname{Int}(\sigma)=\{t_1T_1+\cdots+t_rT_r:t_1,\ldots,t_r>0\}$，而映射
\begin{equation}\label{eq:large-radius-exponential-coordinates}
t_1T_1+\cdots+t_rT_r\longmapsto(q_1,\ldots,q_r)
=(e^{2\pi i t_1},\ldots,e^{2\pi i t_r})
\end{equation}
诱导双全纯同构
\[
\mathcal D_\sigma\simeq(\Delta^*)^r.
\]
利用该同构及包含映射 $(\Delta^*)^r\subset\Delta^r$，得到部分紧化
\[
\mathcal D_\sigma\subset\mathcal D_\sigma^-\simeq\Delta^r.
\]
特别地，原点 $0\in\Delta^r$ 成为一个特殊边界点，记为 $0\in\mathcal D_\sigma^-$。

对 $\omega\in H^2(V,\mathbb R)+i\operatorname{Int}(\sigma)$，当 $\omega=\sum_jt_jT_j$ 的所有 $\operatorname{Im}(t_j)$ 都趋于 $+\infty$ 时，在映射~\eqref{eq:large-radius-exponential-coordinates} 下，$\omega$ 的像趋于 $0$。因此，称 $0\in\mathcal D_\sigma^-$ 为\emph{大半径极限点}。这是 Kähler 一侧与最大幺幂边界点对应的概念。注意，每个大半径极限点都具有上述典范坐标 $q_1,\ldots,q_r$。
\index{大半径极限点|hyperpage}

对于\hyperref[subsec:nef-gkz]{第 5.5.3 小节}定义的 nef 完全交，命题~\ref{prop:nef-i-gkz} 中已出现坐标 $q_1,\ldots,q_r$ 的一种版本。那里令 $q_j=e^{t_j}$，并对 $\beta\in H_2(V,\mathbb Z)$ 置
\[
q^\beta=\prod_jq_j^{\int_\beta T_j}
=\exp\left(\int_\beta(t_1T_1+\cdots+t_rT_r)\right)
\]
（见式~\eqref{eq:nef-gkz-q-coordinates} 后面的讨论）。若将 $t_j$ 换成 $2\pi i t_j$，就恰好得到上面的变量 $t_j,q_j$。所以，除多出的因子 $2\pi i$ 外，命题~\ref{prop:nef-i-gkz} 的形式函数 $\widetilde I$ 内在地定义于复化 Kähler 空间 $K_{\mathbb C}(V)$ 。由于 $\widetilde I$ 与镜像 $V^\circ$ 的 Picard--Fuchs 方程有关，这已经提示了\IndexLocation{indexloc:0830:151}镜像对称给出的复结构模空间与 Kähler 模空间之间的关系。

\SourcePage{152}
定义了大半径极限点以后，接着注意到，$\sigma$ 的不同选择给出许多这样的点。将 $\mathcal D_\sigma^-$ 上相应的坐标记为 $q_1^\sigma,\ldots,q_r^\sigma$，就可以研究改变 $\sigma$ 时这些坐标如何变化。例如，设有两个这样的锥，且 $\tau\subset\sigma$。这给出包含关系 $\mathcal D_\tau\subset\mathcal D_\sigma$。为联系相应的部分紧化，将 $\tau$ 的生成元用 $\sigma$ 的生成元表示为
\[
T_i^\tau=\sum_{j=1}^r a_{ij}T_j^\sigma,\qquad a_{ij}\geq0.
\]
则包含映射 $\mathcal D_\tau\subset\mathcal D_\sigma$ 延拓为映射 $\mathcal D_\tau^-\to\mathcal D_\sigma^-$，其表达式为
\[
q_j^\sigma=\prod_{i=1}^r(q_i^\tau)^{a_{ij}}.
\]
这令人想起式~\eqref{eq:mum-log-coordinate-change} 和式~\eqref{eq:mum-q-monomial-change}，而且相似之处比预想的更深。首先，\hyperref[subsec:mum-boundary-points]{第 6.1.1 小节}也有锥，即由单值性算子的对数生成的锥；式~\eqref{eq:mum-log-coordinate-change} 来自一个锥包含于另一个锥。此外，\hyperref[chap:quantum-cohomology]{第~\ref*{chap:quantum-cohomology} 章}研究量子上同调时，将看到这里锥 $\sigma$ 的生成元 $T_i$ 自然对应于 A 模型联络的单值性的对数。因此，两种情形极其相似，这正是镜像对称的预言。

大半径极限点与最大幺幂边界点之间，一个可能的差别在于：前者的局部坐标唯一，后者却因式~\eqref{eq:mum-coordinate-basis-change} 而不唯一。如\hyperref[subsec:mirror-map-definition]{第~\ref*{subsec:mirror-map-definition} 小节}将说明的，假定式~\eqref{eq:mum-canonical-coordinates} 中 $g_0,\ldots,g_r$ 具有整性，就能消除这一差别。这正好对应于上面假定 $T_1,\ldots,T_r$ 具有整性。

接着讨论特殊情形的 Kähler 模：$\operatorname{Aut}(V)=\{1\}$，且 $K(V)$ 是 $H^2(V,\mathbb R)$ 中一个有理多面锥的内部。此时 Kähler 模空间为 $\mathcal{KM}=K_{\mathbb C}(V)$，可以如下构造它的环面部分紧化。闭包 $\overline{K(V)}$ 关于格 $N=H^2(V,\mathbb Z)/\text{挠元}$ 为有理多面锥，所以得到\IndexLocation{indexloc:1354:152}仿射环面簇 $X_{\overline{K(V)}}$，简记为 $X_K$。同构
\begin{equation}\label{eq:kahler-torus-isomorphism}
H^2(V,\mathbb C)/N\simeq N\otimes\mathbb C^*=T_N=X_K\text{ 的环面}
\end{equation}
由 $K_{\mathbb C}(V)$ 的定义给出开集的自然包含映射
\begin{equation}\label{eq:kahler-toric-inclusion}
\mathcal{KM}=K_{\mathbb C}(V)\subset T_N\subset X_K.
\end{equation}
利用环面簇的边界 $B=X_K-T_N$，将部分紧化 $\overline{\mathcal{KM}}$ 定义为 $\mathcal{KM}$，加上 $B$ 中所有属于 $\mathcal{KM}$ 在 $X_K$ 中的像的闭包的点。

锥 $\overline{K(V)}$ 确定一点 $0\in X_K$，它是环面作用的不动点，并且可能有很严重的奇点。可以验证 $0\in\overline{\mathcal{KM}}$，因为在式~\eqref{eq:kahler-torus-isomorphism} 下，沿 Kähler 锥内部的射线使虚部趋于无穷时，像趋于 $0$。所以 $\overline{\mathcal{KM}}$ 很可能有奇点。

用标准环面方法消解该奇点。令 $\Sigma_+$ 为支集 $\overline{K(V)}$ 上的一个扇，使环面簇 $X_{\Sigma_+}$ 光滑。将 $X_K$ 换成 $X_{\Sigma_+}$ 后，式~\eqref{eq:kahler-toric-inclusion} 仍然适用；可以完全按前面的方式定义 $\overline{\mathcal{KM}}$，使它的边界由 $X_{\Sigma_+}-T_N$ 中属于 $\mathcal{KM}$ 的像的闭包的点组成。

\SourcePage{153}
$\Sigma_+$ 的极大锥给出 $X_{\Sigma_+}$ 的特殊点。我们断言，这些点全都是 $\overline{\mathcal{KM}}$ 的大半径极限点。原因在于，极大锥都是满足 $\operatorname{Int}(\sigma)\subset K(V)$ 的单纯锥 $\sigma$；这正是定义大半径极限点时所用的锥。将 $\mathcal{KM}$ 视为 $X_{\Sigma_+}$ 的子集。仿射环面簇 $X_\sigma$ 就是 $\mathbb C^r$；在其中取坐标多圆盘 $U_\sigma=\{|q_j|<1\text{，对所有 }j\}$，由式~\eqref{eq:kahler-toric-inclusion} 得到
\[
\mathcal D_\sigma=U_\sigma\cap\mathcal{KM}.
\]
很容易推出：$U_\sigma\cap\overline{\mathcal{KM}}$ \emph{恰好}是前面讨论的部分紧化 $\mathcal D_\sigma^-$。

因此，得到一个良好的部分紧化，其中有有限多个大半径极限点。此外，若 $\omega\in H^2(V,\mathbb R)+iK(V)$ 的虚部在某个极大锥 $\sigma\in\Sigma_+$ 的全部生成元方向上都趋于 $+\infty$，那么 $\omega$ 在 $\mathcal{KM}$ 中的像就趋于 $\overline{\mathcal{KM}}$ 中相应的大半径极限点。

这幅图景很好，但仍十分特殊。对于一般的 Calabi--Yau 流形，自同构群未必平凡，Kähler 锥也未必为多面锥。$K(V)$ 的结构可能相当复杂；\hyperref[subsec:threefold-kahler-moduli]{第 6.2.2 小节}详细研究三维簇时将看到这一点。一般而言，所能期望的最好情形是\emph{锥猜想} \SourceCite{Morrison3}。为陈述它，令 $K_+$ 为 $\overline{K(V)}\cap H^2(V,\mathbb Q)$ 的凸包。由于 $K(V)$ 为开集，$K(V)\subset K_+$，所以 $K_+$ 由 $K(V)$ 及其边界中由有理点生成的部分组成。
\begin{conjecture}\label{conj:kahler-cone}
设 $V$ 为维数至少为 $2$ 的 Calabi--Yau 流形，Kähler 锥为 $K(V)$。则存在有理多面锥 $\Pi\subset K_+$，使得
\[
\operatorname{Aut}(V)\cdot\Pi=K_+.
\]
\end{conjecture}
\index{锥猜想|hyperpage}
\SourceCite{GM} 已在一个非平凡例子中验证了锥猜想。此外，关于 $\operatorname{Aut}(V)$ 在 Kähler 锥上的作用，一些有趣的有限性结论是猜想~\ref{conj:kahler-cone} 的推论。例如，它表明模去 $\operatorname{Aut}(V)$ 后，$\overline{K(V)}$ 只有有限多条边。\SourceCite{Borcea1} 在一个特定 Calabi--Yau 三维簇 $V$ 上验证了这一点；该三维簇具有非多面的 Kähler 锥及无限的 $\operatorname{Aut}(V)$。锥猜想还表明，某些类的 Calabi--Yau 三维簇仅具有有限多个代数纤维空间结构；这一结论已知成立 \SourceCite{Oguiso,OgP}。\SourceCite{Szendroi} 证明了锥猜想的另一个推论。\SourceCite{Kawamata2} 提出了锥猜想的相对版本，并在其中证明了低维 Calabi--Yau 纤维空间的情形。

对我们而言，锥猜想的关键性质是：对于\emph{任何}满足它的 Calabi--Yau 流形，都能构造 $\mathcal{KM}$ 的一个部分紧化。因此，假定 $V$ 满足猜想~\ref{conj:kahler-cone}。Looijenga 关于管状域（如 $H^2(V,\mathbb R)+iK(V)$）的商的部分紧化的工作 \SourceCite{Looijenga} 表明，可以找到一个 $\operatorname{Aut}(V)$ 不变的扇 $\Sigma_+$（可能无限），其支集为 $K_+$。还可以假定自同构群在 $\Sigma_+$ 上自由作用，且 $\Sigma_+$ 的极大锥具有上述 $\sigma$ 的形式。

根据这些数据，Looijenga 从扇 $\Sigma_+$ 确定的各层构造空间 $\widehat{\mathcal D}(\Sigma_+)$。困难之处是赋予 $\widehat{\mathcal D}(\Sigma_+)$ 正确的拓扑，使商 $\widehat{\mathcal D}(\Sigma_+)/\operatorname{Aut}(V)$ 成为正规解析空间。这个商称为\emph{半环面部分紧化}，就是所需的部分紧化 $\overline{\mathcal{KM}}$。\SourceCite{Morrison3} 描述了这一构造的具体方式。

\SourcePage{154}
当 $\operatorname{Aut}(V)$ 正常不连续地作用，且 $K(V)$ 是齐次自伴锥 \SourceCite{AMRT} 时，完全可以在环面几何的范围内完成这一构造。为便于说明，在这一特殊情形中考察它与前面构造的关系。其思想是：无限扇 $\Sigma_+$ 确定一个局部有限型概形 $X_{\Sigma_+}$，将式~\eqref{eq:kahler-toric-inclusion} 中的 $X_K$ 换为 $X_{\Sigma_+}$，得到开集的包含映射
\[
K_{\mathbb C}(V)\subset X_{\Sigma_+}.
\]
由于扇在 $\operatorname{Aut}(V)$ 下不变，这一包含映射关于 $\operatorname{Aut}(V)$ 等变，所以得到包含映射
\[
\mathcal{KM}\subset X_{\Sigma_+}/\operatorname{Aut}(V).
\]
\SourceCite{AMRT} 给出了 $X_{\Sigma_+}/\operatorname{Aut}(V)$ 的一幅很好的图示，\SourceCite{Morrison3} 也转载了它。利用 $X_{\Sigma_+}/\operatorname{Aut}(V)$ 的环面边界，完全按前面的方式得到部分紧化 $\overline{\mathcal{KM}}$。此外，因为在这一情形中，有效自同构群在扇上自由作用，每个极大锥都给出一个大半径极限点，其邻域同构于 $\mathcal D_\sigma^-$。（不过，由于 $\operatorname{Aut}(V)$ 的作用，$\Sigma_+$ 中无限多个不同的锥可能给出同一个大半径极限点。）如 \SourceCite{Morrison3} 所示，这幅图景在 Looijenga 考虑的更一般情形中仍然成立。

所有这些讨论的结果是：得到一个良好的部分紧化 $\overline{\mathcal{KM}}$，其中包含有限多个大半径极限点。不过，这样的紧化有许多种，因为给定一个紧化后，总能进一步细分扇，构造其他紧化。因此，与最大幺幂边界点的情形一样，依照 \SourceCite{Morrison6}，在大半径极限点之间引入如下等价关系。
\begin{definition}\label{def:equivalent-large-radius-points}\index{大半径极限点!等价关系|hyperpage}
设 $\overline{\mathcal{KM}}$ 为 Kähler 模空间 $\mathcal{KM}$ 的一个光滑正规交叉部分紧化。若存在包含 $p,q$ 的连通集 $\Xi\subset\overline{\mathcal{KM}}-\mathcal{KM}$，并且在 $\Xi$ 的一个邻域上存在由 $\overline{\mathcal{KM}}$ 的余切丛中、沿边界有对数极点的形式组成的局部系统 $\mathcal L$，使得在 $p,q$ 各自附近，该局部系统由 $1$ 形式 $d\log q_k$ 张成，其中 $q_k$ 如式~\eqref{eq:large-radius-exponential-coordinates} 所述，则称大半径极限点 $p,q\in\overline{\mathcal{KM}}$ \emph{等价}。
\end{definition}
很容易看出，进一步细分扇 $\Sigma_+$ 所产生的多个大半径极限点，在这一意义下彼此等价。

如\hyperref[sec:complex-moduli]{第 6.1 节}所解释的，镜像对称告诉我们：$V$ 的 Kähler 模应通过镜像映射，同构于镜像 $V^\circ$ 的复结构模的一部分，且使大半径极限点的每个等价类映到最大幺幂边界点的一个等价类。事实上，镜像映射应保持两个等价关系定义中提到的局部系统。$\overline{\mathcal{KM}}$ 的构造有很大的非唯一性，镜像复结构模的紧化也不唯一。因此，镜像猜想的一部分内容，就是能找到相容的紧化。

实际处理中，在 $X_{\Sigma_+}$ 光滑的假设下，选取锥尽可能大的 $\Sigma_+$。这样可以缩小等价类，与\hyperref[sec:complex-moduli]{第 6.1 节}处理复结构模的做法相似。不过，若允许 $\mathcal{KM}$ 的紧化有奇点，则存在一个典范的极小紧化；它由包含 $K_+$ 及其各个面的扇 $\Sigma_{\mathrm{SBB}}$ 确定。由于 $K_+$ 可能有无限多个面，环面方法不适用，但 Looijenga 的构造
\SourcePage{155}
仍然有效，所得极小紧化记为 $\overline{\mathcal{KM}}_{\mathrm{SBB}}$。这一构造推广了有界对称域的商的 Satake--Baily--Borel 紧化构造，这也解释了该记号。当 $\operatorname{Aut}(V)=1$，且 $\overline{K(V)}$ 为有理多面锥时，$\overline{\mathcal{KM}}_{\mathrm{SBB}}$ 就是前面由 $\overline{K(V)}$ 的仿射环面簇 $X_K$ 得到的紧化。

Kähler 模存在 Satake--Baily--Borel 型极小紧化，通过镜像对称意味着：镜像的复结构模也应有类似的极小紧化。\SourceCite{Morrison3} 提出了这一方向的精确猜想。

最后讨论一个读者可能已经想到的问题：为什么对复结构模作紧化，却只对 Kähler 模作\emph{部分}紧化？原因是，$\mathcal{KM}$ 实际上只是整个 Kähler 模空间的一小部分；要找到缺失的部分，就需要“越出 Kähler 锥” \SourceCite{Morrison5}。一般 Calabi--Yau 流形的相应理论尚未建立，但在三维簇及环面超曲面两个特殊情形中，已经进行了广泛研究。因此，接着讨论这两类十分有趣的簇的 Kähler 模。

\subsection{Calabi--Yau 三维簇的 Kähler 模}\label{subsec:threefold-kahler-moduli}\index{Kähler 锥!三维簇|hyperpage}
令 $V$ 为光滑 Calabi--Yau 三维簇。这里有两个目标：首先，描述关于 Kähler 锥 $K(V)$ 及其闭包 $\overline{K(V)}$ 的结构的已知结果；其次，解释如何扩大 Kähler 模空间。

描述 Kähler 锥时，一个重要要素是\emph{三次锥} $W^*\subset H^2(V,\mathbb R)$，它定义为满足 $D^3=0$ 的所有类 $D$ 所组成的集合。基本思想是：在 $W^*$ 以外，$\overline{K(V)}$ 很容易理解。将使用 Wilson 的工作 \SourceCite{Wilson2,Wilson1}，它将 $\overline{K(V)}$ 的结构与本原收缩联系起来。
\begin{definition}\label{def:primitive-contraction}
\emph{收缩}是双有理态射 $V\to\overline V$，其中 $\overline V$ 为射影正规三维簇，且 $\rho(\overline V)<\rho(V)$（通常以 $\rho$ 表示 Picard 数）。若收缩不能进一步分解为正规簇之间的非平凡双有理态射，则称它为\emph{本原收缩}。
\end{definition}
\index{本原收缩|hyperpage}
若 $\phi$ 为本原收缩，则可以证明：除倍数外，$V$ 上被 $\phi$ 收缩的 $1$ 圈的数值类 $C$ 唯一。这定义超平面 $\mathcal H=\{D\in H^2(V,\mathbb R):D\cdot C=0\}$；很容易看出，该超平面支撑 $\overline{K(V)}$ 的一个余维一的面。称这个面\emph{对应于}该本原收缩。这些面决定 $\overline{K(V)}$ 在三次锥 $W^*$ 以外的结构。更确切地说，有如下结果。
\begin{proposition}[\SourceCite{Wilson2}]\label{prop:kahler-cone-contractions}\index{三次锥|hyperpage}
设 $V$ 为光滑 Calabi--Yau 三维簇。
\begin{enumerate}[label=(\roman*)]
\item 在 $W^*$ 以外，锥 $\overline{K(V)}$ 局部为有理多面锥。
\item 在 $W^*$ 以外，$\overline{K(V)}$ 的余维一的面对应于本原双有理收缩。
\end{enumerate}
\end{proposition}
这仍允许 Kähler 锥具有相当复杂的结构。例如，$\overline{K(V)}$ 可能在 $W^*$ 以外有无限多个面，并向 $W^*$ 聚集。还应指出，$W^*$ 可能包含 $\overline{K(V)}$ 的余维一的面。这样的面由将 $V$ 映到一个较低维簇的线性系 $|D|$ 确定，而该线性系的一般成员 $D$ 的类位于这个面上。
\SourcePage{156}
接着描述如何扩大 $V$ 的 Kähler 模空间。思想是将 Kähler 锥 $K(V)$ 跨过 $\overline{K(V)}$ 的某些面延伸出去。将使用对应于本原收缩的面；回顾，这样的收缩 $\phi$ 分为三种类型：
\begin{itemize}
\item I 型：存在曲线 $C\subset V$，使 $\phi$ 在 $V-C$ 上为同构，并将 $C$ 映到一点。
\item II 型：存在曲面 $S\subset V$，使 $\phi$ 在 $V-S$ 上为同构，并将 $S$ 映到一点。
\item III 型：存在曲面 $S\subset V$，使 $\phi$ 在 $V-S$ 上为同构，并将 $S$ 映到一条曲线。
\end{itemize}
现在假定 $\overline{K(V)}$ 的一个面对应于 I 型本原收缩 $\phi:V\to\overline V$。有一条不必不可约的曲线 $C\subset V$ 被 $\phi$ 收缩。三维簇几何的一个基本定理表明，可以对 $C$ 作\emph{翻越} \SourceCite{Reid2,Kollar}。特别地，这意味着存在不同于 $V$ 的 Calabi--Yau 三维簇 $V'$、曲线 $C'\subset V'$ 及映射 $\phi':V'\to\overline V$，它收缩 $C'$，并在 $V'-C'$ 上为同构。所得双有理映射 $V\dashrightarrow V'$ 在余维一处为同构，所以有自然识别 $H^2(V)\simeq H^2(V')$。在这一识别下，$\overline{K(V)}$ 与 $\overline{K(V')}$ 沿各自的一个公共面相交；它们分别由超平面 $\{D:D\cdot C=0\}$ 及 $\{D':D'\cdot C'=0\}$ 支撑。
\index{翻越|hyperpage}

Kawamata 的结果 \SourceCite{Kawamata1} 表明：对于所有可从 $V$ 经一系列翻越得到的 Calabi--Yau 三维簇 $V'$（包括 $V$ 自身），锥 $\overline{K(V')}$ 构成其并集的凸包的一个多面分解的各个腔；该凸包称为\emph{可动锥} $\operatorname{Mov}(V)$。因为 $V$ 的所有双有理模型都由一系列翻越构造，$V$ 的双有理自同构群 $\operatorname{Bir}(V)$ 自然作用于 $\operatorname{Mov}(V)$。锥猜想的双有理版本是\emph{\IndexLocation{indexloc:0158:156}双有理锥猜想} \SourceCite{Morrison5}；它断言存在有理多面锥 $\Pi$，使
$\operatorname{Bir}(V)\cdot\Pi=\operatorname{Mov}(V)_+$，其中 $\operatorname{Mov}(V)_+$ 是 $\operatorname{Mov}(V)\cap H^2(V,\mathbb Q)$ 的凸包。用与前面类似的构造，可以适当细分 $\operatorname{Mov}(V)$，同时构造所有 $V'$ 的 Kähler 模空间的一个部分紧化。\SourceCite{Morrison5} 给出了更详细的说明；\hyperref[subsec:enlarged-kahler-moduli]{第 6.2.3 小节}将结合环面超曲面重新讨论这一思想。
\index{可动锥|hyperpage}

还可以进一步扩大 Kähler 模空间。设 $\phi$ 是 III 型收缩，其例外曲面为一个以亏格 $g\geq2$ 的曲线 $C$ 为底的直纹曲面 $S$，且 $\phi$ 沿直纹将 $S$ 收缩到 $C$。若对 $V$ 的复结构作一般形变，$S\to C$ 的一般纤维就不再保持全纯，但有 $2g-2$ 个纤维可以全纯形变。事实上，一个纤维的一阶形变障碍，就是 Kodaira--Spencer 类 $\rho\in H^1(\Theta_V)$ 在自然映射
$r:H^1(\Theta_V)\to H^0(K_C)$ 下的像；若 $\rho$ 一般，则 $r(\rho)$ 在 $2g-2$ 个不同的点为零。这些点上的纤维可以一阶形变 \SourceCite{Wilson2}。还可计算出没有更高阶障碍。于是可同时翻越这 $2g-2$ 条形变后的曲线。翻越所得 Calabi--Yau 三维簇属于同一个复结构模空间，但复结构不同。Kähler 结构也随之改变。所得 $H^2$ 上的同构可以沿形变追溯回原来的复结构，从而成为 $H^2(V)$ 的自同构。置 $E=[S]$，该自同构为
\[
D\longmapsto D+(D\cdot f)E,\qquad f=\text{纤维的类}.
\]

\SourcePage{157}
由于 $E\cdot f=-2$，这是关于超平面 $\{D:D\cdot f=0\}$ 的反射。用上述所有可能的 III 型收缩给出的反射，对 $\operatorname{Mov}(V)$ 作反射，就得到\emph{反射可动锥} \SourceCite{Morrison5}。若同时扩大双有理自同构群，使它包含上述反射，就得到一个更大的 Kähler 模空间，使用的正是反射可动锥及刚才描述的扩大自同构群。此外，锥猜想的适当版本还能使我们对这一扩大模空间作部分紧化。

反射可动锥给出的 Kähler 模空间，是跨越 $\overline{K(V)}$ 中对应于 I 型和 III 型收缩的墙构造的，不过只使用了将直纹曲面收缩到亏格 $g\geq2$ 的曲线的 III 型收缩。目前，跨过对应于 II 型及其余 III 型收缩的墙后，应如何扩大 Kähler 模空间，还没有方案。不过，在\hyperref[subsec:enlarged-kahler-moduli]{第 6.2.3 小节}将看到，对于 Calabi--Yau 环面超曲面，可以进一步扩大 Kähler 模空间。

最后讨论的主题是：改变 $V$ 的复结构时，Kähler 锥如何变化。\hyperref[sec:physics]{第 1.1 节}的简单图景似乎表明，可以保持 Kähler 结构不变，同时改变复结构。对于三维簇，这在一般情形下成立。更确切地说，Wilson \SourceCite{Wilson2} 证明：Calabi--Yau 三维簇 $V$ 的复结构处于一般位置时，Kähler 锥在复结构的小形变下保持不变。这种一般位置的刚性很强。例如，\SourceCite{Wilson3} 证明：若两个 Calabi--Yau 三维簇 $V_1,V_2$ 在各自的复结构模空间中处于一般位置，且互为辛形变，那么它们的 Kähler 锥相同；\SourceCite{Wilson4} 则表明，在三次锥 $W^*$ 以外，与 $\overline{K(V_1)}$ 的某个余维一的面对应的本原收缩的类型（I、II 或 III 型），与 $\overline{K(V_2)}$ 的相应面对应的本原收缩的类型相同。

然而，实际情形更为复杂，\hyperref[sec:physics]{第 1.1 节}图下的脚注已暗示了这一点。特别是，有些 $V$ 的形变会使 Kähler 锥变化。设 $\phi$ 为 III 型收缩，将直纹曲面 $S$ 映到亏格为 $g$ 的曲线 $C$。当 $g\geq2$ 时，前面可以用反射扩展可动锥。不过，底曲线满足 $g=1$ 时，上面的论证表明：在复结构的一般形变中，没有任何纤维随之形变。这意味着，对含有曲面 $S$（其类为 $E$）的特殊复结构，有效曲线锥严格包含一般复结构的有效曲线锥，因为一般复结构中不存在以类 $f$ 表示的有效曲线。对偶地，这一特殊复结构的 Kähler 锥变小。因此，Kähler 模空间可能依赖于复结构。

\subsection{环面超曲面的 Kähler 模}\label{subsec:enlarged-kahler-moduli}
设 $\Delta$ 为自反多胞形，$\Sigma$ 为 $\Delta$ 的法扇的一个极大射影细分。若 $V\subset X_\Sigma$ 为一般 Calabi--Yau 超曲面，回顾\hyperref[chap:mirror-constructions]{第~\ref*{chap:mirror-constructions} 章}对限制映射 $H^2(X_\Sigma)\to H^2(V)$ 的描述。它的核由除子类 $[D_\rho]$ 生成；这些除子对应于由 $\Delta^\circ$ 某个刻面的内部本原格点张成的一维锥 $\rho\in\Sigma(1)$。限制映射的像按定义为 $H_{\mathrm{toric}}^2(V)$。目标是为 $V$ 的“环面
\SourcePage{158}
部分”构造 Kähler 模空间。从某种意义上说，这一构造与\hyperref[subsec:toric-complex-moduli]{第 6.1.2 小节}中 $\overline{\mathcal M}_{\mathrm{simp}}$ 的构造互为“镜像”。

与\hyperref[subsec:toric-complex-moduli]{第 6.1.2 小节}一样，稍微改变问题，排除 $\Delta^\circ$ 的刻面内部格点。相应地，稍微修改极大射影细分的定义，以避开这些点。
\begin{definition}\label{def:simplified-projective-subdivision}\index{简化射影细分|hyperpage}
给定自反多胞形 $\Delta\subset M_{\mathbb R}$，若 $N_{\mathbb R}$ 中的扇 $\Sigma$ 具有以下性质，则称为\emph{简化射影细分}：
\begin{itemize}
\item $\Sigma$ 细分 $\Delta$ 的法扇。
\item $\Sigma(1)=(\Delta^\circ\cap N)_0-\{0\}$。
\item $X_\Sigma$ 射影且单纯。
\end{itemize}
\end{definition}
现在可以仿照\hyperref[sec:batyrev-mirror]{第 4.1 节}的 Batyrev 构造，将极大射影细分 $\Sigma$ 换成一个简化射影细分，仍记为 $\Sigma$。下面证明这样得到的 Calabi--Yau 簇，与 Batyrev 构造给出的为同一类。首先，给定极大射影细分 $\Sigma$，去掉由 $\Delta^\circ$ 的刻面内部格点生成的所有一维锥，就得到 $X_\Sigma$ 的一个爆缩的扇。所得扇未必单纯，所以可能需要在不加入新一维锥的情况下进一步细分，以得到简化射影细分 $\Sigma'$。得到交换图
\begin{equation}\label{eq:simplified-subdivision-diagram}
\begin{array}{ccc}
X_\Sigma&\overset{\phi}{\dashrightarrow}&X_{\Sigma'}\\
\scriptstyle\pi\searrow&&\swarrow\scriptstyle\pi'\\
&\mathbb P_\Delta&
\end{array}
\end{equation}
按构造，除 $\mathbb P_\Delta$ 中与 $\Delta^\circ$ 的那些含内部格点的刻面对应的点的上方以外，$\phi$ 为同构。一般 Calabi--Yau 超曲面 $\overline V\in|-K_{\mathbb P_\Delta}|$ 与这一有限集不交，因此它在 $\pi,\pi'$ 下的严格变换——即 $X_\Sigma,X_{\Sigma'}$ 中的一般 Calabi--Yau 超曲面——同构，正如所断言的。反过来更简单：任何简化射影细分都可细分为极大射影细分，仍得到同构的 Calabi--Yau 超曲面。

现在固定 $\Delta$ 的法扇的一个简化射影细分 $\Sigma$，考虑 $V\subset X_\Sigma$。由于排除了 $\Delta^\circ$ 的刻面内部格点，自然映射 $H^2(X_\Sigma)\to H^2(V)$ 为单射。其像已知为 $H_{\mathrm{toric}}^2(V)$，所以
\[
H^2(X_\Sigma)\simeq H_{\mathrm{toric}}^2(V).
\]
\hyperref[subsec:toric-kahler-cone]{第 3.3.3 小节}表明，$X_\Sigma$ 的 Kähler 锥的闭包为 $\overline{K(X_\Sigma)}=\operatorname{cpl}(\Sigma)$。因此，可以将 $\operatorname{cpl}(\Sigma)$ 视为位于 $H_{\mathrm{toric}}^2(V)$ 中。将 $V$ 的\emph{Kähler 锥的环面部分}定义为
\[
K(V)_{\mathrm{toric}}=K(V)\cap H_{\mathrm{toric}}^2(V),
\]
则有显然的包含关系 $\operatorname{cpl}(\Sigma)\subset\overline{K(V)}_{\mathrm{toric}}$，因为 $X_\Sigma$ 的丰沛类限制到 $V$ 后仍为丰沛类。

暂且假定总有 $\operatorname{cpl}(\Sigma)=\overline{K(V)}_{\mathrm{toric}}$。此时，很容易构造一个扩大的环面 Kähler 模空间，并用\hyperref[sec:gkz-decomposition]{第 3.4 节}的次级扇将它紧化。置 $\Xi=(\Delta^\circ\cap N)_0-\{0\}$，次级扇自然位于向量空间
$A(\Xi)\simeq H^2(X_\Sigma,\mathbb R)\simeq H_{\mathrm{toric}}^2(V,\mathbb R)$ 中（沿用\hyperref[sec:gkz-decomposition]{第~\ref*{sec:gkz-decomposition} 节}的记号）。次级扇及整上同调给出的格确定一个完备环面簇，暂且将它
\SourcePage{159}
记为 $\overline{\mathcal{KM}}_{\mathrm{toric}}$。对于每个简化射影细分 $\Sigma$，$\operatorname{cpl}(\Sigma)$ 是次级扇的一个极大锥；若 $\operatorname{cpl}(\Sigma)=\overline{K(V)}_{\mathrm{toric}}$，则 $\overline{\mathcal{KM}}_{\mathrm{toric}}$ 中相应的点就是\hyperref[subsec:large-radius-limits]{第 6.2.1 小节}意义下的大半径极限点（可能需要细分 $\operatorname{cpl}(\Sigma)$，以得到光滑空间）。

这幅图景的问题在于：包含关系 $\operatorname{cpl}(\Sigma)\subset\overline{K(V)}_{\mathrm{toric}}$ 中，\emph{不一定取等号}。环面 Kähler 锥的闭包 $\overline{K(V)}_{\mathrm{toric}}$ 可能严格大于 $\operatorname{cpl}(\Sigma)$。原因是，可能存在与 $\Sigma$ 不同的简化射影细分 $\Sigma'$，给出同构的 Calabi--Yau 超曲面。因此，$\overline{K(V)}_{\mathrm{toric}}$ 可能同时包含 $\operatorname{cpl}(\Sigma)$ 和 $\operatorname{cpl}(\Sigma')$。本小节后面将看到一个例子。

事实上，目前还不知道环面 Kähler 锥 $\overline{K(V)}_{\mathrm{toric}}$ 的精确结构。猜想它为有理多面锥；在所有已计算出 $\overline{K(V)}_{\mathrm{toric}}$ 的情形中，确实如此。下面的猜想~\ref{conj:toric-kahler-cone} 将仔细描述我们认为一般情形下 $\overline{K(V)}_{\mathrm{toric}}$ 应有的形态。猜想的实质是：$\overline{K(V)}_{\mathrm{toric}}$ 是某些锥 $\operatorname{cpl}(\Sigma')$ 的并集。若猜想正确，就意味着由次级扇构造的空间 $\overline{\mathcal{KM}}_{\mathrm{toric}}$，是扩大环面 Kähler 模空间的“真正”紧化的一个爆破。

本小节余下的计划是：先考察 $\overline{K(V)}_{\mathrm{toric}}$ 的结构，再构造一个扩大 Kähler 模空间，并将它紧化。

开始研究 $\overline{K(V)}_{\mathrm{toric}}$ 时，需要考察：$\Sigma$ 遍历所有简化射影细分时，各锥 $\operatorname{cpl}(\Sigma)$ 如何拼接。例如，如何从一个细分转到“邻近”的另一个细分？为此使用\emph{线性回路}的语言，现在给出定义。
\begin{definition}\label{def:linear-circuit}
给定 $\Xi=(\Delta^\circ\cap N)_0-\{0\}$，\emph{线性回路}是线性相关的子集 $\Xi'\subset\Xi$，使 $\Xi'$ 的任何真子集都线性无关。
\end{definition}
\index{线性回路|hyperpage}
更多细节可参阅 \SourceCite{GKZ2} 及其参考文献。要注意，\SourceCite{GKZ2} 的情形稍有不同：那里研究\emph{仿射回路}，使用仿射相关与三角剖分，而不是线性相关与扇。

下面用 $v_i$ 表示 $\Xi$ 的元素。给定线性回路 $\Xi'\subset\Xi$，固定其元素 $v_i\in\Xi'$ 之间的非平凡线性关系 $\sum_i b_i v_i=0$，并分解为 $\Xi'=\Xi'_+\cup\Xi'_-$；其中 $\Xi'_+$ 为在线性关系中具有正系数 $b_i$ 的元素 $v_i\in\Xi'$ 构成的集合，$\Xi'_-$ 类似定义。

对任何子集 $S\subset\Xi$，以 $C(S)$ 表示 $S$ 生成的凸锥。给定满足 $|\Xi'_+|\geq1$ 的回路 $\Xi'\subset\Xi$，将 $\Sigma_+(\Xi')$ 定义为极大维锥由 $C(\Xi'-\{v_i\})$（$v_i\in\Xi'_+$）组成的扇，其支集为 $C(\Xi')$。类似地，若 $|\Xi'_-|\geq1$，就得到扇 $\Sigma_-(\Xi')$。举例说明，设 $\Xi'=\{v_1,v_2,v_3,v_4\}$，线性关系为 $v_1+v_2=v_3+v_4$。此时，$C(\Xi')$ 是有四个生成元的锥。将它细分为一对单纯锥有两种办法：$\Sigma_+(\Xi')$ 给出一种，$\Sigma_-(\Xi')$ 给出另一种。

接着描述线性回路与扇 $\Sigma$ 如何相互作用。

\SourcePage{160}
\begin{definition}\label{def:circuit-supported-fan}\index{线性回路!扇的支撑|hyperpage}
设 $\Sigma$ 为 $N_{\mathbb R}$ 中的射影单纯扇，$\Sigma(1)\subset\Xi$；令 $\Xi'\subset\Xi$ 为线性回路。若满足下列条件，则称 $\Sigma$ \emph{支撑于} $\Xi'$（或称 $\Xi'$ \emph{由} $\Sigma$ \emph{支撑}）：
\begin{enumerate}[label=(\roman*)]
\item $\Sigma_+(\Xi')$ 是 $\Sigma$ 的子扇。
\item 设 $\sigma$ 为 $\Sigma_+(\Xi')$ 的一个极大维锥。若存在子集 $\Pi\subset\Xi-\Xi'$，使 $\sigma\cup\Pi$ 生成 $\Sigma$ 的一个极大维锥，则对 $\Sigma_+(\Xi')$ 的其他所有极大维锥 $\sigma'$，$\sigma'\cup\Pi$ 也生成 $\Sigma$ 的一个锥。
\end{enumerate}
\end{definition}
因此，支撑于 $\Xi'$ 意味着：$\Sigma$ 中与 $C(\Xi')$ 相接的锥，由 $\Sigma_+(\Xi')$ 以一种特别强的方式确定。本小节后面重新讨论\hyperref[ex:degree-four-moduli]{例题~\ref*{ex:degree-four-moduli}} 时，将给出一个支撑于回路的射影细分的例子。

定义~\ref{def:circuit-supported-fan} 使我们能够用线性回路修改扇 $\Sigma$。更确切地说，假定 $\Sigma$ 支撑于 $\Xi'$，且 $|\Xi'_-|\geq1$。将 $\Sigma$ 中由 $\sigma\cup\Pi$ 张成的单纯锥——其中 $\sigma$ 是 $\Sigma_+(\Xi')$ 的锥，$\Pi\subset\Xi-\Xi'$——换成由 $\sigma'\cup\Pi$ 张成的单纯锥，其中 $\sigma'$ 是 $\Sigma_-(\Xi')$ 的锥，就得到新的单纯扇 $\Sigma'=\operatorname{flip}_{\Xi'}(\Sigma)$。例如，设 $\Sigma$ 支撑于线性回路 $\Xi'=\{v_1,v_2,v_3,v_4\}$，且 $v_1+v_2=v_3+v_4$（即上述例子）。此时，$\Sigma$ 使用将 $C(\Xi')$ 分为两个单纯锥的一种方式，翻转所得的 $\Sigma'$ 则使用另一种方式。
\index{翻转|hyperpage}

从满足 $\Sigma(1)\subset\Xi$ 的单纯射影扇 $\Sigma$ 出发，翻转 $\Sigma'$ 显然仍单纯，并满足 $\Sigma'(1)\subset\Xi$。在所考虑的相邻射影腔之间的翻转中，利用 \SourceCite{OdP} 的技巧还能证明 $\Sigma'$ 也射影。因此，$\operatorname{cpl}(\Sigma')$ 是 GKZ 分解中的一个锥。不过，即使 $\Sigma$ 细分 $\Delta$ 的法扇（如简化射影细分的定义所要求的），$\Sigma'$ 也未必如此。但若 $\Xi'$ 位于 $\Delta^\circ$ 的一个真面中，则 $\Sigma'$ 也细分 $\Delta$ 的法扇。

接着考察翻转 $\Sigma'=\operatorname{flip}_{\Xi'}(\Sigma)$ 的几何意义。设 $\Sigma$ 为单纯射影扇，$\Sigma(1)\subset\Xi$，且 $|\Xi'_+|,|\Xi'_-|\geq2$。由定义立即可知，$X_{\Sigma'}$ 是从 $X_\Sigma$ 爆缩子簇 $\bigcap_{v_i\in\Xi'_-}D_i$，再爆破所得奇点，以 $X_{\Sigma'}$ 的子簇 $\bigcap_{v_i\in\Xi'_+}D_i$ 代替它而得到的。将 $X_\Sigma$ 到 $X_{\Sigma'}$ 的这一双有理变换称为\emph{广义翻越}，记为 $\operatorname{flop}_{\Xi'}$。例如，余维二的经典翻越可看作与回路 $\{v_1,v_2,v_3,v_4\}$、关系 $v_1+v_3=v_2+v_4$ 对应的广义翻越。

现在令 $\Xi'$ 为由 $\Sigma$ 支撑的线性回路，满足 $|\Xi'_+|,|\Xi'_-|\geq2$。再假定 $\operatorname{flop}_{\Xi'}$ 的基点集与 $V$ 不交，即对一般 Calabi--Yau 超曲面 $V\subset X_\Sigma$，有
\[
\bigcap_{v_i\in\Xi'_-}D_i\cap V=\varnothing.
\]
立即可知，$\operatorname{flop}_{\Xi'}$ 诱导 Calabi--Yau 超曲面之间的同构 $V\simeq V'$。将这种翻越称为\emph{\IndexLocation{indexloc:0335:160}平凡翻越}，相应的翻转 $\Sigma'=\operatorname{flip}_{\Xi'}(\Sigma)$ 称为 $\Sigma$ 的\emph{平凡翻转}。

当 $\Sigma'$ 为 $\Sigma$ 的平凡翻转时，得到映射
\[
H^2(X_{\Sigma'})\longrightarrow H^2(V')\simeq H^2(V),
\]
它把 Kähler 类映到 Kähler 类。因此，$\operatorname{cpl}(\Sigma')\subset\overline{K(V)}_{\mathrm{toric}}$。由此引出如下定义。给定 Calabi--Yau 环面
\SourcePage{161}
超曲面 $V\subset X_\Sigma$，考虑锥
\begin{equation}\label{eq:trivial-flip-kahler-union}\index{平凡翻转|hyperpage}
\widetilde{K(V)}_{\mathrm{toric}}=\bigcup_{\Sigma'}\operatorname{cpl}(\Sigma'),
\end{equation}
其中并集遍及所有可由 $\Sigma$ 经一系列平凡翻转得到的扇 $\Sigma'$。由上述讨论立即可知
\[
\widetilde{K(V)}_{\mathrm{toric}}\subset\overline{K(V)}_{\mathrm{toric}}.
\]
\begin{conjecture}\label{conj:toric-kahler-cone}\index{环面 Kähler 锥|hyperpage}
$\widetilde{K(V)}_{\mathrm{toric}}=\overline{K(V)}_{\mathrm{toric}}$。
\end{conjecture}
据我们所知，在所有已知例子中，$\widetilde{K(V)}_{\mathrm{toric}}$ 都与 $\overline{K(V)}_{\mathrm{toric}}$ 相等。有趣的是，在这些例子中，$\widetilde{K(V)}_{\mathrm{toric}}$ 总是单纯锥，尽管 $\overline{K(X_\Sigma)}$ 未必如此。下面重新讨论\hyperref[ex:degree-four-moduli]{例题~\ref*{ex:degree-four-moduli}} 时，将说明猜想~\ref{conj:toric-kahler-cone}，以及一些可能有助于最终证明它的思想。然而，\SourceCite{Szendroi2} 表明，猜想~\ref{conj:toric-kahler-cone} 一般情形下不成立。

\emph{本小节余下部分都假定猜想~\ref{conj:toric-kahler-cone} 成立。}

来探究这对 Kähler 模的环面部分意味着什么。将尚未扩大的 Kähler 模空间的环面部分定义为商
\begin{equation}\label{eq:toric-kahler-moduli}\index{环面 Kähler 模空间|hyperpage}
\mathcal{KM}_{\mathrm{toric}}
=\frac{\{\omega\in H_{\mathrm{toric}}^2(V,\mathbb C):
\operatorname{Im}(\omega)\in K(V)_{\mathrm{toric}}\}}
{\operatorname{im}H_{\mathrm{toric}}^2(V,\mathbb Z)}.
\end{equation}
严格地说，还应对 $V$ 的保持环面部分的自同构取商。不过，这些自同构保持有理多面锥 $\overline{K(V)}_{\mathrm{toric}}$，每个自同构都置换其本原射线生成元。因此，它们通过一个有限自同构群作用，真正的环面 Kähler 模空间就是式~\eqref{eq:toric-kahler-moduli} 的有限商。这类似于\hyperref[subsec:toric-complex-moduli]{第 6.1.2 小节}中 $\mathcal M_{\mathrm{poly}}$ 是 $\mathcal M_{\mathrm{simp}}$ 的有限商的情形。正如使用 $\mathcal M_{\mathrm{simp}}$ 更方便，这里考虑按上述方式定义的 $\mathcal{KM}_{\mathrm{toric}}$ 也更简单。

有了这一环面 Kähler 模空间的定义，可以看到，式~\eqref{eq:trivial-flip-kahler-union} 中的锥 $\operatorname{cpl}(\Sigma')$ 给出一个细分 $\overline{K(V)}_{\mathrm{toric}}$ 的扇。与\hyperref[subsec:large-radius-limits]{第 6.2.1 小节}一样，所得环面簇给出 $\mathcal{KM}_{\mathrm{toric}}$ 的一个部分紧化；若进一步细分以得到光滑部分紧化，就会得到许多大半径极限点。其中包括来自 $\operatorname{cpl}(\Sigma)$ 的点，但也可能有更多点，因为 $\overline{K(V)}_{\mathrm{toric}}$ 中可能出现其他锥 $\operatorname{cpl}(\Sigma')$。所有这些点都应对应于镜像一侧的最大幺幂边界点。

若要寻找 $\mathcal{KM}_{\mathrm{toric}}$ 的\emph{极小}部分紧化，显然应使用 $\overline{K(V)}_{\mathrm{toric}}$ 给出的仿射环面簇。在镜像一侧，本章后面将看到：需要将 $\overline{\mathcal M}_{\mathrm{simp}}$ 略作爆缩，才能得到简化复结构模空间 $\mathcal M_{\mathrm{simp}}$ 的相应极小紧化；\hyperref[ex:degree-four-moduli]{例题~\ref*{ex:degree-four-moduli}} 已暗示了这一点。

下面讨论线性回路的一些其他性质，它们在计算例子时很有用。首先观察到，线性回路给出对偶空间 $A(\Xi)^*$ 中有趣的元素。由于 $A(\Xi)=\mathbb R^\Xi/M_{\mathbb R}$，$A(\Xi)^*$ 是线性关系的向量空间
\[
A(\Xi)^*=\left\{\ell=(\ell_i)\in\mathbb R^\Xi:\sum_i\ell_i v_i=0\right\}.
\]
若 $\Xi'$ 为线性回路，有关系 $\sum_i b_i v_i=0$，由此得到分解 $\Xi'=\Xi'_+\cup\Xi'_-$。该关系还给出 $\ell_{\Xi'}=(b_i)\in A(\Xi)^*$；
\SourcePage{162}
固定分解 $\Xi'=\Xi'_+\cup\Xi'_-$ 后，它确定到一个正标量倍。将元素 $D\in A(\Xi)$ 视为除子类 $D=[\sum_i a_iD_i]$，则
\[
\ell_{\Xi'}(D)=\sum_i a_i b_i.
\]
现在可用 $\Sigma$ 所支撑的线性回路刻画 $\operatorname{cpl}(\Sigma)$。若 $\Sigma$ 支撑于 $\Xi'$，且 $\Sigma'=\operatorname{flip}_{\Xi'}(\Sigma)$，很容易证明：对 $D\in\operatorname{cpl}(\Sigma)$，有 $\ell_{\Xi'}(D)\geq0$；对 $D'\in\operatorname{cpl}(\Sigma')$，有 $\ell_{\Xi'}(D')\leq0$。因此，超平面 $\ell_{\Xi'}=0$ 将两个锥的内部分开。此外，可以证明
\[
\operatorname{cpl}(\Sigma)=\{D\in A(\Xi):\ell_{\Xi'}(D)\geq0
\text{，对所有由 }\Sigma\text{ 支撑的 }\Xi'\}.
\]
事实上，$\operatorname{cpl}(\Sigma)$ 的刻面与由 $\Sigma$ 支撑的线性回路 $\Xi'$ 一一对应。

还可用线性回路缩小 GKZ 分解。这很有用，因为我们真正感兴趣的，并不是那些 $\Sigma(1)$ 为 $\Xi=(\Delta^\circ\cap N)_0-\{0\}$ 的真子集所对应的锥 $\operatorname{cpl}(\Sigma)$。思想如下。令 $\Sigma$ 为简化射影细分，设 $\Xi'\subset\Xi$ 为任意满足 $|\Xi'_-|=1$ 的线性回路，不必由 $\Sigma$ 支撑。我们断言，$\operatorname{cpl}(\Sigma)$ 上有 $\ell_{\Xi'}\geq0$。为此，重标 $\Xi'$，使线性关系为 $b_0v_0=\sum_{i>0}b_i v_i$，所有系数都为正。由于 $v_0\in\Xi=\Sigma(1)$，某个极大锥 $\sigma\in\Sigma$ 含有 $v_0$。取 $D=[\sum_i a_iD_i]\in\operatorname{cpl}(\Sigma)$。因为 $\sigma$ 单纯，可以找到 $m_\sigma\in M_{\mathbb R}$，使对所有 $v_i\in\sigma$ 都有 $\langle m_\sigma,v_i\rangle=-a_i$；还注意到，对任意 $v_i\in\Xi$，都有 $\langle m_\sigma,v_i\rangle\geq-a_i$，因为 $D\in\operatorname{cpl}(\Sigma)$（见\hyperref[subsec:toric-kahler-cone]{第 3.3.1 小节}）。于是
\[
-b_0a_0=\langle m_\sigma,b_0v_0\rangle
=\left\langle m_\sigma,\sum_{i>0}b_i v_i\right\rangle
\geq-\sum_{i>0}a_i b_i,
\]
立即得到 $\ell_{\Xi'}(D)\geq0$。为使用这一结论，回顾锥
\[
A^+(\Xi)=\left\{\left[\sum_i a_iD_i\right]\in A(\Xi):a_i\geq0\text{，对所有 }i\right\}
\]
是 GKZ 分解的支集。将 \emph{Calabi--Yau 锥}\footnote{B. Sturmfels 向我们指出了这个锥的重要性。}定义为
\[
C_{\mathrm{CY}}=\{D\in A^+(\Xi):\ell_{\Xi'}(D)\geq0
\text{，对所有满足 }|\Xi'_-|=1\text{ 的 }\Xi'\}.
\]
关键在于，若 $\Sigma$ 为简化射影细分，由前一段知 $\operatorname{cpl}(\Sigma)\subset C_{\mathrm{CY}}$。所以，只需研究 GKZ 分解在锥 $C_{\mathrm{CY}}$ 上的限制。

下面用一个前面见过的例子，说明迄今建立的概念。
\begingroup
\SourceExampleNumber{6.1.4}{2}
\renewcommand{\theHexam}{6.1.4.revisited.\arabic{exam}}
\begin{example}[重访]\label{ex:revisited-toric-kahler}
\hyperref[ex:degree-four-moduli]{例题~\ref*{ex:degree-four-moduli}} 研究了由自反多胞形 $\Delta$ 得到的 Calabi--Yau 三维簇 $V\subset X_\Sigma$ 的复结构模；$\Delta$ 的顶点 $v_1,\ldots,v_6$ 为
\begin{equation}\label{eq:revisited-polytope}
\begin{gathered}
(-1,-2,-3,-7),\ (1,0,0,0),\ (0,1,0,0),\\
(0,0,1,0),\ (0,0,0,1),\ (0,0,-1,-2).
\end{gathered}
\end{equation}
现在要研究镜像 $V^\circ\subset X_{\Sigma^\circ}$ 的 Kähler 模。因为处理的是 $V^\circ$，需要交换 $N,M$、$\Delta,\Delta^\circ$ 等。

首先观察到，$K(V^\circ)=K(V^\circ)_{\mathrm{toric}}$。这由\hyperref[sec:batyrev-mirror]{第 4.1 节}的公式，特别是式~\eqref{eq:toric-hodge-subspace} 和式~\eqref{eq:batyrev-h11} 推出，因为它们表明 $h^2(V^\circ)=h_{\mathrm{toric}}^2(V^\circ)=2$。

\SourcePage{163}
$\Delta^\circ$ 的法扇的简化射影细分 $\Sigma^\circ$ 位于 $M_{\mathbb R}$，满足 $\Sigma^\circ(1)=(\Delta\cap M)_0-\{0\}$。这恰好是\hyperref[ex:degree-four-moduli]{例题~\ref*{ex:degree-four-moduli}} 的集合 $\Xi$；注意 $\Xi=\{v_1,v_2,v_3,v_4,v_5,v_6\}$。锥 $\operatorname{cpl}(\Sigma^\circ)$ 是 $\Xi$ 的 GKZ 分解的一部分。前面用 $\Xi^+=(\Xi\cup\{0\})\times\{1\}=(\Delta\cap M)_0\times\{1\}$ 构造了次级扇~\eqref{eq:second-moduli-secondary-fan}。改用 $\Xi$ 重复该构造，就给出 GKZ 分解。回顾 $A(\Xi)^*$ 是 $\Xi$ 各元素之间的线性关系向量空间。在本例中，以下矩阵的各行给出 $A(\Xi)^*$ 的一组基：
\[
\begin{pmatrix}-2&-2&-4&1&0&7\\1&1&2&0&1&-3\end{pmatrix}.
\]
使用 $A(\Xi)$ 的对偶基，该矩阵的各列给出向量 $e_{v_1}^*,\ldots,e_{v_6}^*\in A(\Xi)$。\hyperref[sec:gkz-decomposition]{第 3.4 节}于是表明，GKZ 分解如下图所示：
\begin{figure}[H]
\centering
\refstepcounter{equation}\label{eq:revisited-gkz-fan}
\makebox[\textwidth]{\makebox[0pt][l]{(\theequation)}\hfill
\raisebox{-.5\height}{\includegraphics[width=0.43375\textwidth]{6-2-gkz-fan.png}}\hfill\mbox{}}
\end{figure}
每条边都标出张成它的点 $e_{v_i}^*$。GKZ 分解有三个极大锥，记为 $\sigma_1,\sigma_2,\sigma_3$。$\Xi$ 的线性回路为
\begin{equation}\label{eq:revisited-circuits}\index{线性回路!例题|hyperpage}
\begin{gathered}
\{v_1,v_2,v_3,v_5,v_6\},\quad\{v_1,v_2,v_3,v_4,v_6\},\\
\{v_4,v_5,v_6\},\quad\{v_1,v_2,v_3,v_4,v_5\},
\end{gathered}
\end{equation}
相应的线性关系为
\begin{equation}\label{eq:revisited-circuit-relations}
\begin{aligned}
v_1+v_2+2v_3+v_5-3v_6&=0,\\
-2v_1-2v_2-4v_3+v_4+7v_6&=0,\\
v_4+2v_5+v_6&=0,\\
v_1+v_2+2v_3+3v_4+7v_5&=0.
\end{aligned}
\end{equation}
观察式~\eqref{eq:revisited-gkz-fan} 最容易找出这些线性回路。每个回路都可以如下得到：沿图中一条边确定的直线，将平面分为两个半平面。$\Xi'$ 由不在该直线上的顶点组成；$\Xi'_+$ 的点位于直线一侧，$\Xi'_-$ 的点位于另一侧。

很容易看出，唯一满足 $|\Xi'_-|=1$ 的线性回路 $\Xi'$（使用式~\eqref{eq:revisited-circuit-relations} 的关系或其相反关系），是按所写符号取的第一个关系。因此，Calabi--Yau 锥 $C_{\mathrm{CY}}$ 是 $A^+(\Xi)=\sigma_1\cup\sigma_2\cup\sigma_3$ 中由单个不等式 $a_1+a_2+2a_3+a_5-3a_6\geq0$ 确定的锥；与通常一样，$A^+(\Xi)$ 的一般点表示为 $[\sum_i a_iD_i]$。得到 $C_{\mathrm{CY}}=\sigma_1\cup\sigma_2$。

接着构造对应于 $\sigma_1,\sigma_2$ 的扇。先从 $\sigma_2$ 开始。一种做法是在 $\sigma_2$ 中取一个除子类，用丰沛性的凸性条件确定扇（\SourceCite{KS} 已将这一过程自动化）。不过，更容易的办法是认识到 $V^\circ$ 来自\IndexLocation{indexloc:1432:163}加权射影空间 $\mathbb P(1,1,2,3,7)$（最初正是这样构造 $V^\circ$ 的）。式~\eqref{eq:revisited-polytope} 的六个顶点中，前五个确定一个单纯形；取该单纯形各个面上的锥，给出
\SourcePage{164}
$M_{\mathbb R}$ 中的一个扇。对应环面簇为 $\mathbb P(1,1,2,3,7)$，它不是 Fano 簇。不过，加入由 $(0,0,-1,-2)$ 生成的一维锥来爆破 $\mathbb P(1,1,2,3,7)$ 后，就得到单纯扇 $\Sigma^\circ$，其极大维锥为
\[
\begin{gathered}
C(v_2,v_3,v_4,v_5),\ C(v_1,v_3,v_4,v_5),\ C(v_1,v_2,v_4,v_5),\ C(v_1,v_2,v_3,v_4),\\
C(v_2,v_3,v_5,v_6),\ C(v_1,v_3,v_5,v_6),\ C(v_1,v_2,v_5,v_6),\ C(v_1,v_2,v_3,v_6).
\end{gathered}
\]
这个扇显然细分 $\Delta^\circ$ 的法扇。还可验证，$\Sigma^\circ$ 支撑于线性回路

\[
\Xi'_1=\{v_1,v_2,v_3,v_5,v_6\},\qquad\Xi'_2=\{v_1,v_2,v_3,v_4,v_6\}.
\]
（符号由式~\eqref{eq:revisited-circuit-relations} 确定）。不等式 $\ell_{\Xi'_i}(D)\geq0$（$i=1,2$）定义 $\sigma_2$，因此 $\operatorname{cpl}(\Sigma^\circ)=\sigma_2$。

此外，$\Xi'_2$ 满足 $|(\Xi'_2)_+|=2$、$|(\Xi'_2)_-|=3$，所以得到扇 $\Sigma^{\circ\prime}=\operatorname{flip}_{\Xi'_2}(\Sigma^\circ)$ 及双有理映射 $\operatorname{flop}_{\Xi'_2}$。利用式~\eqref{eq:revisited-gkz-fan}，很容易验证 $\operatorname{cpl}(\Sigma^{\circ\prime})=\sigma_1$。还注意到，$\Xi'_2$ 包含在 $\Delta$ 的一个刻面中，该刻面由 $M_{\mathbb R}$ 中方程
\[
\langle m,(-1,-1,-1,1)\rangle=-1
\]
确定的超平面支撑，对应于顶点 $(-1,-1,-1,1)\in\Delta^\circ$（见\hyperref[ex:degree-four-moduli]{例题~\ref*{ex:degree-four-moduli}}）。由前面的讨论，$\Sigma^\circ,\Sigma^{\circ\prime}$ 都是简化射影细分。

现在证明，$\Sigma^{\circ\prime}$ 是 $\Sigma^\circ$ 的平凡翻转。因为二者都是简化射影细分，所以处于式~\eqref{eq:simplified-subdivision-diagram} 的情形。此外，可验证
$\operatorname{flop}_{\Xi'_2}:X_{\Sigma^\circ}\dashrightarrow X_{\Sigma^{\circ\prime}}$
除 $\mathbb P_{\Delta^\circ}$ 中与顶点 $(-1,-1,-1,1)\in\Delta^\circ$ 对应的点的上方外，为同构。一般 $\overline{V^\circ}\subset\mathbb P_{\Delta^\circ}$ 避开这一点，因此一般 $V^\circ\subset X_{\Sigma^\circ}$ 避开 $\operatorname{flop}_{\Xi'_2}$ 的基点集。所以，$X_{\Sigma^\circ},X_{\Sigma^{\circ\prime}}$ 中的一般 Calabi--Yau 超曲面同构，$\Sigma^{\circ\prime}$ 是 $\Sigma^\circ$ 的平凡翻转。

这样就证明了 $\widetilde{K(V^\circ)}_{\mathrm{toric}}=\sigma_1\cup\sigma_2$。下一步是证明它等于 $\overline{K(V^\circ)}_{\mathrm{toric}}$，即本例中猜想~\ref{conj:toric-kahler-cone} 成立。取 $D\in\overline{K(V^\circ)}_{\mathrm{toric}}$。先将 $V^\circ$ 视为 $X_{\Sigma^\circ}$ 的子簇，用其齐次坐标在 $V^\circ$ 上施加方程 $x_1=x_6=0$，定义曲线 $C\subset V^\circ$。再将 $V^\circ$ 视为 $X_{\Sigma^{\circ\prime}}$ 的子簇，用后者的齐次坐标及方程 $x_1=x_2=0$ 定义曲线 $C'\subset V^\circ$。因为 $D\in\overline{K(V^\circ)}_{\mathrm{toric}}$，有 $C\cdot D\geq0$ 和 $C'\cdot D\geq0$。可以证明，这两个不等式恰好定义 $\sigma_1\cup\sigma_2$。事实上，第一个不等式等价于 $\ell_{\Xi'_1}(D)\geq0$，其中 $\Xi'_1$ 为上述由 $\Sigma^\circ$ 支撑的线性回路；第二个不等式等价于 $\ell_{\Xi'_3}(D)\geq0$，其中 $\Xi'_3$ 是式~\eqref{eq:revisited-circuits} 的第三个回路，符号取式~\eqref{eq:revisited-circuit-relations} 所写的。因此，$\overline{K(V^\circ)}_{\mathrm{toric}}\subset\sigma_1\cup\sigma_2$，随即得到等号。

一般情形下，证明猜想~\ref{conj:toric-kahler-cone} 的一种思路，是构造类似上述 $C,C'$ 的曲线，使它们给出的线性泛函，与定义 $\widetilde{K(V)}_{\mathrm{toric}}$ 所用各细分的适当回路给出的泛函 $\ell_{\Xi'}$ 一致。

现在已得到 Kähler 锥的闭包
$\overline{K(V^\circ)}=\overline{K(V^\circ)}_{\mathrm{toric}}=\sigma_1\cup\sigma_2$，接着讨论 $V^\circ$ 的 Kähler 模空间。由\hyperref[subsec:large-radius-limits]{第 6.2.1 小节}，$\mathcal{KM}$ 的一个\IndexLocation{indexloc:0893:164}部分紧化由 $\overline{K(V^\circ)}$ 的仿射环面簇给出，很容易看出它光滑。此外，\IndexLocation{indexloc:0698:164}大半径极限点附近的自然坐标由 $\overline{K(V^\circ)}$ 的生成元确定。所需生成元 $D,D'$ 为式~\eqref{eq:revisited-gkz-fan} 的向量 $D=(-2,1)=e_{v_1}^*$、$D'=(1,0)=e_{v_4}^*$，所以 $D=D_1$ 且
\SourcePage{165}
$D'=D_4$，这里 $D_i$ 表示与 $v_i$ 对应的环面除子在 $V^\circ$ 上的限制。可以计算交数
\[
D^3=0,\qquad D^2\cdot D'=1,\qquad D\cdot D'^2=3,\qquad D'^3=9,
\]
由此得到 $V^\circ$ 上同调中的关系 $D^3=0$、$D'^2-3D\cdot D'=0$。下面将用到这些关系。

下一步是与镜像 $V$ 的复结构模比较。\hyperref[ex:degree-four-moduli]{例题~\ref*{ex:degree-four-moduli}} 用次级扇~\eqref{eq:second-moduli-secondary-fan} 紧化了 $\mathcal M_{\mathrm{simp}}$。现在知道，需要将锥 $\sigma_1,\sigma_2$ 合并为一个锥，这对应于爆缩原来定义的 $\overline{\mathcal M}_{\mathrm{simp}}$。这样，Kähler 锥 $\overline{K(V^\circ)}$ 就对应于紧化后的复结构模空间中的一个边界点；可以证明它为最大幺幂。事实上，使用锥 $\overline{K(V^\circ)}$ 的生成元提供的自然坐标，就得到两个单值性变换，其对数为 $N_1,N_2$。计算得 $N_1^3=0$、$N_2^2-3N_1N_2=0$。\footnote{计算使用了 \SourceCite{CdFKM,CFKM} 的方法；在这一用途下，它们效率更高。}结合相应周期空间的维数与非零最高阶项，足以证明极小爆缩中的相应边界点为最大幺幂。注意，它们与 $D,D'$ 满足的关系 $D^3=0$、$D'^2-3D\cdot D'=0$ 相似。这并非偶然，因为镜像映射应当将单值性的对数，与同 $\overline{K(V^\circ)}$ 中相应锥的生成元作杯积对应起来。于是得到本例中镜像对称如何起作用的一个显式例子。
\end{example}
\endgroup

接着转向 Kähler 模空间的整体描述。目标是使它与\hyperref[subsec:toric-complex-moduli]{第 6.1.2 小节}中镜像流形的简化多项式复结构模空间的整体描述，直接按定义相同。

更正式地，设 $V\subset X_\Sigma$ 为 Calabi--Yau 环面超曲面，$\Sigma$ 为自反多胞形 $\Delta$ 的法扇的一个简化射影细分，并且如通常一样，$\Xi=(\Delta^\circ\cap N)_0-\{0\}$。将\emph{紧化环面 Kähler 模空间} $\overline{\mathcal{KM}}_{\mathrm{toric}}$ 定义为 $A(\Xi)$ 中以下扇对应的环面簇：从 $\Xi^+=(\Delta^\circ\cap N)_0\times\{1\}$ 的次级扇出发，凡一组锥 $\operatorname{cpl}(\Sigma')$ 可通过一系列\IndexLocation{indexloc:0331:165}平凡翻转联系起来，就用它们的并集代替这一组锥（如猜想~\ref{conj:toric-kahler-cone} 所述）。

因此，紧化环面模空间与镜像族的紧化复结构模空间 $\overline{\mathcal M}_{\mathrm{simp}}$ 经极小爆缩后的空间，本质上相同。注意，按定义，紧化 Kähler 模空间包含与 Kähler 锥对应的仿射环面簇。\hyperref[subsec:large-radius-limits]{第 6.2.1 小节}已经说明，这个环面簇自然包含 Kähler 模空间。如前所述，忽略了自同构群的作用；在始终假定猜想~\ref{conj:toric-kahler-cone} 成立的前提下，它将给出 $\overline{\mathcal{KM}}_{\mathrm{toric}}$ 的一个有限商。

$\overline{\mathcal{KM}}_{\mathrm{toric}}$ 的定义有一个很好的特点：它自然地扩大 Kähler 模空间，同时将它紧化。\hyperref[subsec:threefold-kahler-moduli]{第 6.2.2 小节}中的一般 Calabi--Yau 三维簇，必须先扩大到可动锥，再进一步扩大到\IndexLocation{indexloc:0914:165}反射可动锥；即便如此，也只得到\IndexLocation{indexloc:0893:165}部分紧化。因此，虽然描述环面 Kähler 模空间花了一些功夫，Calabi--Yau 环面超曲面仍然比一般 Calabi--Yau 三维簇容易处理得多，更不用说任意维数的 Calabi--Yau 簇了。

\SourcePage{166}
紧化环面 Kähler 模空间与\hyperref[subsec:toric-reduction]{第~\ref*{subsec:toric-reduction} 小节}中讨论的矩映射，为规范线性 $\sigma$ 模型（GLSM）的相结构提供了数学描述；\hyperref[sec:gauged-linear-sigma-models]{附录 B.5} 将概述这一模型。GLSM 依赖于一些参数，可将它们与 $H^2(X_\Sigma,\mathbb R)$ 的元素识别起来。该空间自然识别为矩映射的目标空间。与讨论辛约化时一样，在上同调类 $D\in A^+(\Xi)$ 上方，对矩映射的一个纤维取商，得到辛商。若 $D$ 位于 $A^+(\Xi)$ 的 GKZ 分解中 $\operatorname{cpl}(\Sigma)$ 的内部，则定理~\ref{thm:gkz-symplectic-reduction} 已证明该辛商就是 $X_\Sigma$。对一般 $D\in A(\Xi)$，辛商仍有定义，但未必为射影环面簇。由此，可能的辛商的同构类与次级扇中的极大维锥一一对应。

GLSM 的额外要素是从辛商中得到一个子簇 $V$。在这里，对应于次级扇不同部分的不同辛商，称为 GLSM 的\emph{相}。若 $D$ 位于某个简化射影细分对应的 $\operatorname{cpl}(\Sigma)$ 的内部，则子簇 $V$ 为 Calabi--Yau 超曲面。一般而言，$V$ 甚至未必为超曲面。\hyperref[sec:gauged-linear-sigma-models]{附录 B.5} 的例子有两个相：相应辛商为 $\mathbb P^4$ 和 $\mathbb C^5/\mathbb Z_5$，对应的子簇 $V$ 为 Fermat 五次三维簇和原点。后一情形是一个 \emph{Landau--Ginzburg 轨形}的例子，\hyperref[sec:landau-ginzburg-models]{附录 B.4} 也将讨论它。

环面 Kähler 模空间的定义带来的一个结果是：若 $V\subset X_\Sigma$ 和 $V^\circ\subset X_{\Sigma^\circ}$ 为 Batyrev 镜像，则 $\overline{\mathcal{KM}}_{\mathrm{toric}}(V)$ 与 $\overline{\mathcal M}_{\mathrm{simp}}(V^\circ)$ 的相应极小爆缩，是同一个空间。这可能使读者觉得镜像映射就是恒等映射；不过，如例题~\ref{ex:weighted-extremal-transition} 将展示的，镜像映射并没有这么简单。真正的镜像映射将在\hyperref[sec:mirror-map]{第~\ref*{sec:mirror-map} 节}详细讨论。

\hypertarget{source:166-batch4}{}
\subsection{Kähler 模空间的边界}\label{subsec:kahler-moduli-boundary}
本小节的讨论仍然在很大程度上具有猜测性质。我们先考察 Calabi--Yau 三维簇，再专门讨论作为环面超曲面出现的情形。出发点是以下观察：常常可以将一个已知的 Calabi--Yau 三维簇爆缩为奇异簇，再将它光滑化，得到一个新的 Calabi--Yau 三维簇。更准确地说，\SourceCite{Morrison9} 给出了以下定义。

\begin{definition}\label{def:extremal-transition}\index{极端跃迁|hyperpage}
设 $V$ 为 Calabi--Yau 三维簇，$\phi:V\to\overline V$ 为双有理收缩。若 $\overline V$ 可以光滑化为 Calabi--Yau 三维簇 $\widetilde V$，则从 $V$ 到 $\widetilde V$ 的这一过程称为\emph{极端跃迁}。
\end{definition}

若 $\phi$ 为本原双有理收缩，为了强调这一点，有时称它为\emph{本原极端跃迁}。回忆\hyperref[subsec:threefold-kahler-moduli]{第 6.2.2 小节}，$\phi$ 分为 I、II、III 三型，因此，例如可以说一个极端跃迁属于 I 型。

在物理学文献中，I 型极端跃迁称为\emph{锥形奇点跃迁}，因为物理学家常将三维结点称为锥形奇点。通常所说的“Reid 的幻想”中出现的跃迁，正是这种锥形奇点跃迁。在 \SourceCite{Reid5} 中，Reid 猜想：若将所有可能的 $V$ 与 $\widetilde V$ 的复模空间，沿所有可能的 $\overline V$ 的模空间粘合起来，所得空间或许是连通的。

下面是取自 \SourceCite{GMS} 的一个漂亮的锥形奇点跃迁例子。

\SourcePage{167}
\SourceExampleNumber{6.2.4}{1}
\begin{example}\label{ex:quintic-conifold-transition}
设 $\overline V\subset\mathbb P^4$ 为一个一般的奇异五次三维簇，包含 $\mathbb P^4$ 中由 $x_3=x_4=0$ 定义的平面。这样的 $\overline V$ 的方程为
$x_3g(x_0,\ldots,x_4)+x_4h(x_0,\ldots,x_4)=0$，其中 $g,h$ 为一般四次多项式。超曲面 $\overline V$ 在 $x_3=x_4=g(x_0,\ldots,x_4)=h(x_0,\ldots,x_4)=0$ 的 16 个点处有结点。考虑 $\mathbb P^4\times\mathbb P^1$ 中由
\[
y_0x_4-y_1x_3=0,\qquad
y_0g(x_0,\ldots,x_4)+y_1h(x_0,\ldots,x_4)=0
\]
定义的光滑 Calabi--Yau 完全交 $V$；这里 $y_0,y_1$ 为 $\mathbb P^1$ 的齐次坐标，$x_0,\ldots,x_4$ 仍为 $\mathbb P^4$ 的齐次坐标。容易看出，$V$ 的 Kähler 锥同构于 $\mathbb P^4\times\mathbb P^1$ 的 Kähler 锥，即由 $\mathbb P^4$ 与 $\mathbb P^1$ 的超平面类的拉回生成。来自 $\mathbb P^4$ 的超平面类 $H$ 张成 Kähler 锥的一个面，且 $H^3=5\ne0$，所以投影 $\mathbb P^4\times\mathbb P^1\to\mathbb P^4$ 限制为本原收缩 $\phi:V\to\overline V$。由于 $\phi$ 收缩的是 16 个结点上方的曲线，它属于 I 型。奇异三维簇 $\overline V$ 可光滑化为光滑五次三维簇 $\widetilde V\subset\mathbb P^4$，从而完成这一锥形奇点跃迁。
\end{example}

我们组织从 $V$ 到 $\widetilde V$ 的极端跃迁的讨论时，所要提示的是：同时粘合 $V$ 与 $\widetilde V$ 的复模空间和 Kähler 模空间，比只粘合复模空间更自然。换言之，我们是在完整的 SCFT 模空间中工作。下面将看到，识别 $V$ 与 $V^\circ$ 的 SCFT 模空间，会带来出人意料的丰富结构。

在 SCFT 模空间中，一个极端跃迁自然可看作两步过程：先固定复结构，使 Kähler 结构发生形变；再固定 Kähler 结构，使复结构发生形变。具体地，根据\hyperref[subsec:threefold-kahler-moduli]{第 6.2.2 小节}的讨论，$\overline V$ 的丰沛类空间经 $\phi$ 拉回，成为 $\overline{K(V)}$ 的一个面 $F$。我们先在 $\overline{\mathcal{KM}}$ 中选取一条趋近于其边界中与 $F$ 关联的分支的路径，同时保持复结构不变。然后，将这一边界分支与 $V$ 的复模空间的乘积，粘合到 $\widetilde V$ 的紧化 Kähler 模空间与 $\overline V$ 的复模空间的乘积上。注意，$\overline V$ 的复模空间是 $\widetilde V$ 的复模空间边界的一个子集。接下来在 $\widetilde V$ 的复模空间中自由移动，将 $\overline V$ 光滑化，即可完成跃迁。

\SourceCite{Morrison9} 猜想，这一跃迁的“镜像”是 $V^\circ$ 的 SCFT 模空间中沿反方向进行的极端跃迁。镜像对称应将 $V$ 的 SCFT 模空间中的路径，与 $V^\circ$ 的 SCFT 模空间中的路径识别起来。若 $\widetilde V^\circ$ 为 $\widetilde V$ 的镜像，那么描述这一跃迁的路径应当如下：先在 $V^\circ$ 的复模空间中移动，固定 Kähler 结构；其极限奇异复结构，则来自 $\widetilde V^\circ$ 的一个本原收缩，可以表示为 $\widetilde V^\circ$ 的 Kähler 模空间中的路径。猜想从 $V^\circ$ 到 $\widetilde V^\circ$ 的跃迁，是一个极端跃迁的逆过程。

正如 \SourceCite{BKK,Morrison9} 所观察到的，Batyrev 镜像构造提供了这一机制的一个漂亮例子。

\SourceExampleNumber{6.2.4}{2}
\begin{example}\label{ex:toric-extremal-transition}
设 $\Delta$ 与 $\widetilde\Delta$ 为两个不同的自反多胞形，且 $\Delta\subset\widetilde\Delta$。为简便起见，假定二者都为四维。我们先用这些数据构造一个从 $V$ 到 $\widetilde V$ 的极端跃迁。

\SourcePage{168}
由于 $\widetilde\Delta^\circ\subset\Delta^\circ$，$\widetilde\Delta$ 的法扇的任一细分 $\widetilde\Sigma$，都可进一步细化为 $\Delta$ 的法扇的一个细分 $\Sigma$。这一细分给出态射 $\pi:X_\Sigma\to X_{\widetilde\Sigma}$。Calabi--Yau 超曲面 $V\subset X_\Sigma$ 由方程
\[
\sum_{m_i\in\Delta\cap M}\lambda_i t^{m_i}=0
\]
定义。类似地，由方程
\begin{equation}\label{eq:extremal-transition-polynomial}
\sum_{m_i\in\widetilde\Delta\cap M}\lambda_i t^{m_i}=0
\end{equation}
得到 Calabi--Yau 超曲面 $\widetilde V\subset X_{\widetilde\Sigma}$。令 $\phi=\pi|_V$，即得到双有理收缩 $\phi:V\to\overline V$；其中 $\overline V$ 由式~\eqref{eq:extremal-transition-polynomial} 定义，只是当 $m_i\in(\widetilde\Delta-\Delta)\cap M$ 时，所有系数 $\lambda_i$ 都为零。让式~\eqref{eq:extremal-transition-polynomial} 中全部 $\lambda_i$ 任意取值，就将 $\overline V$ 光滑化为 $\widetilde V$，完成极端跃迁。从 $V$ 到 $\widetilde V$ 的这一过程称为\emph{环面极端跃迁}。

关键的观察在于：由于 $\widetilde\Delta^\circ\subset\Delta^\circ$，同样的构造给出镜像之间的环面极端跃迁，不过方向是从 $\widetilde V^\circ$ 到 $V^\circ$。这恰好符合 Morrison 猜想的预言。
\end{example}

还应提到，\hyperref[ex:toric-extremal-transition]{例题~\ref*{ex:toric-extremal-transition}}的方法，已经应用于\hyperref[sec:batyrev-mirror]{第 4.1 节}讨论的\IndexLocation{indexloc:0109:168}加权射影空间超曲面所对应的 7555 个\IndexLocation{indexloc:1218:168}自反多胞形。可以证明 \SourceCite{CGGK,ACJM}，用环面极端跃迁连接模空间的各个分支后，所得空间的这 7555 个不可约分支全部位于同一个连通分支中。

一般地，设有一个从 $V$ 到 $\widetilde V$ 的极端跃迁。对于 SCFT 模空间中与 $V$、$\widetilde V$ 关联的两个模空间分支交会之处的结构，目前还没有良好的数学定义。现有的解释只有物理学中的解释：在 I 型极端跃迁的情形下，利用 II 型弦理论和无质量黑洞 \SourceCite{GMS}。II 型弦理论的模空间大于 SCFT 模空间；极端跃迁将这些更大的模空间粘合起来，物理学给出的描述如下。被 $V\to\overline V$ 收缩的曲线对应于黑洞。当我们沿极端跃迁趋近 $V$ 的 Kähler 模空间边界时，这些黑洞变为无质量。另一方面，II 型弦理论中的上同调类也对应于无质量粒子，而 $V$ 与 $\widetilde V$ 的 Hodge 数之差，正由这些无质量黑洞说明。基于 \SourceCite{KMP} 所解释的物理，对某些 III 型极端收缩也能给出类似说明。第 \hyperref[subsec:conifold-transitions-conclusion]{\ref*{subsec:conifold-transitions-conclusion}} 小节将简要回到这些思想。（这里所说的 II 型弦理论，在英文中使用小写的“type”；它与极端收缩的 I、II、III 型，即英文使用大写“Type”的分类，没有直接关系。）

接下来讨论取自 \SourceCite{CdFKM,KMP} 的一个极端跃迁例子。它揭示了有关模空间细致而丰富的结构。令人意外的是，这个例子我们已经很熟悉了。

\SourceExampleNumber{6.2.4}{3}
\begin{example}\label{ex:weighted-extremal-transition}
继续讨论\hyperref[ex:weighted-pf-method]{例题~\ref*{ex:weighted-pf-method}}和\hyperref[ex:weighted-complex-moduli]{例题~\ref*{ex:weighted-complex-moduli}}开始研究的镜像族，它来自 $\mathbb P(1,1,2,2,2)$ 的环面爆破中的反典范超曲面。
\SourcePage{169}
有自反多胞形 $\Delta$，其极多胞形 $\Delta^\circ$ 满足
\[
\begin{aligned}
\Xi=(\Delta^\circ\cap N)_0-\{0\}=\{&(-1,-2,-2,-2),(1,0,0,0),(0,1,0,0),\\
&(0,0,1,0),(0,0,0,1),(0,-1,-1,-1)\}.
\end{aligned}
\]
将这些点记为 $m_1,\ldots,m_6$。\hyperref[ex:weighted-resolution]{例题~\ref*{ex:weighted-resolution}}中的 GKZ 分解表明，满足 $\Sigma(1)=\Xi$ 的简化射影细分 $\Sigma$ 唯一。它给出 Calabi--Yau 三维簇 $V\subset X_\Sigma$。我们将构造一个从 $V$ 出发的极端跃迁，并研究这一跃迁在 $V$ 的 Kähler 模空间中的表现。

利用\hyperref[sec:batyrev-mirror]{第 4.1 节}的公式，可以算出
\[
h^{1,1}(V)=h^{1,1}_{\mathrm{toric}}(V)=2,\qquad
h^{2,1}(V)=86,\quad h^{2,1}_{\mathrm{poly}}(V)=83.
\]
因此 Batyrev 镜像 $V^\circ$ 满足
\[
h^{2,1}(V^\circ)=h^{2,1}_{\mathrm{poly}}(V^\circ)=2,\qquad
h^{1,1}(V^\circ)=86,\quad h^{1,1}_{\mathrm{toric}}(V^\circ)=83.
\]
令 $D_1,\ldots,D_6$ 为与 $\Sigma(1)=\Xi$ 关联的环面除子。由于 $H^2(X_\Sigma)\simeq H^2(V)$，既可将它们视为 $X_\Sigma$ 上的类，也可视为 $V$ 上的类。$V$ 的 Kähler 锥由 $D_1,D_3$ 的类张成。\hyperref[ex:weighted-complex-moduli]{例题~\ref*{ex:weighted-complex-moduli}}中矩阵的前六列给出线性等价关系
$D_2\sim D_1$、$D_4\sim D_5\sim D_3$、$D_6\sim D_3-2D_1$，它们对后续计算很有用。特别地，由此得到 $V\in|\sum_iD_i|=|4D_3|$。

接着寻找 $V$ 的双有理收缩。利用\hyperref[subsec:threefold-kahler-moduli]{第 6.2.2 小节}的理论，研究 Kähler 锥的各个面。使用 $X_\Sigma$ 中与 $\Sigma(1)=\Xi$ 对应的齐次坐标 $x_1,\ldots,x_6$ 很方便。可验证，$|D_1|$ 是由截面 $x_1,x_2$ 定义的线性系，它们给出 K3 纤维化 $V\to\mathbb P^1$，且 $D_1^3=0$。因此 $|D_1|$ 不给出双有理收缩。另一方面，$|D_3|$ 由截面 $x_1^2x_6,x_1x_2x_6,x_2^2x_6,x_3,x_4,x_5$ 定义。由于 $D_3^3=8$，得到本原收缩 $\phi=\phi_{|D_3|}$，其像为 $\mathbb P^5$ 中次数为 8 的三维簇 $\overline V$。事实上，$\overline V$ 是一个二次超曲面与一个四次超曲面的完全交。为看清这一点，在 $\mathbb P^5$ 上引入坐标 $y_0,\ldots,y_5$。按上述顺序用 $|D_3|$ 的截面定义 $\phi$，可见 $\overline V$ 位于秩为 3 的奇异二次超曲面 $y_0y_2=y_1^2$ 上；而定义 $V\subset X_\Sigma$ 的 $|4D_3|$ 的截面，由 $y_i$ 的一个四次多项式诱导。可以验证，$\phi$ 为 III 型本原收缩，将例外曲面 $x_6=0$ 收缩到 $\overline V$ 中由 $y_0=y_1=y_2=0$ 截出的轨迹；它实际上是一条平面四次曲线。例外除子 $x_6=0$ 是这条亏格为 3 的曲线上的直纹曲面。将 $\overline V$ 光滑化为 $\mathbb P^5$ 中光滑的 $(2,4)$ 完全交 $\widetilde V$，便完成了本原极端跃迁。$\widetilde V$ 的有关 Hodge 数为 $h^{1,1}(\widetilde V)=1$、$h^{2,1}(\widetilde V)=89$。

在 Kähler 模的语言中，这一跃迁描述如下：由于 $\overline V$ 来自 $V$ 的 Kähler 锥的一个面，显然应将 $\overline{\mathcal{KM}}(\widetilde V)$ 沿这一面对应的子集，接到 $\overline{\mathcal{KM}}(V)$ 上。不过，这还没有交代全部情形，因为 Kähler 锥的一个面的实余维数为一，而 $\overline{\mathcal{KM}}(\widetilde V)\subset\overline{\mathcal{KM}}(V)$ 的复余维数为一。因此，还需要进一步确定粘接 $\overline{\mathcal{KM}}(\widetilde V)$ 的精确轨迹。下面将看到，镜像会告诉我们该怎样做。

\SourcePage{170}
至此，SCFT 模空间的图景已近乎完整。迄今我们只看了 $V$ 的复模空间中多项式的部分，但 $V$ 的 Hodge 数表明，存在非多项式形变。例外除子 $x_6=0$ 在 $V$ 的一般复结构形变下不会继续存在；不过，由于它是亏格为 3 的曲线上的直纹曲面，有 $2g-2=4$ 条纤维会随复结构的一般形变而形变。在这一邻近复结构上，形变后的 $|D_3|$ 仍会诱导本原收缩，但这次属于 I 型。此外，这 4 条孤立曲线可以同时翻越，得到不同的复结构和 Kähler 结构。可以证明，这一翻越诱导 SCFT 模空间上的一个对合。

于是得到 $V$ 的 SCFT 模空间的如下描述。对于 86 维复模空间中的任一复结构，都有一个二维 Kähler 模空间，其中含有与本原收缩 $\phi_{|D_3|}$ 对应的一条特殊曲线。另外，考虑刚才构造的 SCFT 模空间上的翻越对合：83 维的多项式子空间逐点不动，而它在 Kähler 模空间上的作用非平凡，是\hyperref[subsec:threefold-kahler-moduli]{第 6.2.2 小节}所述的反射。$\widetilde V$ 的 89 维复模空间中，包含定义它的二次式秩为 3 的 $(2,4)$ 完全交所构成的 83 维子空间，以及二次式秩为 4 的 $(2,4)$ 完全交所构成的 86 维子空间。它们分别与 $V$ 的多项式模空间及复模空间识别；$\widetilde V$ 的 Kähler 模空间则与 $V$ 的 Kähler 模空间中的一条曲线识别，下面会描述这条曲线。

接着考察从 $\widetilde V^\circ$ 到 $V^\circ$ 的“对偶”极端跃迁。\hyperref[ex:weighted-pf-method]{例题~\ref*{ex:weighted-pf-method}}和\hyperref[ex:weighted-complex-moduli]{例题~\ref*{ex:weighted-complex-moduli}}已描述 $V^\circ$；$\widetilde V^\circ$ 则可用\hyperref[sec:toric-complete-intersections]{第 4.3 节}讨论的 Batyrev--Borisov 构造得到。这样，根据 \SourceCite{CdFKM} 的计算，可以显式写出从 $\widetilde V^\circ$ 到 $V^\circ$ 的极端跃迁。它诱导模空间上的映射 $\overline{\mathcal M}_{\mathrm{simp}}(\widetilde V^\circ)\to\overline{\mathcal M}_{\mathrm{simp}}(V^\circ)$。利用 $\overline{\mathcal M}_{\mathrm{simp}}$ 上的坐标 $z_1,z_2$，可证明 \SourceCite{CdFKM,KMP}，$\overline{\mathcal M}_{\mathrm{simp}}(\widetilde V^\circ)$ 沿\hyperref[ex:weighted-complex-moduli]{例题~\ref*{ex:weighted-complex-moduli}}中由 $D_{A\cap\Gamma}(f_\Gamma)=1-4z_2=0$ 定义的判别式轨迹分支粘接。\hyperref[subsec:toric-mirror-map]{第~\ref*{subsec:toric-mirror-map} 小节}将简要说明这一论证。判别式轨迹的出现并不意外，因为从 $\widetilde V^\circ$ 到 $V^\circ$ 的过程，首先收缩到一个奇异簇。

由于从 $\widetilde V^\circ$ 到 $V^\circ$ 的跃迁，经镜像对称与从 $V$ 到 $\widetilde V$ 的跃迁相容，现在终于能完成对 $V$ 的 SCFT 模空间的描述。即用镜像映射将 $\overline{\mathcal M}_{\mathrm{simp}}(V^\circ)$ 与 $\overline{\mathcal{KM}}(V)$ 识别后，根据上一段，$\widetilde V$ 的 Kähler 模空间应与 $V$ 的 Kähler 模空间中由 $1-4z_2=0$ 定义的曲线识别起来。

$V^\circ$ 的 SCFT 模空间具有与 $V$ 相同的复杂几何。此外，也能看出翻越对合的作用。由定理~\ref{thm:batyrev-hodge} 的证明，$V^\circ$ 构造为环面簇 $X_{\Sigma^\circ}$ 中的超曲面；该环面簇含有 3 个环面除子 $D_i$，每个限制到 $V^\circ$ 后都有两个分支 $D_i',D_i''$。除子类 $D_i'+D_i''$ 属于 $H^{1,1}_{\mathrm{toric}}(V^\circ)$，但单独的 $D_i',D_i''$ 不属于这个子空间。交换 $D_i'$ 与 $D_i''$（$i=1,2,3$）的对合，对应于翻越对合在 $V$ 的复模空间上的作用。它在 $H^{1,1}_{\mathrm{toric}}(V^\circ)$ 上的作用平凡。
\end{example}

\SourcePage{171}
这个例子对环面情形的镜像映射有若干有趣的启示。通过次级扇的构造，Batyrev 镜像 $V,V^\circ$ 使 $\overline{\mathcal{KM}}_{\mathrm{toric}}(V)$ 与 $\overline{\mathcal M}_{\mathrm{simp}}(V^\circ)$ 成为同一个模空间。这似乎暗示镜像映射是恒等映射。然而，仔细考察\hyperref[ex:weighted-extremal-transition]{例题~\ref*{ex:weighted-extremal-transition}}就知道并非如此。基本思想是：要使镜像对称与极端跃迁相容，镜像映射就必须不同于恒等映射。

具体说明如下。回到\hyperref[ex:weighted-extremal-transition]{例题~\ref*{ex:weighted-extremal-transition}}，$V$ 的 Kähler 锥闭包为 $\{t_1D_3+t_2D_1:t_1,t_2\geq0\}$。注意，这里的 $t_1,t_2$ 是\hyperref[chap:toric]{第~\ref*{chap:toric} 章}式~\eqref{eq:weighted-gkz-picture} 中记为 $r_1,r_2$ 的坐标。因此，Kähler 锥的面由 $t_j=0$ 定义。转到 Kähler 模空间时，让 $t_j$ 位于上半平面，则 $q_j=e^{2\pi it_j}$ 是 $\mathcal{KM}(V)$ 的相应局部坐标。因此，$\operatorname{Im}(t_j)=0$ 对应于 $V$ 的 Kähler 模空间中的 $|q_j|=1$。注意，这个集合在 $\overline{\mathcal{KM}}(V)$ 中的实余维数为一。

镜像一侧有 $V^\circ$ 的简化复模，它由同一个空间给出。\hyperref[ex:weighted-complex-moduli]{例题~\ref*{ex:weighted-complex-moduli}}中的自然坐标 $z_1,z_2$ 来自矩阵~\eqref{eq:weighted-moduli-relations}。该矩阵的行给出 $A^*(\Xi)$ 的一组基，以及 $A(\Xi)$ 的对偶基；GKZ 分解就在后一个空间中。这正是上面描述 $V$ 的 Kähler 锥时用到的基 $D_3,D_1$。因此，将两边都与 Chow 商识别之后，$\overline{\mathcal M}_{\mathrm{simp}}(V^\circ)$ 的局部坐标 $z_1,z_2$ 就是 $\overline{\mathcal{KM}}(V)$ 的坐标 $q_1,q_2$；\hyperref[subsec:toric-mirror-map]{第~\ref*{subsec:toric-mirror-map} 小节}将仔细证明这一点。

为说明镜像映射为何不是 $q_j=z_j$，使用\hyperref[ex:weighted-extremal-transition]{例题~\ref*{ex:weighted-extremal-transition}}中从 $V$ 到 $\widetilde V$ 的极端跃迁。如上所见，由 $D_3$ 生成的 Kähler 锥面给出\IndexLocation{indexloc:1089:171}本原收缩 $\phi=\phi_{|D_3|}:V\to\overline V$。在 Kähler 模中，这个面由 $\operatorname{Im}(t_2)=0$ 定义，也就是 $|q_2|=1$。若镜像映射为恒等映射，则从 $\widetilde V^\circ$ 到 $V^\circ$ 的“对偶”极端跃迁，应在 $|z_2|=1$ 这一轨迹的某处粘合。然而，\hyperref[ex:weighted-extremal-transition]{例题~\ref*{ex:weighted-extremal-transition}}说明，这条轨迹由 $1-4z_2=0$，即 $z_2=1/4$ 定义。两者不一致，表明镜像映射不能是恒等映射。

物理学家会用 Kähler 模空间“受到瞬子修正”来解释这一差异。我们以前已见过这种现象，最突出的是：瞬子，即从 Riemann 曲面到 $V$ 的全纯映射，使 A 模型关联函数不同于交数。\hyperref[sec:mirror-map]{第~\ref*{sec:mirror-map} 节}将最终解释镜像映射。

还有一点很有趣：尽管 $|z_2|=1$ 与 $1-4z_2=0$ 不相容，它们至少涉及同一个变量。这并非偶然，而是 Morrison 关于极端跃迁中复模与 Kähler 模识别位置的一般猜想 \SourceCite{Morrison9} 的一部分。

第 \hyperref[subsec:conifold-transitions-conclusion]{\ref*{subsec:conifold-transitions-conclusion}}、\hyperref[subsec:other-physics-topics]{\ref*{subsec:other-physics-topics}} 小节将讨论极端跃迁的更多方面。第 \hyperref[subsec:flag-cy-complete-intersections]{\ref*{subsec:flag-cy-complete-intersections}} 小节还会说明，极端跃迁可用于构造 Grassmann 簇和部分旗流形中的 Calabi--Yau 完全交的镜像。

\section{镜像映射}\label{sec:mirror-map}
本节描述 Calabi--Yau 流形 $V$ 的复模空间与其镜像流形 $V^\circ$ 的 Kähler 模空间之间的镜像映射。预期这些模空间彼此同构；\SourcePage{172}
镜像映射是一个显式的局部同构，将 $V^\circ$ 的 Kähler 模空间在大半径极限点处的一组自然坐标，与 $V$ 的复模空间在相应最大幺幂边界点附近的一组自然坐标联系起来。

先给出镜像映射的一般定义，再将它应用于环面超曲面，最后讨论几个显式例子。

\subsection{镜像映射的定义}\label{subsec:mirror-map-definition}\index{镜像映射!定义|hyperpage}
\SourceCite{CdGP} 为五次三维簇写出了镜像映射；其他 Kähler 模空间为一维的例子，基本上也用同样的方法处理 \SourceCite{Font,KT1,KT2,Morrison1}。更高维模空间的镜像映射见 \SourceCite{CdFKM,HKTY1,CFKM,HKTY2}，而 \SourceCite{BCOV2} 给出了镜像映射的物理论证。\hyperref[chap:quantum-cohomology]{第~\ref*{chap:quantum-cohomology} 章}将基于量子上同调，从数学上说明镜像映射的动机；\hyperref[chap:mirror-theorem]{第~\ref*{chap:mirror-theorem} 章}给出另一种计算它的方法，例题~\ref{ex:quintic-hypergeometric-mirror-map} 末尾将透露一点线索。镜像映射的纯数学定义归于 Morrison；下面遵循 \SourceCite{Morrison6} 的论述。

为定义镜像映射，先在复模空间 $\overline{\mathcal M}(V)$ 中取一个最大幺幂边界点。按定义~\ref{def:maximal-unipotent-point}，取 $W_2\subset H^d(V)$ 的一组基 $g_0,\ldots,g_r$；这里 $W_2$ 来自单值性权滤过，$d=\dim V$。如式~\eqref{eq:mum-canonical-coordinates}，定义 $q_k$ 为
\begin{equation}\label{eq:mirror-map-period-definition}
\frac{1}{2\pi i}\log q_k
=\frac{1}{\langle g_0,\Omega\rangle}\sum_{j=1}^r\langle g_j,\Omega\rangle m^{jk},
\end{equation}
其中 $(m^{jk})$ 为定义~\ref{def:maximal-unipotent-point} 中矩阵 $(m_{ij})$ 的逆矩阵。由\hyperref[subsec:mum-boundary-points]{第 6.1.1 小节}，当 $p$ 位于所取最大幺幂边界点的邻域中时，$q_i=q_i(p)$ 是全纯函数。

根据定义~\ref{def:maximal-unipotent-point} 之后的讨论，还可如下改写 $q_k$ 的公式。首先，$y_0=\langle g_0,\Omega\rangle$ 是在最大幺幂边界点处全纯的 Picard--Fuchs 方程解，并且在相差常数倍的意义下唯一。再令
\[
y_k=\sum_{j=1}^r\langle g_j,\Omega\rangle m^{jk}.
\]
根据\hyperref[chap:hodge-yukawa]{第~\ref*{chap:hodge-yukawa} 章}，$y_k$ 为 Picard--Fuchs 方程的一个解，满足
\[
y_k=y_0\frac{\log z_k}{2\pi i}+\widetilde y_k,
\]
其中 $\widetilde y_k$ 在边界点处全纯。于是式~\eqref{eq:mirror-map-period-definition} 可写为
$\frac{1}{2\pi i}\log q_k=y_k/y_0$，从而
\begin{equation}\label{eq:mirror-map-holomorphic-correction}
q_k=z_k\exp(2\pi i\,\widetilde y_k/y_0).
\end{equation}
注意，固定 $y_0$ 后，$y_k$ 在相差 $y_0$ 的常数倍的意义下唯一。因此，$q_k$ 在相差常数倍的意义下唯一。

\SourcePage{173}
现在令 $V^\circ$ 为 $V$ 的镜像，假定 $V^\circ$ 的 Kähler 锥满足锥猜想。再取\hyperref[subsec:large-radius-limits]{第 6.2.1 小节}中锥 $K_+$ 的一个细分 $\Sigma_+$，使它给出未扩大的 Kähler 模空间的光滑部分紧化 $\overline{\mathcal{KM}}(V^\circ)$。每个极大维锥 $\sigma\in\Sigma_+$ 都给出 $\overline{\mathcal{KM}}(V^\circ)$ 的一个大半径极限点。由式~\eqref{eq:large-radius-exponential-coordinates}，$\sigma$ 的生成元 $T_i$ 确定典范坐标 $q_i$；它们除次序外是唯一的。

由于 $\overline{\mathcal M}(V)$ 与 $\overline{\mathcal{KM}}(V^\circ)$ 有相同的维数 $r$，这是二者互为镜像的条件之一，可以在 $\overline{\mathcal M}(V)$ 的最大幺幂边界点附近取 $p$，定义
\begin{equation}\label{eq:mirror-map-local}
M(p)=(q_1(p),\ldots,q_r(p))\in\overline{\mathcal{KM}}(V^\circ).
\end{equation}
这里 $q_i(p)$ 是上面定义的全纯函数，而目标空间使用镜像一侧所选大半径极限点处的典范局部坐标。因此，定义镜像映射 $M$ 需要两类数据：
\begin{itemize}
\item 式~\eqref{eq:mirror-map-period-definition} 中的函数 $q_1,\ldots,q_r$；
\item 给出大半径极限点的锥 $\sigma$ 的生成元 $T_1,\ldots,T_r\in H^2(V^\circ)$。
\end{itemize}
第一项完全由 $V$ 的最大幺幂边界点确定，第二项由 $V^\circ$ 的大半径极限点确定。

选定函数 $q_1,\ldots,q_r$ 及锥生成元 $T_1,\ldots,T_r$ 后，式~\eqref{eq:mirror-map-period-definition}、\eqref{eq:mirror-map-local} 就定义了映射 $M$。对于多数选择，$M$ 并不是我们所要的映射；然而，镜像对称的核心正在于：若满足下列条件，
\begin{itemize}
\item $W_2$ 中的类 $g_0,g_1,\ldots,g_r$ 选取得正确，等价地，Picard--Fuchs 方程的解 $y_0,y_1,\ldots,y_r$ 选取得正确；
\item 在两个模空间的相容紧化中，最大幺幂边界点及大半径极限点选取得正确，
\end{itemize}
那么式~\eqref{eq:mirror-map-local} 给出的就是镜像映射。

为理解怎样才算“选取得正确”，需要进一步解释这两点。第一点回到式~\eqref{eq:mum-coordinate-basis-change}：将 $g_0,\ldots,g_r$ 换为 $g_j'=\sum_{k=0}^rc_{jk}g_k$ 后，所得新局部坐标 $(q_1',\ldots,q_r')$ 满足
\begin{equation}\label{eq:mirror-map-constant-ambiguity}
q_k'=e^{2\pi i(c^{k0}/c_{00})}q_k,\qquad
c^{k0}=\sum_{j=1}^rc_{j0}m'^{jk},
\end{equation}
其中 $(m'^{jk})$ 为由 $N_j(g_k')=m'_{jk}g_0'$ 确定的矩阵的逆矩阵。

处理这种非唯一性，有两种已知的方法。第一种使用整数结构 \SourceCite{Morrison6}。单值性权滤过中的空间 $W_0\subset W_2$ 定义在 $\mathbb Q$ 上。用 $W_i\cap H^d(V,\mathbb Z)$ 表示 $W_i$ 与 $H^d(V,\mathbb Z)$ 在 $H^d(V,\mathbb Q)$ 中的像的交，则可以取
$g_0\in W_0\cap H^d(V,\mathbb Z)$、$g_1,\ldots,g_r\in W_2\cap H^d(V,\mathbb Z)$，使 $g_0$ 在 $\mathbb Z$ 上张成前一个格，而 $g_0,g_1,\ldots,g_r$ 在 $\mathbb Z$ 上张成后一个格。称这样的 $g_0,g_1,\ldots,g_r$ 为 $W_0\subset W_2$ 的一组 $\mathbb Z$ 基。

给定这样一组 $\mathbb Z$ 基，容易证明，由 $N_j(g_k)=m_{jk}g_0$ 定义的数 $m_{jk}$ 是整数。\hyperref[subsec:higher-dimensional-mum]{第 5.2.2 小节}的整性猜想断言，矩阵 $(m_{jk})$ 在 $\mathbb Z$ 上可逆；等价地，式~\eqref{eq:mirror-map-period-definition} 中的 $m^{jk}$ 为整数。五次三维簇的镜像已验证这一猜想 \SourceCite{Morrison2}；\SourceCite{Morrison6} 讨论了其他情形。

\SourcePage{174}
下面假定整性猜想成立。设两组 $\mathbb Z$ 基 $g_0,g_1,\ldots,g_r$ 与 $g_0',g_1',\ldots,g_r'$ 满足 $g_j'=\sum_{k=0}^rc_{jk}g_k$，则 $c_{00}=\pm1$ 且 $c_{jk}\in\mathbb Z$。由于还假定 $m'^{jk}\in\mathbb Z$，式~\eqref{eq:mirror-map-constant-ambiguity} 给出 $q_k'=q_k$。因此，只要使用 $\mathbb Z$ 基，式~\eqref{eq:mirror-map-period-definition} 的函数除次序外就是唯一的。

\SourceExampleNumber{6.3.1}{1}
\begin{example}\label{ex:mirror-map-integral-connection}
设 $V$ 为光滑 Calabi--Yau 三维簇，在一个最大幺幂边界点处满足整性猜想；将该点看作 $0\in\Delta^r$。命题~\ref{prop:normalized-yukawa-pf} 之后的讨论说明，可以找到变量 $q_1,\ldots,q_r$，使 Gauss--Manin 联络 $\nabla$ 为
\begin{equation}\label{eq:mirror-map-flat-connection}
\begin{aligned}
\nabla_{\delta_i}e^0&=0,\\
\nabla_{\delta_i}e^k&=\delta_{ik}e^0,&&1\leq k\leq r,\\
\nabla_{\delta_i}e_j&=\sum_{k=1}^rY_{ijk}e^k,&&1\leq j\leq r,\\
\nabla_{\delta_i}e_0&=e_i,
\end{aligned}
\end{equation}
其中 $\delta_j=2\pi iq_j\partial/\partial q_j$，见式~\eqref{eq:multimoduli-connection}。下面说明，此时 $q_1,\ldots,q_r$ 确定镜像映射。

命题~\ref{prop:normalized-yukawa-pf} 的证明中已说明，$e_0$ 是归一化的三形式 $\widetilde\Omega$，平坦截面 $e^0=g_0$ 是 $W_0$ 的整数生成元，且 $\langle g_0,\widetilde\Omega\rangle=1$。此外，$e^1,\ldots,e^r$ 为 $W_2$ 的截面。写 $q_j=\exp(2\pi iu_j)$，容易看出
\[
g_j=-e^j+u_je^0,\qquad1\leq j\leq r
\]
是平坦截面。逆时针绕 $(\Delta^*)^r$ 的第 $i$ 个因子一周，将 $u_j$ 变为 $u_j+\delta_{ij}$，所以 $\mathcal T_i(g_j)=g_j+\delta_{ij}g_0$，进而 $N_i(g_j)=\delta_{ij}g_0$。

由于 $e^0,e^1,\ldots,e^r$ 为 $W_2$ 的截面，$g_0,g_1,\ldots,g_r$ 也是。我们还断言这些截面都是整的。对 $g_0$ 已知如此；对其他截面，注意
$\exp(-\sum_i u_iN_i)g_j=-e^j$。与命题~\ref{prop:normalized-yukawa-pf} 的证明一样，$e^j$ 在 $0\in\Delta^r$ 上为整；由\hyperref[subsec:limiting-hodge]{第 5.1.4 小节}的讨论，这意味着平坦截面 $g_j$ 也为整。因此，$g_0,g_1,\ldots,g_r$ 构成 $W_2$ 的一组平坦整数截面基，且 $g_0$ 张成 $W_0$。

现在可以计算镜像映射。由于 $N_i(g_j)=\delta_{ij}g_0$，矩阵 $(m_{ij})=(\delta_{ij})$ 为单位矩阵，因此式~\eqref{eq:mirror-map-period-definition} 的右端化为 $\langle g_k,\widetilde\Omega\rangle$；回忆 $\langle g_0,\widetilde\Omega\rangle=1$。而
\[
\langle g_j,\widetilde\Omega\rangle
=\langle-e^j+u_je^0,e_0\rangle=u_j,
\]
因为式~\eqref{eq:multimoduli-integral-pairings} 给出 $\langle e^j,e_0\rangle=0$（$1\leq j\leq r$）。由式~\eqref{eq:mirror-map-period-definition}，$q_k$ 确实是镜像映射的第 $k$ 个坐标。
\end{example}

环面情形还有第二种方法，能在不依赖整性猜想的情况下，选出式~\eqref{eq:mirror-map-period-definition} 中正确的 $g_0,g_1,\ldots,g_r$。基本思想是：环面几何自然给出 $z_i$ 的选择；对式~\eqref{eq:mirror-map-holomorphic-correction} 中 $\widetilde y_k/y_0$ 的常数项作适当归一化，就唯一选出 $q_k$。\hyperref[subsec:toric-mirror-map]{第~\ref*{subsec:toric-mirror-map} 小节}将说明这一方法。

\SourcePage{175}
算出 $q_1,\ldots,q_r$ 后，“选取得正确”的另一面，是寻找 $\overline{\mathcal{KM}}(V^\circ)$ 中对应的大半径极限点。这可能并不简单：如\hyperref[subsec:mum-boundary-points]{第 6.1.1 小节}和\hyperref[subsec:large-radius-limits]{第 6.2.1 小节}所见，紧化可能不唯一，还可能有许多不同的点的等价类可供选择。一种办法是寻找适当的极小紧化，尽可能爆缩边界；这或许能减少将最大幺幂边界点与大半径极限点配对的选择。不过，多数时候的问题是，必须相当完整地了解模空间的形态。事实上，$V^\circ$ 为 $V$ 的镜像所包含的要求之一，就是\emph{要充分理解 $\overline{\mathcal M}(V)$ 和 $\overline{\mathcal{KM}}(V^\circ)$，从而知道怎样将最大幺幂边界点与大半径极限点相匹配}。最好的例子是环面情形：两个模空间本质上相同，所以能够准确知道边界点如何匹配。\hyperref[subsec:toric-mirror-map]{第~\ref*{subsec:toric-mirror-map} 小节}将作更详细的讨论。

此外，正确选定大半径极限点后，得到一个由 $T_1,\ldots,T_r$ 生成的锥 $\sigma$，但这些生成元应如何排序？\hyperref[subsec:mirror-map-derivative]{第~\ref*{subsec:mirror-map-derivative} 小节}讨论镜像映射的导数时，将看到这基本上由\hyperref[chap:introduction]{第~\ref*{chap:introduction} 章}式~\eqref{eq:mirror-cohomology-forms} 中的同构 $H^{d-1,1}(V)\simeq H^{1,1}(V^\circ)$ 决定。

所有选择都正确完成后，关键思想是：从 $\overline{\mathcal M}(V)$ 到 $\overline{\mathcal{KM}}(V^\circ)$ 的镜像映射，应与 A 模型及 B 模型关联函数相容。进一步还猜想，这种相容性可推广到 Hodge 结构变分。$\overline{\mathcal M}(V)$ 上有一个来自 $V$ 的复结构变化的自然 Hodge 结构变分；它的渐近行为决定某个边界点是否最大幺幂。\hyperref[chap:quantum-cohomology]{第~\ref*{chap:quantum-cohomology} 章}将定义 $\overline{\mathcal{KM}}(V^\circ)$ 上的 A 模型 Hodge 结构变分，并证明它在大半径极限点处最大幺幂。镜像猜想断言，镜像映射与这些 Hodge 结构变分相容；这将蕴含相应关联函数的相容性。\hyperref[chap:mirror-theorem]{第~\ref*{chap:mirror-theorem} 章}将说明，在环面背景下这一猜想能证明到什么程度。

\subsection{镜像映射的导数}\label{subsec:mirror-map-derivative}
设 $V,V^\circ$ 为维数 $d\geq3$ 的 Calabi--Yau 流形，假定镜像映射已按\hyperref[subsec:mirror-map-definition]{第 6.3.1 小节}定义。接下来要联系以下三个主题：
\begin{itemize}
\item 镜像映射 $M$ 的导数；
\item \hyperref[chap:introduction]{第~\ref*{chap:introduction} 章}式~\eqref{eq:mirror-cohomology-forms} 中的同构 $H^{d-1,1}(V)\simeq H^{1,1}(V^\circ)$；
\item 已知 $q_1,\ldots,q_r$ 后，怎样“正确选取”$T_1,\ldots,T_r$。
\end{itemize}
为研究镜像映射，用复化 Kähler 空间
\[
K_{\mathbb C}(V^\circ)
=(H^2(V^\circ,\mathbb R)+iK(V^\circ))/\operatorname{im}H^2(V^\circ,\mathbb Z)
\]
代替 Kähler 模空间 $\mathcal{KM}(V^\circ)$ 很方便。若假定\hyperref[subsec:large-radius-limits]{第 6.2.1 小节}的锥猜想成立，前者是后者的离散覆盖。为避免处理这一商，只将镜像映射看作 $M:S\to K_{\mathbb C}(V^\circ)$；其中 $S$ 是 $V$ 的复模空间在一个最大幺幂边界点附近的内侧邻域。

\SourcePage{176}
$V^\circ$ 的一个大半径极限点，由 Kähler 锥闭包中整类 $T_1,\ldots,T_r\in H^2(V^\circ,\mathbb C)=H^{1,1}(V^\circ)$ 生成的锥 $\sigma$ 确定。它在所选大半径极限点附近给出坐标
\begin{equation}\label{eq:mirror-map-kahler-coordinates}
(e^{2\pi it_1},\ldots,e^{2\pi it_r})
\longleftrightarrow\left[\sum_jt_jT_j\right]\in K_{\mathbb C}(V^\circ).
\end{equation}
由式~\eqref{eq:mirror-map-local}，$p\in S$ 时 $M(p)=(q_1(p),\ldots,q_r(p))$；式~\eqref{eq:mirror-map-holomorphic-correction} 随即说明，$M$ 可延拓到边界点，并在该处为局部双全纯映射。

为描述 $M$ 的导数，需要理解 $S$ 和 $K_{\mathbb C}(V^\circ)$ 的切丛。先看后者。每个 $T\in H^2(V^\circ,\mathbb C)$ 都在 $H^2(V^\circ,\mathbb R)+iK(V^\circ)$ 上给出一个平移不变的向量场 $\partial/\partial T$，因而可将它看作 $K_{\mathbb C}(V^\circ)$ 上的向量场。由此得到平凡化
\[
T_{K_{\mathbb C}(V^\circ)}\simeq H^2(V^\circ,\mathbb C)\times K_{\mathbb C}(V^\circ).
\]
坐标~\eqref{eq:mirror-map-kahler-coordinates} 给出向量场 $\partial/\partial t_j$，它自然与 $\partial/\partial T_j$ 识别；在上述同构下，$\partial/\partial t_j$ 对应于 $T_j$。利用这一平凡化，局部双全纯映射 $M$ 诱导
\begin{equation}\label{eq:mirror-map-differential}
D_M:T_S\longrightarrow H^2(V^\circ,\mathbb C)\times K_{\mathbb C}(V^\circ).
\end{equation}
由于 $q_1,\ldots,q_r$ 可用作 $S$ 的局部坐标，向量场 $\delta_j=2\pi iq_j\partial/\partial q_j$ 在内部构成 $T_S$ 的一组基；$M$ 的上述公式说明，$D_M$ 将 $\delta_j$ 映到 $\partial/\partial t_j=\partial/\partial T_j$。以式~\eqref{eq:mirror-map-differential} 的语言说，就是映到 $T_j$。

再描述 $S$ 的切丛 $T_S$。$V$ 的复模空间为 $\mathcal M(V)$；在最大幺幂边界点 $p_0$ 附近的内侧邻域 $S$ 上，假定有万有族 $\pi:\mathcal V\to S$；严格地说，应先取 $S$ 的有限覆盖。由于 $V$ 为 Calabi--Yau，Kodaira--Spencer 映射为同构 \SourceCite{Bogomolov,Tian,Todorov}：
\[
\kappa:T_S\simeq R^1\pi_*\Theta_{\mathcal V/S},
\]
其中 $\Theta_{\mathcal V/S}$ 为相对切丛。

还有 $\mathcal H=R^d\pi_*\mathbb C\otimes_{\mathbb C}\mathcal O_S$ 的 Hodge 滤过 $\mathcal F^p$。我们最关心
$\mathcal F^d=\pi_*\Omega^d_{\mathcal V/S}$ 与
$\mathcal F^{d-1}/\mathcal F^d=R^1\pi_*\Omega^{d-1}_{\mathcal V/S}$。内乘给出映射
\begin{equation}\label{eq:mirror-map-contraction}
R^1\pi_*\Theta_{\mathcal V/S}
\longrightarrow\operatorname{Hom}(\mathcal F^d,\mathcal F^{d-1}/\mathcal F^d).
\end{equation}
众所周知，它与 $\kappa$ 的复合
\begin{equation}\label{eq:mirror-map-kodaira-spencer}
T_S\longrightarrow\operatorname{Hom}(\mathcal F^d,\mathcal F^{d-1}/\mathcal F^d)
\end{equation}
将向量场 $X$ 与全纯 $d$ 形式 $\Omega$ 映为 $[\nabla_X(\Omega)]$；方括号表示在商 $\mathcal F^{d-1}/\mathcal F^d$ 中的类。

由于 $V$ 为 Calabi--Yau，可以在 $S$ 上固定 $\mathcal F^d$ 的一个处处非零截面 $\Omega$。照常令 $\widetilde\Omega=\Omega/\langle g_0,\Omega\rangle$ 作归一化；$g_0$ 来自\IndexLocation{indexloc:0772:176}最大幺幂单值性的定义。$\widetilde\Omega$ 诱导 $\mathcal F^d\simeq\mathcal O_S$，因而式~\eqref{eq:mirror-map-kodaira-spencer} 给出同构
\begin{equation}\label{eq:mirror-map-infinitesimal-period}
T_S\simeq R^1\pi_*\Theta_{\mathcal V/S}
\simeq R^1\pi_*\Omega^{d-1}_{\mathcal V/S},
\end{equation}
由 $X\mapsto[\nabla_X(\widetilde\Omega)]$ 定义。它的含义如下：设 $p\in S$、$V=V_p$，将 $X$ 看作把 $V_p$ 变为 $V_\varepsilon$ 的\IndexLocation{indexloc:0209:176}无穷小形变，则在此形变下，$\widetilde\Omega_p$ 变为
$\widetilde\Omega_\varepsilon=\widetilde\Omega_p+\varepsilon\nabla_X(\widetilde\Omega)_p+O(\varepsilon^2)$。

\SourcePage{177}
现在可以联系前面提到的三点。首先，对 $p\in S$、$V=V_p$，结合式~\eqref{eq:mirror-map-differential} 与~\eqref{eq:mirror-map-infinitesimal-period}，得到同构
\begin{equation}\label{eq:mirror-map-hodge-differential}
H^1(V,\Omega_V^{d-1})\simeq T_{S,p}
\simeq H^2(V^\circ,\mathbb C)=H^{1,1}(V^\circ).
\end{equation}
因此，镜像映射的导数就是\hyperref[chap:introduction]{第~\ref*{chap:introduction} 章}同构~\eqref{eq:mirror-cohomology-forms} 的 $d$ 维版本。此外，前面已看到 $D_M$ 将 $\delta_j=2\pi iq_j\partial/\partial q_j$ 映为 $T_j$。在式~\eqref{eq:mirror-map-hodge-differential} 中，这表示在 $V=V_p$ 处，将 $[\nabla_{\delta_j}\widetilde\Omega]$ 映为 $T_j$，即由 $2\pi iq_j\partial/\partial q_j\mapsto T_j$ 定义相应同构。因此，知道式~\eqref{eq:mirror-map-hodge-differential}，就等于知道怎样“正确选取”$T_1,\ldots,T_r$。

最后给出一个\hyperref[chap:quantum-cohomology]{第~\ref*{chap:quantum-cohomology} 章}会用到的引理。它将镜像映射的导数与丛 $\mathcal H=\mathcal F^0=R^d\pi_*\mathbb C\otimes_{\mathbb C}\mathcal O_S$ 联系起来。

\begin{lemma}\label{lem:mirror-map-hodge-splitting}
设 $\widetilde\Omega$ 为最大幺幂边界点 $p_0$ 附近的归一化 $d$ 形式。则：
\begin{enumerate}[label=(\roman*)]
\item 在 $p_0$ 附近，映射
$\psi:\mathcal O_S\oplus T_S\to\mathcal F^{d-1}\subset\mathcal H$，
$\psi(f,X)=f\widetilde\Omega+\nabla_X(\widetilde\Omega)$，为与式~\eqref{eq:mirror-map-infinitesimal-period} 相容的同构。边界处使用对数切丛及 Hodge 丛的典范延拓。
\item 经 $M$ 识别底空间后，有丛同构
\[
\phi:\mathcal F^{d-1}\simeq
\bigl(H^0(V^\circ,\mathbb C)\oplus H^2(V^\circ,\mathbb C)\bigr)\times M(S),
\]
它将 $\widetilde\Omega$ 映为 $1\in H^0(V^\circ,\mathbb C)$；将 $T_S$ 经 $\psi$ 视为 $\mathcal F^{d-1}$ 的子丛后，$\phi$ 在其上的限制就是镜像映射的导数。
\end{enumerate}
\end{lemma}
\begin{proof}
最大幺幂单值性的定义给出 $W_0=W_1$ 且维数为 1。杯积与单值性权滤过相容，因此 $\operatorname{Gr}^W_{2d}$ 的维数为 1，且 $W_{2d-1}=W_{2d-2}$。用于证明式~\eqref{eq:mum-top-splitting} 的论证给出
\[
F^d_{\mathrm{lim}}\oplus W_{2d-1}
=F^d_{\mathrm{lim}}\oplus W_{2d-2}=H^d(V,\mathbb C),
\]
其中 $F^d_{\mathrm{lim}}$ 来自极限 Hodge 滤过。由此，$\operatorname{Gr}^W_{2d-2}$ 上的 Hodge 结构纯为 $(d-1,d-1)$ 型，进而
$F^{d-1}_{\mathrm{lim}}/F^d_{\mathrm{lim}}\simeq\operatorname{Gr}^W_{2d-2}$。

结合最大幺幂单值性的定义与杯积在 $\operatorname{Gr}^W_\bullet$ 上给出的对偶性，可知：若 $g\in W_{2d}$ 生成 $\operatorname{Gr}^W_{2d}$，则 $N_j(g)$ 构成 $\operatorname{Gr}^W_{2d-2}$ 的一组基。再结合式~\eqref{eq:connection-monodromy-residue}，便知转到典范延拓后，引理在 $p_0$ 处成立。于是 $\mathcal O_S\oplus T_S\simeq\mathcal F^{d-1}$ 立即成立；其余结论也随之得到。
\end{proof}

式~\eqref{eq:mirror-map-flat-connection} 提供了这一引理的具体说明：$e_0=\widetilde\Omega$、$e_i=\nabla_{\delta_i}(\widetilde\Omega)$。若用 $1$ 作 $\mathcal O_S$ 的基，用 $\delta_i$ 作 $T_S$ 的基，则引理第一部分的同构将 $1$ 映为 $\widetilde\Omega$，将 $\delta_i$ 映为 $e_i$。特别地，式~\eqref{eq:mirror-map-flat-connection} 显式包含了 $T_S$ 的一个副本。第二部分的同构将 $e_0$ 映为 $1\in H^0(V^\circ,\mathbb C)$，将 $e_i$ 映为 $T_i$（$1\leq i\leq r$）。\hyperref[chap:quantum-cohomology]{第~\ref*{chap:quantum-cohomology} 章}将把 $\phi$ 扩大为 $\mathcal H=\mathcal F^0$ 与由 $\bigoplus_pH^{2p}(V^\circ,\mathbb C)$ 给出的平凡丛之间的同构。我们将看到，这属于镜像对称的 Hodge 理论表述。

\SourcePage{178}
\subsection{环面超曲面的镜像映射}\label{subsec:toric-mirror-map}
对于环面超曲面，可以更加精确地说明镜像映射应是什么，也能给出处理式~\eqref{eq:mirror-map-constant-ambiguity} 中非唯一性的另一种方案。

照常固定\IndexLocation{indexloc:1218:178}自反多胞形 $\Delta$，令 $V\subset X_\Sigma$ 为 Calabi--Yau 超曲面，$\Sigma$ 为 $\Delta$ 的法扇的简化射影细分。再令 $V^\circ\subset X_{\Sigma^\circ}$ 为 $V$ 的一个 Batyrev 镜像，置
\[
\Xi=(\Delta\cap M)_0-\{0\},\qquad
\Xi^+=(\Xi\cup\{0\})\times\{1\}=(\Delta\cap M)_0\times\{1\}.
\]
注意，$\Sigma^\circ(1)=\Xi$，因为 $\Sigma^\circ$ 是 $\Delta^\circ$ 的法扇的简化射影细分。还要回忆，$\Xi$ 的 GKZ 分解与 $\Xi^+$ 的次级扇，都位于向量空间 $A(\Xi)$ 中。

为了确定镜像映射，还需要更多数据。锥 $\operatorname{cpl}(\Sigma^\circ)$ 是 $\Xi$ 的 GKZ 分解中的一个锥。将它细分为光滑锥 $\sigma$；于是每个 $\sigma$ 都是单纯锥，其生成元构成格的一组整基。

我们的基本主张是：每个锥 $\sigma$ 都给出了确定镜像映射所需的全部信息。换言之，$\sigma$ 确定相互对应的最大幺幂边界点与\IndexLocation{indexloc:0696:178}大半径极限点，并且\emph{不必假定\IndexLocation{indexloc:0636:178}整性猜想，也不必知道整数上同调的任何信息}，就能选定式~\eqref{eq:mirror-map-constant-ambiguity} 中的常数。此外，从 $\sigma$ 出发，还会自动“正确选取”$T_1,\ldots,T_r$。

先看\hyperref[subsec:toric-complex-moduli]{第 6.1.2 小节}用 $\Xi^+$ 的次级扇构造的简化多项式模空间，见式~\eqref{eq:moduli-secondary-configuration}。由于希望紧化光滑，采用次级扇的一个细化，其中包含上述细分 $\operatorname{cpl}(\Sigma^\circ)$ 的各锥 $\sigma$。将这一紧化记为 $\overline{\mathcal M}_{\mathrm{simp}}$。假定猜想~\ref{conj:toric-mum-boundary} 成立，本章余下部分都作此假定，则每个 $\sigma$ 都给出它的一个最大幺幂边界点。

固定一个 $\sigma$ 后，可如下得到该最大幺幂边界点处的自然坐标。$\sigma$ 的生成元构成 $A(\Xi)$ 中的格的一组基；\hyperref[subsec:toric-complex-moduli]{第 6.1.2 小节}将该格记为 $N_0$。对偶基 $\ell_1,\ldots,\ell_r$ 生成 $A(\Xi)^*$ 中的格，即式~\eqref{eq:moduli-relation-lattice} 中的 $M_0$。它是 $\Xi^+$ 的关系格，所以每个基元素可写为 $\ell_k=(\ell_{kj})$。由式~\eqref{eq:moduli-monomial-coordinate}，
\begin{equation}\label{eq:toric-mirror-moduli-coordinates}
z_k=\prod_j\lambda_j^{\ell_{kj}}
\end{equation}
为最大幺幂边界点处的局部坐标。回忆，$\lambda_i$ 为 Laurent 多项式
\begin{equation}\label{eq:toric-mirror-laurent-polynomial}
\sum_{m_i\in(\Delta\cap M)_0}\lambda_i t^{m_i}
=\lambda_0+\sum_{m_i\in\Xi}\lambda_i t^{m_i}
\end{equation}
的系数；仍取 $m_0=0$。这样书写局部坐标，强调我们所处理的是模；更准确地说是简化模，因为\IndexLocation{indexloc:0217:178}占优猜想意味着 $\overline{\mathcal M}_{\mathrm{simp}}$ 为真正多项式模空间的有限覆盖。

镜像映射将是式~\eqref{eq:toric-mirror-moduli-coordinates} 中坐标 $z_1,\ldots,z_r$ 的函数。下面将看到，每个 $\ell_k=(\ell_{kj})$ 都贡献一个符号
\begin{equation}\label{eq:toric-mirror-sign}
(-1)^{\ell_{k0}},\qquad k=1,\ldots,r,
\end{equation}
它对确定镜像映射有重要作用。

再看\hyperref[subsec:enlarged-kahler-moduli]{第 6.2.3 小节}的环面 Kähler 模空间 $\overline{\mathcal{KM}}_{\mathrm{toric}}(V^\circ)$。假定猜想~\ref{conj:toric-kahler-cone} 成立，则在相差有限群作用的意义下，
\SourcePage{179}
环面 Kähler 模的一个紧化由 $\Xi^+$ 的次级扇给出，其中需将对应于平凡翻转的锥合并。为了得到与 $V$ 的复模相容的紧化，使用上述次级扇细化给出的光滑爆破，记为 $\overline{\mathcal{KM}}_{\mathrm{toric}}(V^\circ)$。特别地，细分 $\operatorname{cpl}(\Sigma^\circ)$ 的锥 $\sigma$ 给出该空间的大半径极限点。

对每个 $\sigma$，可如下得到相应大半径极限点处的局部坐标 $q_1,\ldots,q_r$。由定理~\ref{thm:batyrev-hodge} 的证明，首先有自然同构
$A(\Xi)\simeq H^2_{\mathrm{toric}}(V^\circ,\mathbb R)$。若 $\sigma$ 的生成元为 $T_1,\ldots,T_r$，则映射
\[
(\Delta^*)^r\longrightarrow
\frac{H^2_{\mathrm{toric}}(V^\circ,\mathbb R)+iK_{\mathrm{toric}}(V^\circ)}
{\operatorname{im}H^2_{\mathrm{toric}}(V^\circ,\mathbb Z)},
\qquad(e^{2\pi it_1},\ldots,e^{2\pi it_r})\longmapsto\left[\sum_{j=1}^rt_jT_j\right]
\]
延拓为开嵌入 $\Delta^r\subset\overline{\mathcal{KM}}_{\mathrm{toric}}(V^\circ)$，原点映到所选大半径极限点。该点处的局部坐标于是为 $q_j=e^{2\pi it_j}$；定义镜像映射时将使用它们。

紧化 $\overline{\mathcal M}_{\mathrm{simp}}(V)$ 与 $\overline{\mathcal{KM}}_{\mathrm{toric}}(V^\circ)$ 相同；固定 $\sigma$ 给出相容的\IndexLocation{indexloc:0766:179}最大幺幂边界点与大半径极限点。前者有局部坐标 $z_1,\ldots,z_r$，后者有 $q_1,\ldots,q_r$。这些坐标其实相同，虽然这一点未必显然。原因是 $z_j$ 对应于模空间中环面 $T_0$ 的特征标群的一组基，沿用\hyperref[subsec:toric-complex-moduli]{第 6.1.2 小节}的记号；而 $q_j$ 由 $T_0$ 的单参数子群的一组基确定。两组基互为对偶，因而给出相同坐标。为使记号清楚，仍在 $\overline{\mathcal M}_{\mathrm{simp}}(V)$ 上使用 $z_j$，在 $\overline{\mathcal{KM}}_{\mathrm{toric}}(V^\circ)$ 上使用 $q_j$。

现在可以确定镜像映射。按式~\eqref{eq:mirror-map-period-definition} 定义 $q_k=q_k(z_1,\ldots,z_r)$，并如下消除其非唯一性。由式~\eqref{eq:mirror-map-holomorphic-correction}，有 $q_k(z_1,\ldots,z_r)\sim cz_k$，其中 $c\ne0$ 为常数。利用式~\eqref{eq:mirror-map-constant-ambiguity}，可使
\begin{equation}\label{eq:toric-mirror-normalization}
q_k(z_1,\ldots,z_r)\sim(-1)^{\ell_{k0}}z_k,
\end{equation}
其中符号来自式~\eqref{eq:toric-mirror-sign}。这显然唯一确定函数 $q_1,\ldots,q_r$，也就定义了唯一的镜像映射~\eqref{eq:mirror-map-local}。用式~\eqref{eq:mirror-map-holomorphic-correction} 前的记号，就是寻找 Picard--Fuchs 方程的解 $y_k$，使
\begin{equation}\label{eq:toric-mirror-normalized-period}
y_k=y_0\frac{\log\bigl((-1)^{\ell_{k0}}z_k\bigr)}{2\pi i}+\widetilde y_k,
\end{equation}
其中 $\widetilde y_k$ 全纯，且在 $z_1=\cdots=z_r=0$ 处为零。于是
$q_k=(-1)^{\ell_{k0}}z_k\exp(2\pi i\,\widetilde y_k/y_0)$。

这一公式虽然看起来有些特别，但符号 $(-1)^{\ell_{k0}}$ 已出现在例题~\ref{ex:quintic-toric-mirror-map} 中，下面加以说明。

\SourceExampleNumber{6.3.3}{1}
\begin{example}\label{ex:quintic-toric-mirror-map}
\hyperref[chap:quintic]{第~\ref*{chap:quintic} 章}讨论\IndexLocation{indexloc:1195:179}五次三维簇 $V$ 时，镜像映射从镜像 $V^\circ$ 的复模，映到 $V$ 的 Kähler 模。
\SourcePage{180}
用齐次坐标，$V^\circ$ 的方程为
\[
\lambda_1x_1^5+\lambda_2x_2^5+\lambda_3x_3^5+\lambda_4x_4^5+\lambda_5x_5^5
+\lambda_0x_1x_2x_3x_4x_5=0,
\]
而 $\Xi^+$ 的元素之间的关系为 $\ell=(1,1,1,1,1,-5)$。如式~\eqref{eq:quintic-gkz-coordinate}，$V^\circ$ 的模坐标为 $z=\lambda_1\cdots\lambda_5/\lambda_0^5$。该式之后的讨论表明，$z$ 正是\hyperref[chap:quintic]{第~\ref*{chap:quintic} 章}称为 $x$ 的坐标。$V$ 的 Kähler 模还有坐标 $q$。

例题~\ref{ex:quintic-hypergeometric-mirror-map} 及其后续讨论将证明，$V^\circ$ 的 Picard--Fuchs 方程有全纯解
\[
y_0(z)=\sum_{n=0}^\infty\frac{(5n)!}{(n!)^5}(-z)^n,
\]
以及解
\[
y_0\frac{\log(-z)}{2\pi i}
+\frac5{2\pi i}\sum_{n=1}^\infty\frac{(5n)!}{(n!)^5}
\left[\sum_{j=n+1}^{5n}\frac1j\right](-z)^n,
\]
具有式~\eqref{eq:toric-mirror-normalization} 所要求的形式。因此，五次三维簇的环面镜像映射为
\[
q=-z\exp\left(\frac5{y_0(z)}\sum_{n=1}^\infty\frac{(5n)!}{(n!)^5}
\left[\sum_{j=n+1}^{5n}\frac1j\right](-z)^n\right).
\]
这就是式~\eqref{eq:q-series-c1} 中的镜像映射。\hyperref[chap:quintic]{第~\ref*{chap:quintic} 章}通过计算 $V^\circ$ 的 Yukawa 耦合，并利用一般 $V$ 含有 2875 条直线，确定右端前面的负号。现在知道了式~\eqref{eq:toric-mirror-normalization}，就直接由 $(-1)^{\ell_{k0}}=\allowbreak(-1)^{-5}=-1$ 得到这个负号。

还可与\hyperref[subsec:mirror-map-definition]{第 6.3.1 小节}定义的镜像映射比较。由于五次镜像满足整性猜想 \SourceCite{Morrison2}，能取满足 $N(g_1)=g_0$ 的整类 $g_0,g_1$，使式~\eqref{eq:mirror-map-period-definition} 中的镜像映射为
\[
\frac1{2\pi i}\log q=\frac{\langle g_1,\Omega\rangle}{\langle g_0,\Omega\rangle}.
\]
它与上述环面定义一致，并非显然。不过，如\hyperref[ex:quintic-normalized-yukawa]{例题~\ref*{ex:quintic-normalized-yukawa}}所述，这由 \SourceCite{CdGP,Morrison2} 的显式计算得到。令 $\gamma_0,\gamma_1$ 分别为 $g_0,g_1$ 的 Poincaré 对偶，上式还意味着
\[
q=\exp\left(2\pi i\frac{\int_{\gamma_1}\Omega}{\int_{\gamma_0}\Omega}\right).
\]
这正是\hyperref[sec:quintic-mirror-map]{第 2.3 节}开头的定义。因此，对五次三维簇，各种不同的镜像映射定义都一致。
\end{example}

理解式~\eqref{eq:toric-mirror-normalization} 的一种方式，是改变 $\lambda_0$ 的符号：在式~\eqref{eq:toric-mirror-laurent-polynomial} 中用 $-\lambda_0$ 代替 $\lambda_0$，则由式~\eqref{eq:toric-mirror-moduli-coordinates}，每个 $z_k$ 乘上符号 $(-1)^{\ell_{k0}}$。若一开始这样安排，就用 $q_k\sim z_k$ 归一化镜像映射，而不必使用式~\eqref{eq:toric-mirror-normalization}。\SourceCite{BvS} 采用这种方法。另一方面，\SourceCite{HKTY1} 使用与我们相同的 $z_k$，随后引入新变量 $x_k=(-1)^{\ell_{k0}}z_k$，用 $x_k$ 定义镜像映射。

\SourcePage{181}
另一种理解是：式~\eqref{eq:toric-mirror-normalization} 等价于指定镜像映射的导数。利用\hyperref[subsec:mirror-map-derivative]{第 6.3.2 小节}的思想，最大幺幂边界点处的切空间可自然识别为定理~\ref{thm:batyrev-hodge} 的证明中记为 $H^1_{\mathrm{poly}}(V,T_V)\simeq H^{n-2,1}_{\mathrm{poly}}(V)$ 的空间；类似地，大半径极限点处的切空间为 $H^2_{\mathrm{toric}}(V^\circ)$。对于坐标 $z_k,q_k$，分别有自然基 $\partial/\partial z_k$ 和 $\partial/\partial q_k$。在这些坐标中，镜像映射的导数为同构
\[
H^{n-2,1}_{\mathrm{poly}}(V)\simeq H^2_{\mathrm{toric}}(V^\circ),\qquad
\frac\partial{\partial z_k}\longmapsto
\sum_j\frac{\partial q_j}{\partial z_k}\frac\partial{\partial q_j}.
\]
在边界点 $z_1=\cdots=z_r=0$，式~\eqref{eq:toric-mirror-normalization} 表明，导数将 $\partial/\partial z_k$ 映为 $(-1)^{\ell_{k0}}\partial/\partial q_k$。理清这些记号，就知道这一映射把单项式确定的 $V$ 的一阶形变，映为来自 $V^\circ$ 的环面除子的上同调类，至多相差符号 $(-1)^{\ell_{k0}}$。因此，它，或者更准确地说，不带这些符号的版本，称为\emph{单项式–除子镜像映射} \SourceCite{AGM1}。

\SourceCite[第 4 节]{AGM1} 的基本猜想之一，是单项式–除子镜像映射在相差一个阶为 2 的自同构的意义下，为镜像映射的导数。所谓阶为 2 的自同构，表示需要指定某些符号；这正是式~\eqref{eq:toric-mirror-normalization} 所做的事情。在我们的处理下，鉴于环面情形的镜像映射定义，这一猜想由定义成立。不过，要真正称得上“镜像映射”，这里描述的映射必须与\hyperref[subsec:mirror-map-definition]{第 6.3.1 小节}末尾提到的关联函数和 Hodge 结构变分相容。这在完全一般的情形中尚未得到证明，因此仍有工作要做。

最后还可以问，为何式~\eqref{eq:toric-mirror-normalization} 是正确的归一化方式。物理学的解释见 \SourceCite{MP1}。从数学看，符号 $(-1)^{\ell_{k0}}$ 最容易用超几何函数解释；下一小节将讨论它们。

最后，我们计算一个熟悉例子的部分镜像映射。

\SourceExampleNumber{6.3.3}{2}
\begin{example}\label{ex:weighted-boundary-mirror-map}
$\mathbb P(1,1,2,2,2)$ 的环面爆破中反典范超曲面的镜像 $V^\circ$，已经讨论多次，最近一次是\hyperref[subsec:kahler-moduli-boundary]{第 6.2.4 小节}。这里不计算整个镜像映射，而只研究与\hyperref[ex:weighted-extremal-transition]{例题~\ref*{ex:weighted-extremal-transition}}中的极端跃迁相关的部分。回忆，实际上有两个跃迁：从 $V$ 到 $\widetilde V$ 的原跃迁，以及从 $\widetilde V^\circ$ 到 $V^\circ$ 的“对偶”跃迁。前者涉及 $V$ 的 Kähler 锥的一个面，后者使用 $V^\circ$ 的判别式轨迹的一部分。问题在于：它们在镜像映射下怎样匹配？

使用\hyperref[ex:weighted-complex-moduli]{例题~\ref*{ex:weighted-complex-moduli}}给出的最大幺幂边界点附近的局部坐标 $z_1,z_2$，在轨迹 $z_1=0$ 上研究镜像映射。这里的计算首先见于 \SourceCite{AGM3}，\SourceCite{Morrison9} 作了进一步阐述。目标是假定 $z_1=0$，将 $q_2$ 算为 $z_2$ 的函数。由于 $z_2=\lambda_1\lambda_2/\lambda_6^2$，有 GKZ 算子 $\Box=\partial_1\partial_2-\partial_6^2$，其中 $\partial_i$ 表示对 $\lambda_i$ 的偏导。导出式~\eqref{eq:weighted-logarithmic-identities} 的计算同样适用，给出周期的 Picard--Fuchs 方程
\[
\delta^2y-2z_2\delta(2\delta+1)y=0.
\]
\SourcePage{182}
这里 $y=y(z_2)$、$\delta=z_2\,d/dz_2$。一般解为
\begin{equation}\label{eq:weighted-boundary-period}
y(z_2)=c_1+c_2\log\left(\frac{\sqrt{1-4z_2}+2z_2-1}{2z_2}\right).
\end{equation}
为了计算镜像映射，需要 $z_2=0$ 附近的解析解与对数解。解析解只有常数，因而取 $y_0(z_2)=1$。由于 $\sqrt{1-4z_2}$ 的 Taylor 展开为 $1-2z_2-2z_2^2+\cdots$，式~\eqref{eq:weighted-boundary-period} 中乘 $c_2$ 的函数，在 $z_2=0$ 处按对数增长。为了使 $q_2(z_2)=\exp(2\pi i\,y_1(z_2)/y_0(z_2))$ 具有式~\eqref{eq:toric-mirror-normalization} 预言的渐近行为，取 $c_1=1/2$、$c_2=1/(2\pi i)$。这给出
\[
y_1(z_2)=\frac1{2\pi i}\log\left(\frac{1-2z_2-\sqrt{1-4z_2}}{2z_2}\right),
\]
于是 $z_1=0$ 上的镜像映射为
\begin{equation}\label{eq:weighted-boundary-mirror-coordinate}
q_2(z_2)=\frac{1-2z_2-\sqrt{1-4z_2}}{2z_2}
=z_2+2z_2^2+5z_2^3+\cdots.
\end{equation}
在 $\overline{\mathcal M}_{\mathrm{simp}}(V^\circ)$ 中，\hyperref[ex:weighted-extremal-transition]{例题~\ref*{ex:weighted-extremal-transition}}说明，对偶极端跃迁发生于\IndexLocation{indexloc:0215:182}判别式轨迹的一部分 $1-4z_2=0$。这条轨迹与 $z_1=0$ 交于 $z_2=1/4$。代入式~\eqref{eq:weighted-boundary-mirror-coordinate}，得到 $q_2=1$。而 $q_2=\exp(2\pi it_2)$，所以 $t_2\in\mathbb Z$，特别地，$t_2$ 的虚部为零。坐标为 $q_1,q_2$ 的点，对应于复化 Kähler 类
\[
\begin{aligned}
t_1T_1+t_2T_2
&=\operatorname{Re}(t_1)T_1+\operatorname{Re}(t_2)T_2\\
&\quad+i\bigl(\operatorname{Im}(t_1)T_1+\operatorname{Im}(t_2)T_2\bigr)=B+i\omega.
\end{aligned}
\]
$\operatorname{Im}(t_2)=0$ 意味着 $\omega$ 位于原极端跃迁所发生的 Kähler 锥面。因此，恰好得到镜像对称预言的相容性；也看清了为何出现 $z_2=1/4$ 而非 $z_2=1$。若镜像映射只是 $q_2=z_2$，才会得到后一个值。

下一小节将回到这一例子，用超几何函数计算完整镜像映射。
\end{example}

\subsection{用超几何函数求镜像映射}\label{subsec:hypergeometric-mirror-map}
\hyperref[sec:hypergeometric]{第~\ref*{sec:hypergeometric} 节}和\hyperref[subsec:hypergeometric-moduli]{第~\ref*{subsec:hypergeometric-moduli} 小节}用 GKZ 超几何方程组，研究 Calabi--Yau \IndexLocation{indexloc:0099:182}环面超曲面 $V\subset X_\Sigma$ 的简化模空间上的 Picard--Fuchs 方程。这里继续这一讨论，用超几何方法构造镜像映射。沿用上一小节的记号；特别地，假定所处位置是\IndexLocation{indexloc:0766:182}最大幺幂边界点，其局部坐标 $z_1,\ldots,z_r$ 由式~\eqref{eq:toric-mirror-moduli-coordinates} 给出。

$\overline{\mathcal M}_{\mathrm{simp}}$ 上的 GKZ 方程组由\hyperref[subsec:hypergeometric-moduli]{第~\ref*{subsec:hypergeometric-moduli} 小节}讨论的微分算子 $\Box_\ell\circ\lambda_0^{-1}$ 组成，其中 $\ell$ 属于式~\eqref{eq:moduli-relation-lattice} 中 $\Xi^+$ 的关系格 $M_0$。对 $f=\lambda_0+\sum_{m_i\in\Xi}\lambda_i t^{m_i}$，还知道 $n$ 形式
\[
\widetilde\omega=\frac{\lambda_0}{f}\frac{dt_1}{t_1}\wedge\cdots\wedge\frac{dt_n}{t_n}
\]
是一个解，即对所有 $\ell$，有 $\Box_\ell(\lambda_0^{-1}\widetilde\omega)=0$。

形式 $\widetilde\omega$ 位于 $X_\Sigma-V$ 上，其在 $V$ 上的留数为全纯 $(n-1)$ 形式 $\Omega$。此外，在 $T\subset X_\Sigma$ 上还有由 $|t_1|=\cdots=|t_n|=1$ 定义的圈 $\gamma$，
\SourcePage{183}
它给出上同调类 $g\in H^{n-1}(V)$。于是
\[
\langle g,\Omega\rangle=\frac1{(2\pi i)^n}\int_\gamma
\frac{\lambda_0}{f}\frac{dt_1}{t_1}\wedge\cdots\wedge\frac{dt_n}{t_n}.
\]
这里加入常数 $1/(2\pi i)^n$，以便计算积分。由\hyperref[chap:hodge-yukawa]{第~\ref*{chap:hodge-yukawa} 章}，知道这一函数是 Picard--Fuchs 方程的解。上述积分是 \emph{Euler 积分}的一个例子；这类 GKZ 方程组的解在 \SourceCite{GKZ1} 中有详细研究。\index{Euler积分@Euler 积分|hyperpage}

为了看清这个解的形式，沿用 \SourceCite{Batyrev3} 的方法，注意到
\[
\frac{\lambda_0}{f}
=\frac1{1-\sum_{m_i\in\Xi}\lambda_i(-\lambda_0)^{-1}t^{m_i}}
=\sum_{k=0}^\infty
\left(\sum_{m_i\in\Xi}\lambda_i(-\lambda_0)^{-1}t^{m_i}\right)^k.
\]
设 $\Xi$ 有 $s$ 个元素。把这一恒等式代入上述积分，逐项积分，便得
\begin{equation}\label{eq:toric-euler-period}
\langle g,\Omega\rangle=
\sum_{\substack{\ell=(u_1,\ldots,u_s,u_0)\in M_0\\u_1,\ldots,u_s\geq0}}
\frac{(-u_0)!}{u_1!\cdots u_s!}(-\lambda_0)^{u_0}
\lambda_1^{u_1}\cdots\lambda_s^{u_s}.
\end{equation}
这是因为 $\ell\in M_0$ 蕴含 $\sum_{i=1}^su_i m_i=0$ 及 $-u_0=u_1+\cdots+u_s$。式~\eqref{eq:toric-euler-period} 中 $\lambda_0$ 前的负号至关重要。

下面用一个具体例子说明级数~\eqref{eq:toric-euler-period}。

\SourceExampleNumber{6.3.4}{1}
\begin{example}\leavevmode\phantomsection\label{ex:quintic-hypergeometric-mirror-map}%
对于五次三维簇的镜像 $V^\circ$，$\Xi^+$ 的元素间的关系只有 $\ell=(1,1,1,1,1,-5)$ 的整数倍。因此，对 $(n,n,n,n,n,-5n)$，$n\geq0$，求和，得到
\begin{equation}\label{eq:quintic-euler-period}
\begin{aligned}
\langle g,\Omega\rangle
&=\sum_{n=0}^\infty\frac{(-(-5n))!}{(n!)^5}
(-1)^{-5n}\lambda_0^{-5n}\lambda_1^n\lambda_2^n\lambda_3^n\lambda_4^n\lambda_5^n\\
&=\sum_{n=0}^\infty\frac{(5n)!}{(n!)^5}(-z)^n,
\end{aligned}
\end{equation}
其中照常取 $z=\lambda_1\lambda_2\lambda_3\lambda_4\lambda_5/\lambda_0^5$。这恰好是\hyperref[chap:quintic]{第~\ref*{chap:quintic} 章}给出的 $y_0(z)$ 的公式~\eqref{eq:y0-hypergeometric}。
\end{example}

一般地，可证明解~\eqref{eq:toric-euler-period} 在最大幺幂边界点处全纯。

\begin{proposition}\label{prop:toric-period-holomorphic}
将级数~\eqref{eq:toric-euler-period} 写成最大幺幂边界点处坐标 $z_1,\ldots,z_r$ 的函数，便得到该点某个邻域内处处非零的全纯函数。
\end{proposition}
\begin{proof}
第一步是用 $z_1,\ldots,z_r$ 写出式~\eqref{eq:toric-euler-period}，并证明只出现非负指数。$z_k$ 来自 $\Xi^+$ 的关系格的一组基 $\ell_k=(\ell_{kj})$，所以式~\eqref{eq:toric-euler-period} 中任一 $\ell$ 都形如 $\ell=n_1\ell_1+\cdots+n_r\ell_r$，其中 $n_1,\ldots,n_r$ 为整数。因此，该级数可写为
\begin{equation}\label{eq:toric-period-local-series}
\sum_{n_1,\ldots,n_r}
\frac{\bigl(-\sum_{k=1}^rn_k\ell_{k0}\bigr)!}
{\prod_{j=1}^s\bigl(\sum_{k=1}^rn_k\ell_{kj}\bigr)!}
\prod_{k=1}^r\bigl((-1)^{\ell_{k0}}z_k\bigr)^{n_k},
\end{equation}
其中求和遍及满足 $\sum_{k=1}^rn_k\ell_{kj}\geq0$，$j=1,\ldots,s$，的所有 $n_1,\ldots,n_r$。需要证明，这些不等式蕴含 $n_1,\ldots,n_r\geq0$；也就是，每个满足 $u_1,\ldots,u_s\geq0$ 的
\SourcePage{184}
$\ell=(u_1,\ldots,u_s,u_0)\in M_0$，都是 $\ell_k=(\ell_{kj})$ 的非负整数线性组合。

只要理清 $M_0$ 及其对偶 $N_0$ 的元素的各种解释，这一点其实很容易看出。在 $\mathbb R$ 上，有
$N_0\otimes\mathbb R=A(\Xi)=H^2_{\mathrm{toric}}(V^\circ,\mathbb R)$，而 $M_0\otimes\mathbb R=A(\Xi)^*$ 为其对偶。$X_{\Sigma^\circ}$ 上每个环面除子 $D_j$，在 $V^\circ$ 上给出仍记为 $D_j$ 的除子，因而有 $[D_j]\in H^2_{\mathrm{toric}}(V^\circ,\mathbb R)$。由于 $\Sigma^\circ(1)=\Xi$，$[D_1],\ldots,[D_s]$ 给出 $M_0$ 上的线性泛函。进一步，作为 $M_0$ 上的泛函，$[D_j]$ 取出 $\ell=(u_1,\ldots,u_s,u_0)$ 的第 $j$ 个坐标。于是 $u_1,\ldots,u_s\geq0$ 就是说，$\ell$ 位于各 $[D_j]$ 生成的锥的对偶锥中。各 $[D_j]$ 生成的锥正是 $A(\Xi)$ 的有效锥，也是 GKZ 分解的支集。指定所选最大幺幂边界点的锥 $\sigma$ 位于 $\operatorname{cpl}(\Sigma^\circ)$ 中，故位于有效锥中；因此，其对偶锥包含有效锥的对偶锥。于是 $\ell$ 位于 $\sigma$ 的对偶锥中，而这个锥由 $\ell_k=(\ell_{kj})$ 生成。

最后证明级数在原点的某个邻域收敛。展开 $(1+\cdots+1)^{u_1+\cdots+u_s}$，得到初等估计
\[
\frac{(u_1+\cdots+u_s)!}{u_1!\cdots u_s!}\leq s^{u_1+\cdots+u_s}.
\]
若像上面一样取 $u_j=\sum_{k=1}^rn_k\ell_{kj}$，则
$\sum_{j=1}^su_j\leq C\sum_{k=1}^rn_k$，其中适当的常数 $C$ 与各 $n_k$ 无关，因为 $n_1,\ldots,n_r\geq0$。级数的常数项显然为 $1$，故命题容易得到。
\end{proof}

这一命题表明，同一个级数~\eqref{eq:toric-euler-period} 在\emph{每个}最大幺幂边界点处都全纯。证明还提示了 GKZ 方程组与 GKZ 分解之间一个有趣的联系。更多细节见 \SourceCite{GKZ1}。

在最大幺幂边界点，Picard--Fuchs 方程只有一个全纯解，相差一个常数因子；这就是单值性权滤过满足 $\dim W_0=1$ 的含义。因此，命题~\ref{prop:toric-period-holomorphic} 描述了镜像映射公式~\eqref{eq:mirror-map-period-definition} 中的分母 $\langle g_0,\Omega\rangle$。为简便起见，以 $y_0$ 表示式~\eqref{eq:toric-period-local-series} 给出的级数。

为了求镜像映射的第 $k$ 个分量 $q_k$，需要找到满足式~\eqref{eq:toric-mirror-normalized-period} 的 Picard--Fuchs 方程解。不过，处理超几何函数时，去掉该式中的因子 $1/(2\pi i)$ 比较方便。因此，寻找满足
\begin{equation}\label{eq:hypergeometric-log-period}
y_k=y_0\log\bigl((-1)^{\ell_{k0}}z_k\bigr)+\widetilde y_k,
\qquad1\leq k\leq r,
\end{equation}
的解 $y_k$，其中 $\widetilde y_k$ 在原点全纯，且 $\widetilde y_k(0)=0$。在这一归一化下，$q_k$ 的公式成为
\begin{equation}\label{eq:hypergeometric-mirror-coordinate}
q_k=(-1)^{\ell_{k0}}z_k\exp(\widetilde y_k/y_0)
\sim(-1)^{\ell_{k0}}z_k,\qquad1\leq k\leq r,
\end{equation}
因为 $y_0$ 在原点非零，而 $\widetilde y_k$ 在原点为零。由式~\eqref{eq:toric-mirror-normalization}，这恰好符合要求。

求满足式~\eqref{eq:hypergeometric-log-period} 的 $y_k$，已有若干方法。先介绍 \SourceCite{HLY1} 中 Frobenius 方法的一个变体，它很适合\IndexLocation{indexloc:0615:184}超几何方程。其思想是，将式~\eqref{eq:toric-period-local-series} 中每个 $n_k$ 换成 $n_k+\rho_k$，其中 $\rho_k$ 为实参数。这样，诸如 $\bigl(\sum_{k=1}^rn_k\ell_{kj}\bigr)!$ 的阶乘，
\SourcePage{185}
换成经典的 Euler $\Gamma$ 函数
\[
\Gamma\left(\sum_{k=1}^r(n_k+\rho_k)\ell_{kj}+1\right).
\]
指数也随之改变，即 $\bigl((-1)^{\ell_{k0}}z_k\bigr)^{n_k}$ 变成
$\bigl((-1)^{\ell_{k0}}z_k\bigr)^{n_k+\rho_k}$。将这样由式~\eqref{eq:toric-period-local-series} 得到的级数记为 $y_0(z;\rho)$，其中 $z$ 代表 $z_1,\ldots,z_r$，$\rho$ 代表 $\rho_1,\ldots,\rho_r$。

若 $\partial_{\rho_k}=\partial/\partial\rho_k$，关键结果是
\[
y_k=\left.\partial_{\rho_k}y_0(z;\rho)\right|_{\rho=0},
\qquad1\leq k\leq r,
\]
具有所需的性质~\eqref{eq:hypergeometric-log-period}。所用的恒等式为
\[
\begin{aligned}
\partial_{\rho_k}\bigl((-1)^{\ell_{k0}}z_k\bigr)^{n_k+\rho_k}
&=\log\bigl((-1)^{\ell_{k0}}z_k\bigr)
\bigl((-1)^{\ell_{k0}}z_k\bigr)^{n_k+\rho_k},\\
[\Box_\ell,\partial_{\rho_k}]&=0.
\end{aligned}
\]
这里略去证明细节，改用几个例子说明实际计算。\index{Frobenius方法@Frobenius 方法|hyperpage}

\begingroup
\SourceExampleNumber{6.3.4}{1}
\renewcommand{\theHexam}{6.3.4.continued.\arabic{exam}}
\begin{example}[续]\label{ex:quintic-frobenius-continued}
回到前面开始讨论的五次镜像。式~\eqref{eq:quintic-euler-period} 给出的 $y_0(z)$ 意味着
\begin{equation}\label{eq:quintic-gamma-frobenius}
y_0(z;\rho)=\sum_{n=0}^\infty
\frac{\Gamma(5n+5\rho+1)}{\Gamma(n+\rho+1)^5}(-z)^{n+\rho}.
\end{equation}
于是 $y_0(z)=y_0(z;0)$，而上述理论说明，GKZ 方程的对数解为
\[
\begin{aligned}
y_1(z)&=\left.\partial_\rho y_0(z;\rho)\right|_{\rho=0}\\
&=y_0(z)\log(-z)+
\sum_{n=0}^\infty
\left.\partial_\rho\left(\frac{\Gamma(5n+5\rho+1)}
{\Gamma(n+\rho+1)^5}\right)\right|_{\rho=0}(-z)^n.
\end{aligned}
\]
为了化简，需要了解 $\Gamma$ 函数的导数。幸运的是，$\Gamma'(s)$ 在正整数处有以下公式：
\[
\Gamma'(n+1)=\Gamma(n+1)\left(-\gamma+\sum_{j=1}^n\frac1j\right),
\]
其中 $\gamma$ 为 Euler 常数。用归纳法及经典事实 $\Gamma'(1)=-\gamma$，可证明该式，见 \SourceCite[第 236 页]{WW}。由此容易得到
\[
y_1(z)=y_0(z)\log(-z)+5\sum_{n=1}^\infty
\frac{(5n)!}{(n!)^5}\left[\sum_{j=n+1}^{5n}\frac1j\right](-z)^n,
\]
式~\eqref{eq:hypergeometric-mirror-coordinate} 于是给出
\[
q=-z\exp\left(\frac5{y_0(z)}\sum_{n=1}^\infty
\frac{(5n)!}{(n!)^5}\left[\sum_{j=n+1}^{5n}\frac1j\right](-z)^n\right).
\]
这恰好是\hyperref[chap:quintic]{第~\ref*{chap:quintic} 章}式~\eqref{eq:y0-hypergeometric} 之后给出的五次三维簇镜像映射公式。

现在改用标准的 Frobenius 方法，它以递推关系为基础，类似于 \SourceCite{BvS} 的做法。为简化记号，令 $x=-z$。
\SourcePage{186}
采用这一变量，五次镜像族的 Picard--Fuchs 方程~\eqref{eq:quintic-pf} 可写为
\[
\delta^4y=5x(5\delta+1)(5\delta+2)(5\delta+3)(5\delta+4)y,
\]
其中 $\delta=x\,d/dx$，见式~\eqref{eq:quintic-gkz-fifth} 及随后的讨论。标准 Frobenius 方法 \SourceCite{CL} 使用函数
\[
\widetilde y(x;\rho)=\sum_{n=0}^\infty a_n(\rho)x^{n+\rho}.
\]
设 $L$ 为 Picard--Fuchs 方程给出的算子，需要选择 $a_n(\rho)$，使 $\widetilde y(x;\rho)$ 满足
\begin{equation}\label{eq:quintic-frobenius-inhomogeneous}
L(\widetilde y(x;\rho))=\rho^4x^\rho,
\end{equation}
因为 Picard--Fuchs 方程的\IndexLocation{indexloc:0621:186}指标方程为 $\rho^4=0$，对应\IndexLocation{indexloc:0772:186}最大幺幂单值性。对 $n>0$，得到递推关系
\[
(\rho+n)^4a_n(\rho)
=5(5\rho+5n-1)(5\rho+5n-2)(5\rho+5n-3)(5\rho+5n-4)a_{n-1}(\rho).
\]
选择 $a_0=1$ 的解，解出递推关系，得到
\[
a_n(\rho)=\frac{\prod_{m=1}^{5n}(5\rho+m)}
{\prod_{m=1}^n(\rho+m)^5}.
\]
于是
\begin{equation}\label{eq:quintic-product-frobenius}
\widetilde y(x;\rho)=\sum_{n=0}^\infty
\frac{\prod_{m=1}^{5n}(5\rho+m)}{\prod_{m=1}^n(\rho+m)^5}x^{n+\rho}.
\end{equation}
由式~\eqref{eq:quintic-frobenius-inhomogeneous} 容易得出，函数
\begin{equation}\label{eq:quintic-frobenius-basis}
y_i(x)=\left.\frac1{i!}\frac{\partial^i\widetilde y(x;\rho)}
{\partial\rho^i}\right|_{\rho=0},\qquad0\leq i\leq3,
\end{equation}
满足 Picard--Fuchs 方程，且给出 $x=0$ 附近解的一组基。注意，可把 $y_i(x)$ 看成 Taylor 级数
\begin{equation}\label{eq:quintic-frobenius-expansion}
\widetilde y(x;\rho)=y_0(x)+y_1(x)\rho+y_2(x)\rho^2+y_3(x)\rho^3+\cdots
\end{equation}
的系数。

读者可能会担心，这与前面使用 $\Gamma$ 函数的 Frobenius 方法不同。不过容易验证，$\widetilde y(x;\rho)$ 与式~\eqref{eq:quintic-gamma-frobenius} 定义的 $y_0(z;\rho)$ 满足
\[
y_0(z;\rho)=\widetilde y(-z;\rho)\frac{\Gamma(5\rho+1)}{\Gamma(\rho+1)^5}.
\]
经典公式，或 Mathematica，给出幂级数展开
\[
\frac{\Gamma(5\rho+1)}{\Gamma(\rho+1)^5}
=1+\frac{5\pi^2}{3}\rho^2+\cdots,
\]
从而
\[
\left.\frac{\partial\widetilde y(-z;\rho)}{\partial\rho}\right|_{\rho=0}
=\left.\frac{\partial y_0(z;\rho)}{\partial\rho}\right|_{\rho=0}.
\]
因此，两种 Frobenius 方法对对数解 $y_1$ 给出相同的公式。

更令人惊讶的是，从更形式化的角度看待这一计算时会发生什么。先将 $x$ 换成 $e^{t_1}$，于是 $\delta=x\,d/dx=d/dt_1$；再将 $\rho$ 换成
\SourcePage{187}
$\mathbb P^4$ 上同调环中的超平面类 $H$。同时将求和指标由 $n$ 改为 $d$，式~\eqref{eq:quintic-product-frobenius} 就成为函数
\[
\widetilde I=e^{t_1H}\sum_{d=0}^\infty
\frac{\prod_{m=1}^{5d}(5H+m)}{\prod_{m=1}^d(H+m)^5}e^{dt_1},
\]
这正是\hyperref[chap:hodge-yukawa]{第~\ref*{chap:hodge-yukawa} 章}式~\eqref{eq:quintic-tilde-i} 考察的函数。因此，从 Frobenius 方法看，乍看奇怪的上同调值函数 $\widetilde I$ 十分自然；感谢 \SourceCite{LLY} 的作者向我们解释这一点。按 $H^*(\mathbb P^4)$ 的通常基展开 $\widetilde I$，得到
\[
\widetilde I=y_0(x)+y_1(x)H+y_2(x)H^2+y_3(x)H^3+y_4(x)H^4.
\]
这就是在式~\eqref{eq:quintic-frobenius-expansion} 中用 $H$ 代替 $\rho$，因为 $H^5=0$。于是 $y_0,y_1,y_2,y_3$ 是式~\eqref{eq:quintic-frobenius-basis} 中的解，而 $y_4$ 不是。这证实了\hyperref[subsec:nef-gkz]{第~\ref*{subsec:nef-gkz} 小节}关于式~\eqref{eq:quintic-i-expansion} 所作的断言。进一步，若像式~\eqref{eq:quintic-iv-hypergeometric} 一样定义 $I_{\mathcal V}=5H\widetilde I$，则
\[
\begin{aligned}
I_{\mathcal V}
&=e^{t_1H}5H\sum_{d=0}^\infty
\frac{\prod_{m=1}^{5d}(5H+m)}{\prod_{m=1}^d(H+m)^5}e^{dt_1}\\
&=5H\bigl(y_0(x)+y_1(x)H+y_2(x)H^2+y_3(x)H^3\bigr),
\end{aligned}
\]
其中 $y_0,y_1,y_2,y_3$ 给出\IndexLocation{indexloc:1176:187}五次镜像 Picard--Fuchs 方程解的一组基。这些公式已经作为式~\eqref{eq:iv-simple} 和~\eqref{eq:iv-period-expansion} 出现在\hyperref[chap:quintic]{第~\ref*{chap:quintic} 章}中。用这种形式写出 Picard--Fuchs 方程的解，首先见于 \SourceCite{Givental1}。如\hyperref[subsec:nef-gkz]{第~\ref*{subsec:nef-gkz} 小节}末尾所述，$\widetilde I$ 在 \SourceCite{LLY} 中以 $\lim_{\lambda\to0}HG[\mathcal I(P)](t_1)$ 出现。上述求解超几何方程的讨论，也说明了为何采用“$HG$”记号。函数
$I_{\mathcal V}=\lim_{\lambda\to0}HG[\mathcal I(\widehat P)](t_1)$ 将在\hyperref[chap:mirror-theorem]{第~\ref*{chap:mirror-theorem} 章}镜像定理的处理中起重要作用。
\end{example}
\endgroup

\SourceExampleNumber{6.3.4}{2}
\begin{example}\label{ex:weighted-hypergeometric-mirror-map}
最后一个例子，回到 $\mathbb P(1,1,2,2,2)$ 的环面爆破中反典范\IndexLocation{indexloc:1020:187}\IndexLocation{indexloc:1029:187}超曲面的镜像 $V^\circ$。例题~\ref{ex:weighted-boundary-mirror-map} 在假定 $z_1=0$ 时计算了镜像映射。用\IndexLocation{indexloc:0616:187}超几何函数，现在能确定所有 $z_1,z_2$ 上的映射。例题~\ref{ex:octic-toric-mirror-theorem} 将再次讨论这一镜像映射，并进一步利用\hyperref[subsec:nef-gkz]{第~\ref*{subsec:nef-gkz} 小节}中本例的函数 $\widetilde I$。

解 $y_0(z_1,z_2)$ 很容易写出，因为格 $M_0$ 的生成元为矩阵~\eqref{eq:weighted-moduli-relations} 的两行
\[
(0,0,1,1,1,1,-4),\qquad(1,1,0,0,0,-2,0).
\]
于是
\[
(u_1,u_2,u_3,u_4,u_5,u_6,u_0)
=n_1(0,0,1,1,1,1,-4)+n_2(1,1,0,0,0,-2,0)
\]
给出式~\eqref{eq:toric-euler-period} 中的 $u_1=u_2=n_2$、
$u_3=u_4=u_5=n_1$、$u_6=n_1-2n_2$ 和 $u_0=-4n_1$，故式~\eqref{eq:toric-period-local-series} 给出级数
\[
y_0(z_1,z_2)=
\sum_{\substack{n_1,n_2\geq0\\n_1\geq2n_2}}
\frac{(4n_1)!}{(n_1!)^3(n_2!)^2(n_1-2n_2)!}z_1^{n_1}z_2^{n_2}.
\]
条件 $n_1\geq2n_2$ 表明 $y_0(0,z_2)=1$，与例题~\ref{ex:weighted-boundary-mirror-map} 的结果一致。

为了描述具有对数增长的解，按上述方式定义 $y_0(z_1,z_2;\rho_1,\rho_2)$，并计算
\[
y_k(z_1,z_2)=\left.\partial_{\rho_k}y_0(z_1,z_2;\rho_1,\rho_2)\right|_{\rho_1=\rho_2=0}.
\]
可以采用上一个例子的方法，
\SourcePage{188}
但必须适当处理条件 $n_1\geq2n_2$。这里微妙的一点是，$y_0(z_1,z_2)$ 实际上为级数
\[
y_0(z_1,z_2)=\sum_{n_1,n_2\geq0}
\frac{\Gamma(4n_1+1)}
{\Gamma(n_1+1)^3\Gamma(n_2+1)^2\Gamma(n_1-2n_2+1)}
z_1^{n_1}z_2^{n_2}.
\]
这与先前的级数相同，因为 Euler $\Gamma$ 函数在 $0,-1,-2,\ldots$ 处有单极点，其倒数在这些点为零。但由此构造 $y_0(z_1,z_2;\rho_1,\rho_2)$ 时，$n_1\geq2n_2$ 与 $n_1<2n_2$ 两类项都出现，对 $\rho_1,\rho_2$ 的导数中也都会出现。

用 $y_2=\left.\partial_{\rho_2}y_0\right|_{\rho_1=\rho_2=0}$ 说明这一过程。令 $\rho=\rho_2$，只需处理
\begin{equation}\label{eq:weighted-gamma-frobenius}
\sum_{n_1,n_2\geq0}
\frac{\Gamma(4n_1+1)}
{\Gamma(n_1+1)^3\Gamma(n_2+\rho+1)^2\Gamma(n_1-2n_2-2\rho+1)}
z_1^{n_1}z_2^{n_2+\rho}.
\end{equation}
对 $\rho$ 求导后置 $\rho=0$，得到 $y_0\log z_2+\widetilde y_2$，其中 $\widetilde y_2$ 来自式~\eqref{eq:weighted-gamma-frobenius} 的系数的导数。它有两类项，分别满足 $n_1\geq2n_2$ 和 $n_1<2n_2$。前者可用上一例题的方法处理，后者则需多做一点计算。使用 \SourceCite[第 239 页]{WW} 中的恒等式
\[
\Gamma(s)\Gamma(1-s)=\frac\pi{\sin(\pi s)},
\]
当 $n_1-2n_2<0$ 时，得到
\[
\left.\partial_\rho\bigl(\Gamma(n_1-2n_2-2\rho+1)^{-1}\bigr)\right|_{\rho=0}
=2(-1)^{n_1}(2n_2-n_1-1)!.
\]
因此，得到所需公式：
\[
\begin{aligned}
y_2(z_1,z_2)
&=y_0(z_1,z_2)\log z_2\\
&\quad+2\sum_{\substack{n_1,n_2\geq0\\n_1<2n_2}}
\frac{(4n_1)!(2n_2-n_1-1)!}{(n_1!)^3(n_2!)^2}
(-1)^{n_1}z_1^{n_1}z_2^{n_2}\\
&\quad+2\sum_{\substack{n_1,n_2\geq0\\n_1\geq2n_2}}
\frac{(4n_1)!}{(n_1!)^3(n_2!)^2(n_1-2n_2)!}
\left[\sum_{j=1}^{n_1-2n_2}\frac1j-\sum_{j=1}^{n_2}\frac1j\right]
z_1^{n_1}z_2^{n_2}.
\end{aligned}
\]
写成 $y_2=y_0\log z_2+\widetilde y_2$ 后，由式~\eqref{eq:hypergeometric-mirror-coordinate}，镜像映射有 $q_2=z_2\exp(\widetilde y_2/y_0)$。这给出了 $q_2$ 的显式公式，虽然相当复杂；也能为 $q_1$ 求出类似公式。至此，镜像映射已完全确定。

作为检验，像例题~\ref{ex:weighted-boundary-mirror-map} 一样置 $z_1=0$。所有 $n_1\geq2n_2$ 的求和项都消失；又因为 $y_0(0,z_2)=1$，得到
\begin{equation}\label{eq:weighted-boundary-log-series}
y_2(0,z_2)=\log z_2+
2\sum_{n_2=1}^\infty\frac{(2n_2-1)!}{(n_2!)^2}z_2^{n_2}.
\end{equation}
另一方面，例题~\ref{ex:weighted-boundary-mirror-map} 给出
\begin{equation}\label{eq:weighted-boundary-log-closed}
y_2(0,z_2)=\log\left(\frac{1-2z_2-\sqrt{1-4z_2}}{2z_2}\right).
\end{equation}
这是式~\eqref{eq:weighted-boundary-mirror-coordinate} 前一公式的表达式；因子 $1/(2\pi i)$ 消失，是因为本小节采用了不同的归一化方式。由式~\eqref{eq:weighted-boundary-period}，这两个公式相同，因为 $y_2(0,z_2)$ 的级数满足例题~\ref{ex:weighted-boundary-mirror-map} 的 Picard--Fuchs 方程
$\delta^2y-2z_2\delta(2\delta+1)y=0$，并具有相同的归一化。因此，虽然这里采用了与例题~\ref{ex:weighted-boundary-mirror-map} 很不同的方法求镜像映射，所得结果完全一致。
\end{example}

\SourcePage{189}
最后补充一点：上述两个例题介绍的镜像映射计算方法，已经用于大量环面超曲面，见 \SourceCite{HKTY1}。适当修改后，这一方法也已用于许多完全交，见 \SourceCite{BvS,HKTY2}。
